\documentclass[a4paper]{article}
\usepackage[margin=0cm,top=2cm]{geometry}
\usepackage[utf8]{inputenc}
\usepackage[german]{babel}
\usepackage{microtype}
\usepackage{graphicx}
\usepackage{color}
\usepackage{tikz}
\usepackage{rotating}
\usepackage[scaled]{helvet}
\usepackage[babel,german=quotes]{csquotes}
\usepackage{fancyhdr}

%   Colors for basic elements
\definecolor{framebg}{RGB}{0,0,0}
\definecolor{contentbg}{RGB}{255,255,255}

%   Colors to distinguish the extensions
\definecolor{basicgame}{RGB}{152,170,149}
\definecolor{intrigue}{RGB}{225,106,85}
\definecolor{seaside}{RGB}{100,145,189}
\definecolor{alchemy}{RGB}{255,245,158}
\definecolor{prosperity}{RGB}{125,195,170}
\definecolor{cornucopia}{RGB}{154,199,119}
\definecolor{hinterlands}{RGB}{175,159,104}
\definecolor{darkages}{RGB}{198,58,69}
\definecolor{guilds}{RGB}{225,111,159}
\definecolor{adventures}{RGB}{199,186,152}
\definecolor{empires}{RGB}{105,181,103}
\definecolor{nocturne}{RGB}{59,75,189}
\definecolor{promo}{RGB}{241,177,59}

%  Change emph command to writing text bold
\DeclareTextFontCommand{\emph}{\boldmath\bfseries}

% Change section numbering to Alph
\renewcommand\thesection{\Alph{section}}

% Define Heading
\pagestyle{fancy}
\fancyhf{}
\fancyhead[C]{Dominion Card Dividers - Created by Danilo Gasdzik based on the works of Eiko Wagenknecht. \\ Dominion Kartentrenner - Erstellt von Danilo Gasdzik basierend auf der Arbeit von Eiko Wagenknecht. \\ \leftmark}
\fancyfoot[L]{\vspace{-8em}\tiny{\hspace{1cm}Entlang der schwarzen Linien ausschneiden. Je nach gewünschter\\\hspace{1cm}Größe der Kartentrenner die obere oder untere Linie wählen.}}
\fancyfoot[C]{\vspace{-8em}\thepage}
\fancyfoot[R]{\vspace{-8em}\tiny{Cut along the black edge. For your desired size choose the upper or lower line.\hspace{1cm}}}
\renewcommand{\headrulewidth}{0pt}

% Enum without indent
\newcounter{mycounter}  
\newenvironment{noindlist}
{\begin{list}{\arabic{mycounter}.~~}{\usecounter{mycounter} \labelsep=0em \labelwidth=2em \leftmargin=2em \itemindent=0em \itemsep=-0.5em \topsep=1em}}
{\end{list}}

% Itemize without spacing
\usepackage{enumitem}
\setlist[itemize]{leftmargin=2em, itemsep=-1em, topsep=1em}

% Shortcuts for often used icons
\newcommand{\coin}[1]{\includegraphics[height=1em]{banner/coin.png}\hspace{-0.7em}\raisebox{0.25em}{\tiny{#1}}\hspace{0.3em}}
\newcommand{\hex}{\includegraphics[height=1em]{banner/hex.png}}
\newcommand{\potion}{\includegraphics[height=1em]{banner/potion.png}}

%   Draw cards in portrait format
%   Card in portrait
%   ---------------------------------------

%   TikZ/PGF size settings for cards
\pgfmathsetmacro{\cardwidth}{6.3cm}
\pgfmathsetmacro{\cardheight}{11.1cm}
\pgfmathsetmacro{\bannerwidth}{3.2cm}
\pgfmathsetmacro{\iconwidth}{0.5cm}
\pgfmathsetmacro{\iconoffset}{0.175cm}
\pgfmathsetmacro{\stripwidth}{0.2cm}
\pgfmathsetmacro{\stripheight}{0.7cm}
\pgfmathsetmacro{\contentoffset}{0.25cm}
\pgfmathsetmacro{\contentoffsettop}{1.6cm}
\pgfmathsetmacro{\contentwidth}{\cardwidth-\contentoffset-\contentoffset}
\pgfmathsetmacro{\contentheight}{\cardheight-\contentoffset-\contentoffsettop}
\pgfmathsetmacro{\cardoffset}{0.5cm}

%   Stylings for elements
\tikzset{
    %   runde Ecken für die Karten
    cardcorners/.style={
        rounded corners=0.4cm
    }
}


%   Create card with black background, white box and extension on top right
\newcommand{\card}{
	\draw ([xshift=0.1cm,yshift=-\cardoffset]0,0) -- ([xshift=0.1cm,yshift=\cardheight+\cardoffset]0,0);
	\draw ([xshift=\cardwidth-0.1cm,yshift=-\cardoffset]0,0) -- ([xshift=\cardwidth-0.1cm,yshift=\cardheight+\cardoffset]0,0);
	\draw ([xshift=-\cardoffset,yshift=0.1cm]0,0) -- ([xshift=\cardwidth+\cardoffset,yshift=0.1cm]0,0);
	\draw ([xshift=-\cardoffset,yshift=\cardheight-0.1cm]0,0) -- ([xshift=\cardwidth+\cardoffset,yshift=\cardheight-0.1cm]0,0);
	\draw ([xshift=-\cardoffset,yshift=\cardheight-0.3cm]0,0) -- ([xshift=\cardwidth+\cardoffset,yshift=\cardheight-0.3cm]0,0);
	\node [anchor=south west,fill=framebg,minimum width=\cardwidth,minimum height=\cardheight] (main) at (0,0) {};
	\node [anchor=south west,fill=contentbg,minimum width=\contentwidth,minimum height=\contentheight,cardcorners] (content) at ([yshift=\contentoffset,xshift=\contentoffset]0,0) {};
	\node [anchor=south east,\cardcolor,minimum width=\stripwidth,minimum height=\stripheight] at ([xshift=-0.6*\stripwidth,yshift=\contentoffset]content.north east) {\scriptsize{\textsc{\cardextension}}};
	\node [anchor=south east,\cardcolor,minimum width=\stripwidth,minimum height=\stripheight] at ([xshift=-0.6*\stripwidth]content.north east) {\scriptsize{\textsc{\cardextensiontitle}}};
}

%   Add banner on top left
\newcommand{\cardbanner}[1]{
	\node [anchor=south west] (cardbanner) at ([yshift=0.2*\contentoffset]content.north west) {
		\includegraphics[width=\bannerwidth]{#1}
	};
}

%   Add icon to banner
\newcommand{\cardicon}[1]{
	\node [anchor=south west] (cardicon) at ([xshift=\iconoffset,yshift=\iconoffset]cardbanner.south west) {
		\includegraphics[width=\iconwidth]{#1}
	};
}

%   Add price to icon
\newcommand{\cardprice}[1]{
	\node [] at (cardicon) {
		\small{
			\textsc{#1}
		}
	};
}


%   Add second icon
\newcommand{\cardiconaddition}[1]{
	\node [anchor=south west] (cardiconaddition) at ([xshift=\iconwidth]cardicon.south west) {
		\includegraphics[width=\iconwidth]{#1}
	};
}

%   Add price to second icon
\newcommand{\cardpriceaddition}[1]{
	\node [] at (cardiconaddition) {
		\small{
			\textsc{#1}
		}
	};
}

%   Add title of the card to banner
\newcommand{\cardtitle}[1]{
	\node [] at ([xshift=0.55*\iconwidth,yshift=1.1*\iconwidth]cardbanner.south) {
		\small{
			\textsc{#1}
		}
	};
}

%   Add colored strip on top right
\newcommand{\cardstrip}{
	\node [anchor=north east,fill=\cardcolor,minimum width=\stripwidth,minimum height=\stripheight] at ([xshift=-\contentoffset]main.north east) {};
}

%   Write content into white area
\newcommand{\cardcontent}[1]{
	\node [] at (content) {
		\rotatebox{270}{
			\parbox[t][\contentwidth][c]{\contentheight}{
				\begin{minipage}{\contentheight}
					\centering
					\begin{minipage}{0.9\textwidth}
						\raggedright
						\scriptsize{
							\textsf{#1}
						}
					\end{minipage}
				\end{minipage}
			}
		}
	};
}
%   Draw cards in landscape format
%%   Card in landscape
%   ---------------------------------------

%   TikZ/PGF size settings for cards
\pgfmathsetmacro{\cardwidth}{9.7cm}
\pgfmathsetmacro{\cardheight}{7.7cm}
\pgfmathsetmacro{\bannerwidth}{3.2cm}
\pgfmathsetmacro{\iconwidth}{0.5cm}
\pgfmathsetmacro{\iconoffset}{0.175cm}
\pgfmathsetmacro{\stripwidth}{0.2cm}
\pgfmathsetmacro{\stripheight}{0.7cm}
\pgfmathsetmacro{\contentoffset}{0.25cm}
\pgfmathsetmacro{\contentoffsettop}{1.6cm}
\pgfmathsetmacro{\contentwidth}{\cardwidth-\contentoffset-\contentoffset}
\pgfmathsetmacro{\contentheight}{\cardheight-\contentoffset-\contentoffsettop}
\pgfmathsetmacro{\cardoffset}{0.5cm}

%   Stylings for elements
\tikzset{
	%   runde Ecken für die Karten
	cardcorners/.style={
		rounded corners=0.4cm
	}
}


%   Create card with black background, white box and extension on top right
\newcommand{\card}{
	\draw ([xshift=0.1cm,yshift=-\cardoffset]0,0) -- ([xshift=0.1cm,yshift=\cardheight+\cardoffset]0,0);
	\draw ([xshift=\cardwidth-0.1cm,yshift=-\cardoffset]0,0) -- ([xshift=\cardwidth-0.1cm,yshift=\cardheight+\cardoffset]0,0);
	\draw ([xshift=-\cardoffset,yshift=0.1cm]0,0) -- ([xshift=\cardwidth+\cardoffset,yshift=0.1cm]0,0);
	\draw ([xshift=-\cardoffset,yshift=\cardheight-0.1cm]0,0) -- ([xshift=\cardwidth+\cardoffset,yshift=\cardheight-0.1cm]0,0);
	\draw ([xshift=-\cardoffset,yshift=\cardheight-0.3cm]0,0) -- ([xshift=\cardwidth+\cardoffset,yshift=\cardheight-0.3cm]0,0);
	\node [anchor=south west,fill=framebg,minimum width=\cardwidth,minimum height=\cardheight] (main) at (0,0) {};
	\node [anchor=south west,fill=contentbg,minimum width=\contentwidth,minimum height=\contentheight,cardcorners] (content) at ([yshift=\contentoffset,xshift=\contentoffset]0,0) {};
	\node [anchor=south east,\cardcolor,minimum width=\stripwidth,minimum height=\stripheight] at ([xshift=-0.6*\stripwidth,yshift=1.25*\contentoffset]content.north east) {\textsc{\cardextension}};
	\node [anchor=south east,\cardcolor,minimum width=\stripwidth,minimum height=\stripheight] at ([xshift=-0.6*\stripwidth]content.north east) {\small{\textsc{\cardextensiontitle}}};
}

%   Add banner on top left
\newcommand{\cardbanner}[1]{
	\node [anchor=south west] (cardbanner) at ([yshift=0.2*\contentoffset]content.north west) {
		\includegraphics[width=\bannerwidth]{#1}
	};
}

%   Add icon to banner
\newcommand{\cardicon}[1]{
	\node [anchor=south west] (cardicon) at ([xshift=\iconoffset,yshift=\iconoffset]cardbanner.south west) {
		\includegraphics[width=\iconwidth]{#1}
	};
}

%   Add price to icon
\newcommand{\cardprice}[1]{
	\node [] at (cardicon) {
		\small{
			\textsc{#1}
		}
	};
}


%   Add second icon
\newcommand{\cardiconaddition}[1]{
	\node [anchor=south west] (cardiconaddition) at ([xshift=\iconwidth]cardicon.south west) {
		\includegraphics[width=\iconwidth]{#1}
	};
}

%   Add price to second icon
\newcommand{\cardpriceaddition}[1]{
	\node [] at (cardiconaddition) {
		\small{
			\textsc{#1}
		}
	};
}

%   Add title of the card to banner
\newcommand{\cardtitle}[1]{
	\node [] at ([xshift=0.55*\iconwidth,yshift=1.1*\iconwidth]cardbanner.south) {
		\small{
			\textsc{#1}
		}
	};
}

%   Add colored strip on top right
\newcommand{\cardstrip}{
	\node [anchor=north east,fill=\cardcolor,minimum width=\stripwidth,minimum height=\stripheight] at ([xshift=-\contentoffset]main.north east) {};
}

\newcommand{\cardcontent}[1]{
	\node [] at (content) {
		\parbox[t][\contentheight][c]{\contentwidth}{
			\begin{minipage}{\contentwidth}
				\centering
				\begin{minipage}{0.9\textwidth}
					\raggedright
					\scriptsize{
						\textsf{#1}
					}
				\end{minipage}
			\end{minipage}
		}
	};
}

% Basic settings for this card set will be overwritten on each card set
\newcommand{\cardcolor}{}
\newcommand{\cardextension}{}
\newcommand{\cardextensiontitle}{}

\begin{document}
	\begin{center}
		\begin{minipage}{0.75\textwidth}
			\Huge 
			Dominion Kartentrenner
			\normalsize
			\tableofcontents
		\end{minipage}
	    	
	    % Basic settings for this card set
\renewcommand{\cardcolor}{}
\renewcommand{\cardextension}{}
\renewcommand{\cardextensiontitle}{}

\clearpage
\newpage
\section{Anleitung und Grundausstattung}

\begin{tikzpicture}
	\card
	\cardbanner{banner/white.png}
	\cardtitle{Anleitung (1)\quad}
	\cardcontent{\tiny{\emph{Spielablauf:} Dominion wird zugweise gespielt. Der Spieler an der Reihe hat am Beginn seines Zuges normalerweise 5 Karten auf der Hand. Er führt nun seinen Zug aus, der aus den 3 folgenden Phasen besteht, die immer in dieser Reihenfolge gespielt werden müssen.

	\medskip

	\emph{1. Phase:} Aktion - Der Spieler \emph{darf} Aktionskarten ausspielen.

	\emph{2. Phase:} Kauf - Der Spieler \emph{darf} Karten kaufen.

	\emph{3. Phase:} Aufräumen - Der Spieler \emph{muss} alle  ausgespielten \emph{und} alle Handkarten offen auf seinen Ablagestapel legen und \emph{sofort} 5 Karten für den nächsten Zug nachziehen.

	\medskip

	Die 1. und die 2. Phase darf, die 3. Phase muss gespielt werden. Wenn der Spieler seinen Zug beendet hat, ist der nächste Spieler an der Reihe. Das Spiel verläuft in dieser Weise bis Spielende.
	
	\medskip

	\begin{tabbing}
		Stapel der anderen Spieler: \= Links \kill
		Eigener Ablagestapel: \> Darf weder durchgezählt noch durchgesehen werden. \\
		Eigener Nachziehstapel: \> Darf durchgezählt, nicht aber durchgesehen werden. \\
		Stapel der anderen Spieler: \> Dürfen weder durchgezählt noch durchgesehen werden. \\
		Stapel im Vorrat: \> Dürfen jederzeit durchgezählt und durchgesehen werden. \\
		Müllstapel: \> Darf jederzeit durchgezählt und durchgesehen werden. \\
	\end{tabbing}}}
\end{tikzpicture}
\hspace{-0.6cm}
\begin{tikzpicture}
	\card
	\cardbanner{banner/white.png}
	\cardtitle{Anleitung (2)\quad}
	\cardcontent{\tiny{\begin{tabbing}
	Spielende:xxx \=  18x Provinzen,xxx \= 12 andere Punktekarten,xxx \= 50 Flüche,xxx \= Geld: \kill
	2 Spieler: \> 8x Provinzen, \> 8 andere Punktekarten, \> 10 Flüche, \> Geld: 46 K, 40 S, 30 G\\
	3 Spieler: \> 12x Provinzen, \> 12 andere Punktekarten, \> 20 Flüche, \> Geld: 39 K, 40 S, 30 G\\
	4 Spieler: \> 12x Provinzen, \> 12 andere Punktekarten, \> 30 Flüche, \> Geld: 32 K, 40 S, 30 G\\
	5 Spieler: \> 15x Provinzen, \> 12 andere Punktekarten, \> 40 Flüche, \> Geld: 85 K, 80 S, 60 G\\
	6 Spieler: \> 18x Provinzen, \> 12 andere Punktekarten, \> 50 Flüche, \> Geld: 78 K, 80 S, 60 G\\
	Spielende: \> Provinzstapel, Kolonienstapel (Dominion – Blütezeit) \emph{oder} \\
					\>3 Stapel (1 - 4 Spieler) bzw. 4 Stapel (5 - 6 Spieler) aus dem Vorrat leer \\
	\end{tabbing}
	Punktekarten aus Erweiterungen werden in Anzahl der \enquote{anderen Punktekarten} ausgelegt. Zur Spielende-Bedingung zählen alle Karten im Vorrat, also auch Fluch-, Geld- und Punktekarten, nicht jedoch z. B. der Müllstapel.

	\medskip

	\emph{Platin und Kolonie (Dominion – Blütezeit):} Die im Spiel befindlichen Königreich-Platzhalterkarten werden gemischt. Ist die erste gezogene Königreichskarte eine Karte aus der Blütezeit-Erweiterung, so wird mit Platin und Kolonie gespielt, ansonsten ohne.

	\medskip

	\emph{Unterschlupf-Karten (Dominion – Dark Ages):} Die im Spiel befindlichen Königreich-Platzhalterkarten werden gemischt. Ist die erste gezogene Königreichskarte eine Karte aus der Dark Ages-Erweiterung, so startet jeder Spieler mit 7 Kupfer, 1 Hütte, 1 Totenstadt und 1 Verfallenes Anwesen, andernfalls erhält jeder Spieler 7 Kupfer und 3 Anwesen.}}
\end{tikzpicture}
\hspace{-0.6cm}
\begin{tikzpicture}
	\card
	\cardbanner{banner/white.png}
	\cardtitle{Platzhalter\quad}
\end{tikzpicture}
\hspace{-0.6cm}
\begin{tikzpicture}
	\card
	\cardbanner{banner/gold.png}
	\cardicon{icons/coin.png}
	\cardprice{0}
	\cardtitle{Kupfer}
\end{tikzpicture}
\hspace{-0.6cm}
\begin{tikzpicture}
	\card
	\cardbanner{banner/gold.png}
	\cardicon{icons/coin.png}
	\cardprice{3}
	\cardtitle{Silber}
\end{tikzpicture}
\hspace{-0.6cm}
\begin{tikzpicture}
	\card
	\cardbanner{banner/gold.png}
	\cardicon{icons/coin.png}
	\cardprice{6}
	\cardtitle{Gold}
\end{tikzpicture}
\hspace{-0.6cm}
\begin{tikzpicture}
	\card
	\cardbanner{banner/green.png}
	\cardicon{icons/coin.png}
	\cardprice{2}
	\cardtitle{Anwesen}
\end{tikzpicture}
\hspace{-0.6cm}
\begin{tikzpicture}
	\card
	\cardbanner{banner/green.png}
	\cardicon{icons/coin.png}
	\cardprice{5}
	\cardtitle{Herzogtum}
\end{tikzpicture}
\hspace{-0.6cm}
\begin{tikzpicture}
	\card
	\cardbanner{banner/green.png}
	\cardicon{icons/coin.png}
	\cardprice{8}
	\cardtitle{Provinz}
\end{tikzpicture}
\hspace{-0.6cm}
\begin{tikzpicture}
	\card
	\cardbanner{banner/purple.png}
	\cardicon{icons/coin.png}
	\cardprice{0}
	\cardtitle{Fluch}
\end{tikzpicture}	
\hspace{0.6cm}

	    % Basic settings for this card set
\renewcommand{\cardcolor}{basicgame}
\renewcommand{\cardextension}{Edition I}
\renewcommand{\cardextensiontitle}{Das Basisspiel}

\clearpage
\newpage
\section{\cardextension \ - \cardextensiontitle}

\begin{tikzpicture}
	\card
	\cardstrip
	\cardbanner{banner/blue.png}
	\cardicon{banner/coin.png}
	\cardprice{2}
	\cardtitle{Burggraben}
	\cardcontent{Spielt ein anderer Spieler eine Angriffskarte (mit der Aufschrift AKTION -- ANGRIFF), kannst du die Karte \emph{BURGGRABEN} vorzeigen, falls du sie in diesem Moment auf der Hand hast. In diesem Fall bist \emph{du} von der Auswirkung des Angriffst nicht betroffen, d.h. du musst mit der \emph{HEXE} keine Fluchkarte nehmen usw. Haben mehrer Spieler einen \emph{BURGGRABEN} auf der Hand, dürfen diese auch eingesetzt und vorgezeigt werden. Danach nehmen die Spieler ihre Karte zurück auf die Hand. Der Spieler, der den Angriff gespielt hat, darf unabhängig davon, ob ein oder mehrere \emph{BURGGRÄBEN} gespielt wurden, die weiteren Anweisungen seiner Aktionskarte ausführen. Der \emph{BURGGRABEN} darf auch in der eigenen Aktionsphase gespielt werden -- dann ziehst du 2 Karten nach.}
\end{tikzpicture}
\hspace{-1cm}
\begin{tikzpicture}
	\card
	\cardstrip
	\cardbanner{banner/white.png}
	\cardicon{banner/coin.png}
	\cardprice{2}
	\cardtitle{Keller}
	\cardcontent{Der ausgespielte \emph{KELLER} selber darf nicht abgelegt werden, da er sich nicht mehr in deiner Hand befindet. Sage an wie viele Karten du ablegst und lege diese auf deinen Ablagestapel. Danach ziehst du die gleiche Anzahl Karten vom Nachziehstapel. Sollte während dieses Vorgangs der Nachziehstapel aufgebraucht werden, wird dein Ablagestapel zusammen mit den soeben abgelegten Karten gemischt und als neuer Nachziehstapel bereitgelegt.}
\end{tikzpicture}
\hspace{-1cm}
\begin{tikzpicture}
	\card
	\cardstrip
	\cardbanner{banner/white.png}
	\cardicon{banner/coin.png}
	\cardprice{3}
	\cardtitle{Dorf}
	\cardcontent{Spielst du mehrere \emph{DÖRFER} hintereinander, zählst du am besten laut mit, wie viele Aktionen du noch spielen darfst, damit du den Überblick behältst.}
\end{tikzpicture}
\hspace{-1cm}
\begin{tikzpicture}
	\card
	\cardstrip
	\cardbanner{banner/white.png}
	\cardicon{banner/coin.png}
	\cardprice{3}
	\cardtitle{Werkstatt}
	\cardcontent{Nimm dir eine Karte aus dem Vorrat und lege diese sofort auf deinen Ablagestapel. Du kansnt weder Geldkarten noch zusätzlich über Aktionskarten erhaltenes Geld oder Münzen (bei Erweiterungen mit Münzen) einsetzen, um den angegebenen Betrag auf der Karte zu erhöhen.}
\end{tikzpicture}
\hspace{-1cm}
\begin{tikzpicture}
	\card
	\cardstrip
	\cardbanner{banner/white.png}
	\cardicon{banner/coin.png}
	\cardprice{3}
	\cardtitle{Holzfäller}
	\cardcontent{In der Kaufphase darfst du 1 zusätzliche Karte kaufen, also insgesamt 2 Käufe tätigen. Für deine Käufe stehen dir in diesem Zug insgesamt 2 Geld zusätzlich zur Verfügung.}
\end{tikzpicture}
\hspace{-1cm}
\begin{tikzpicture}
	\card
	\cardstrip
	\cardbanner{banner/white.png}
	\cardicon{banner/coin.png}
	\cardprice{4}
	\cardtitle{Schmiede}
	\cardcontent{Du musst 3 Karten vom Nachziehstapel ziehen und auf die Hand nehmen.}
\end{tikzpicture}
\hspace{-1cm}
\begin{tikzpicture}
	\card
	\cardstrip
	\cardbanner{banner/white.png}
	\cardicon{banner/coin.png}
	\cardprice{4}
	\cardtitle{Umbau}
	\cardcontent{Der ausgespielte \emph{UMBAU} selber darf nicht entsorgt werden, da er sich nicht mehr in der Hand befindet. Weitere \emph{UMBAU}-Karten in deiner Hand dürfen entsorgt werden. Wenn du keine Karte zum Entsorgen auf der Hand hast, darfst du dir auch keine neue Karte nehmen. Die neue Karte, die du dir nimmst, darf maximal bis zu \coin{2} mehr als die entsorgte Karte kosten. Der Betrag darf weder durch weitere Geldkarten, Münzen oder zusätzliches Geld von anderen Aktionskarten erhöht werden. Die neue Karte kann die gleiche Karte sein, wie die, die du entsorgt hast. Lege die neue Karte auf deinen Ablagestapel.}
\end{tikzpicture}
\hspace{-1cm}
\begin{tikzpicture}
	\card
	\cardstrip
	\cardbanner{banner/white.png}
	\cardicon{banner/coin.png}
	\cardprice{4}
	\cardtitle{Miliz}
	\cardcontent{Deine Mitspieler müssen Karten aus ihrer Hand ablegen, bis sie nur noch 3 Karten auf der Hand haben. Spieler, die zum Zeitpunkt des Angriffs bereits 3 oder weniger Karten auf der Hand haben, müssen keine weiteren Karten ablegen.}
\end{tikzpicture}
\hspace{-1cm}
\begin{tikzpicture}
	\card
	\cardstrip
	\cardbanner{banner/white.png}
	\cardicon{banner/coin.png}
	\cardprice{5}
	\cardtitle{Markt}
	\cardcontent{Du musst eine Karte vom Nachziehstapel auf die Hand nehmen. Du \emph{darfst} in der Aktionsphase eine weiter Aktionskarte ausspielen. Du \emph{darfst} in der Kaufphase einen zusätzlichen Kauf tätigen und hast dafür ein zusätzliches Geld zur Verfügung.}
\end{tikzpicture}
\hspace{-1cm}
\begin{tikzpicture}
	\card
	\cardstrip
	\cardbanner{banner/white.png}
	\cardicon{banner/coin.png}
	\cardprice{5}
	\cardtitle{Mine}
	\cardcontent{\emph{Errata:} Der Kartentext ist falsch, es sollte \enquote{Du \emph{darfst} eine beliebige Geldkarte aus der Hand entsorgen. Nimm eine Geldkarte vom Vorrat auf die Hand, die bis zu \coin{3} mehr kostet.} heißen.
	\\
	\medskip
	\\
	Normalerweise entsorgst du ein Kupfer und nimmst dir dafür ein Silber, oder du entsorgst ein Silber und nimmst dir ein Gold. Du kannst dir aber auch eine gleichwertige oder billigere Karte nehmen. Die neue Karte nimmst du sofort auf die Hand und darfst sie noch während deines Zuges einsetzen. Wer keine Geldkarte zum Entsorgen hat, erhält keine neue Karte.}
\end{tikzpicture}
\hspace{-1cm}
\begin{tikzpicture}
	\card
	\cardstrip
	\cardbanner{banner/white.png}
	\cardicon{banner/coin.png}
	\cardprice{6}
	\cardtitle{Abenteurer}
	\cardcontent{Sollte dein Nachziehstapel während des Aufdeckens aufgebraucht werden, mische deinen Ablagestapel. Die bereits ufgedeckten Karten werden nicht mit gemischt, sondern bleiben zunächst offen liegen. Solltest du auch mit Hilfe des neuen Nachziehstapels nicht genug Geldkarten aufdecken, bekommst du nur diejenigen Geldkarten, die du aufgedeckt hast.}
\end{tikzpicture}
\hspace{-1cm}
\begin{tikzpicture}
	\card
	\cardstrip
	\cardbanner{banner/white.png}
	\cardicon{banner/coin.png}
	\cardprice{5}
	\cardtitle{Bibliothek}
	\cardcontent{Aktionskarten darfst du zur Seite legen, sobald du sie ziehst, musst dies aber nicht tun. Hast du bereits 7 oder mehr Karten auf der Hand, wenn du die \emph{BIBLIOTHEK} ausspielst, ziehst du keine Karte nach. Wenn dein Nachziehstapel während des Ziehens aufgebraucht ist, mischst du den Ablagestapel, mischst aber die zur Seite gelegten Aktionskarten nicht mit ein. Diese werden erst auf den Ablagestapel gelegt, sobald du 7 Karten auf der Hand hast. Sollten die Karten nicht reichen, ziehst du nur so viele Karten wie möglich.}
\end{tikzpicture}
\hspace{-1cm}
\begin{tikzpicture}
	\card
	\cardstrip
	\cardbanner{banner/white.png}
	\cardicon{banner/coin.png}
	\cardprice{4}
	\cardtitle{Bürokrat}
	\cardcontent{Ist dein Nachziehstapel aufgebraucht, wenn du diese Karte spielst, legst du die Silberkarte verdeckt ab. Sie bildet dann deinen Nachziehstapel. Das Gleiche gilt für alle Mitspieler, die eine Punktekarte verdeckt auf den Nachziehstapel legen müssen.}
\end{tikzpicture}
\hspace{-1cm}
\begin{tikzpicture}
	\card
	\cardstrip
	\cardbanner{banner/white.png}
	\cardicon{banner/coin.png}
	\cardprice{4}
	\cardtitle{Dieb}
	\cardcontent{Jeder Mitspieler legt die beiden aufgedeckten Karten zunächst offen vor sich ab. Wer nur noch 1 Karte im nachziehstapel hat, legt diese vor sich ab und mischt erst dann seinen Ablagestapel. Hat ein Spieler nach dem Mischen noch immer nicht genug Karten, deckt er nur so viele auf wie möglich. Hat ein Spieler 2 Geldkarten offen liegen, wählst du eine davon aus, die der Spieler entsorgen muss. Die andere Karte legt er auf seinen Ablagestapel. Hat ein Spieler 1 Geldkarte offen liegen, muss er diese entsorgen. Hat ein Spieler keine Geldkarte aufgedeckt, muss er keine Karte entsorgen. Von den auf diese Weise entsorgten Karten darfst du eine beliebige Anzahl nehmen.}
\end{tikzpicture}
\hspace{-1cm}
\begin{tikzpicture}
	\card
	\cardstrip
	\cardbanner{banner/white.png}
	\cardicon{banner/coin.png}
	\cardprice{4}
	\cardtitle{Festmahl}
	\cardcontent{Die nimmst dir eine beliebige Karte aus dem Vorrat, die höchstens \coin{5} kostet und legst sie sofort auf deinen Ablagestapel. Du darfst den Betrag weder durch weitere Geldkarten, Münzen oder zusätzliches Geld von anderen Aktionskarten erhöhen. Spielst du das \emph{FESTMAHL} direkt nach dem \emph{THRONSAAL}, erhältst du 2 Karten, obwohl du das \emph{FESTMAHL} nur einmal entsorgen kannst.}
\end{tikzpicture}
\hspace{-1cm}
\begin{tikzpicture}
	\card
	\cardbanner{banner/green.png}
	\cardicon{banner/coin.png}
	\cardprice{4}
	\cardtitle{Gärten}
	\cardcontent{DIese Karte ist die einzige Punktekarte unter den Königreichkarten. SIe hat bis zum Ende des Spiels keine Funktion. Bei der Wertung des Spiels erhält der Spieler, der die Karte in seinem Kartensatz (Nachziehstapel, Handkarten und Ablagestapel) hat, für jeweils 10 Karten einen Siegpunkt. Es wird imemr abgerundet, d.h. 39 Karten ergeben 3 Siegpunkte, ebenso wie 31 Karten 3 Spiegpunkte ergeben. Wer mehrere \emph{GÄRTEN} besitzt, erhält für jeden \emph{GARTEN} die entsprechende Anzahl an Siegpunkten.}
\end{tikzpicture}
\hspace{-1cm}
\begin{tikzpicture}
	\card
	\cardstrip
	\cardbanner{banner/white.png}
	\cardicon{banner/coin.png}
	\cardprice{4}
	\cardtitle{Geldverleiher}
	\cardcontent{\emph{Errata:} Der Kartentext ist falsch, es sollte \enquote{Du \emph{darfst} ein Kupfer aus der Hand entsorgen. Wenn du das tust: + \coin{3}.} heißen.
	\\
	\medskip
	\\
	Wenn du kein Kupfer zum entsorgen auf der Hand hast, erhältst du kein zusätzliches Geld für die Kaufphase.}
\end{tikzpicture}
\hspace{-1cm}
\begin{tikzpicture}
	\card
	\cardstrip
	\cardbanner{banner/white.png}
	\cardicon{banner/coin.png}
	\cardprice{5}
	\cardtitle{Hexe}
	\cardcontent{Wenn du die \emph{HEXE} spielst und nicht mehr genügend Fluchkarten vorrätig sind, werden diese im Uhrzeigersinn (beginnend bei deinem linken Nachbarn) verteilt. Die Mitspieler legen die Fluchkarten sofort auf ihren Ablagestapel. Du ziehst immer 2 Karten von deinem Nachziehstapel, auch wenn keine Fluchkarten mehr im Vorrat sind.}
\end{tikzpicture}
\hspace{-1cm}
\begin{tikzpicture}
	\card
	\cardstrip
	\cardbanner{banner/white.png}
	\cardicon{banner/coin.png}
	\cardprice{5}
	\cardtitle{Jahrmarkt}
	\cardcontent{Spielst du mehrer \emph{JAHRMÄRKTE} hintereinander, zählst du am besten laut mit, wie viele Aktionen du noch ausspielen darfst, damit du den Überblick behältst.}
\end{tikzpicture}
\hspace{-1cm}
\begin{tikzpicture}
	\card
	\cardstrip
	\cardbanner{banner/white.png}
	\cardicon{banner/coin.png}
	\cardprice{3}
	\cardtitle{Kanzler}
	\cardcontent{Legst du deinen Nachziehstapel auf deinen Ablagestapel, musst du dies tun, bevor du weitere Aktionen ausspielst oder zur Kaufphase übergehst. Du darfst deinen Nachziehstapel nicht einsehen, bevor du ihn ablegst.}
\end{tikzpicture}
\hspace{-1cm}
\begin{tikzpicture}
	\card
	\cardstrip
	\cardbanner{banner/white.png}
	\cardicon{banner/coin.png}
	\cardprice{2}
	\cardtitle{Kapelle}
	\cardcontent{Die ausgespielte \emph{KAPELLE} selber darf nicht entsorgt werden, da sie sich nicht mehr auf der Hand befindet. Weitere \emph{KAPELLEN} auf der Hand dürfen entsorgt werden.}
\end{tikzpicture}
\hspace{-1cm}
\begin{tikzpicture}
	\card
	\cardstrip
	\cardbanner{banner/white.png}
	\cardicon{banner/coin.png}
	\cardprice{5}
	\cardtitle{Laboratorium}
	\cardcontent{Du \emph{musst} zuerst zwei Karten vom Nachziehstapel auf die Hand nehmen. Dann \emph{darfst} du eine weiter Aktionskarte ausspielen.}
\end{tikzpicture}
\hspace{-1cm}
\begin{tikzpicture}
	\card
	\cardstrip
	\cardbanner{banner/white.png}
	\cardicon{banner/coin.png}
	\cardprice{5}
	\cardtitle{\scriptsize{Ratsversammlung}}
	\cardcontent{Jeder Spieler \emph{muss} eine Karte von seinem Nachziehstapel auf die Hand nehmen.}
\end{tikzpicture}
\hspace{-1cm}
\begin{tikzpicture}
	\card
	\cardstrip
	\cardbanner{banner/white.png}
	\cardicon{banner/coin.png}
	\cardprice{5}
	\cardtitle{Spion}
	\cardcontent{\emph{Zuerst} musst du eine Karte vom Nachziehstapel auf die Hand nehmen. Dann \emph{muss} jeder Spieler (auch du) die oberste Karte seine Nachziehstapels aufdecken. Du entscheidest für jeden Spieler, ob die aufgedeckte Karte auf den Nachziehstapel zurück- oder auf den Ablagestapel abgelegt werden soll. Spieler, deren Nachziehstapel aufgebraucht ist, mischen ihren Ablagestapel und legen ihn als neuen Nachziehstapel bereit. Nur wer weder einen Nachzieh- noch Ablagestapel hat, braucht keine Karte aufzudecken. Ist den Mitspielern die Reihenfolge des Aufdeckens wichtig, deckst du zuerst auf - die anderen Spieler folgen im Uhrzeigersinn.}
\end{tikzpicture}
\hspace{-1cm}
\begin{tikzpicture}
	\card
	\cardstrip
	\cardbanner{banner/white.png}
	\cardicon{banner/coin.png}
	\cardprice{4}
	\cardtitle{Thronsaal}
	\cardcontent{\emph{Errata:} Der Kartentext ist falsch, es sollte \enquote{Du \emph{darfst} eine beliebige Aktionskarte aus der Hand zweimal ausspielen.} heißen.
	\\
	\medskip
	\\
	Wähle eine Aktionskarte aus deiner Hand und spiele sie zweimal aus, d.h. du legst die Aktionskarte aus, führst die Anweisungen der Karte komplett aus, nimmst die Karte zurück auf die Hand, legst sie noch einmal aus und führst die Anweisungen erneut aus. Für das doppelte ausspielen dieser Aktionskarte muss der Spieler keine zusätzliche Aktion (+1 Aktion) zur Verfügung haben -- sie ist sozusagen \enquote{kostenlos}. Legst du zwei \emph{THRONSAAL}-Karten aus, darfst du zuerst eine Aktion doppelt ausführen und dann eine andere Aktion ebenfalls doppelt ausführen. Du darfst aber nicht ein und und dieselbe Aktion vier mal ausführen. 
	\\
	\medskip
	\\
	Erlaubt die doppelt ausgespielte Karte +1 Aktion (z.B. der \emph{MARKT}), hast du nach der vollständigen Ausführung des \emph{THRONSAALS} zwei weitere Aktionen zur Verfügung. Hättest du zwei \emph{MARKT}-Karten regulär hintereinander ausgespielt, bleibt dir nur noch eine zusätzliche Aktion zur Verfügung, da das Ausspielen der zweiten Marktkarte schon die zusätzliche Aktion der ersten Karte aufgebraucht hätte. Beim \emph{THRONSAAL} ist es besonders wichtig, laut die verbleibenden Anzahlen an Aktionen mitzuzählen. Du darfst keine weitere Aktion ausspielen, bevor der \emph{THRONSAAL} komplett abgearbeitet ist.}
\end{tikzpicture}
\hspace{-1cm}
\begin{tikzpicture}
	\card
	\cardstrip
	\cardbanner{banner/white.png}
	\cardicon{}
	\cardprice{}
	\cardtitle{\scriptsize{Empfohlene 10er Sätze\qquad}}
	\cardcontent{\emph{Erstes Spiel:}
		\\
		Burggraben, Dorf, Holzfäller, Keller, Markt, Miliz, Mine, Schmiede, Umbau, Werkstatt
		\\
		\smallskip
		\\
		\emph{Großes Geld:}
		\\
		Abenteurer, Bürokrat, Festmahl, Geldverleiher, Kanzler, Kapelle, Laboratorium, Mark, Mine, Thronsaal
		\\
		\smallskip
		\\
		\emph{Interaktion:}
		\\
		Bibliothek, Burggraben, Bürokrat, Dieb, Dorf, Jahrmarkt, Kanzler, Miliz, Ratsversammlung, Spion
		\\
		\smallskip
		\\
		\emph{Im Wandel:}
		\\
		Dieb, Dorf, Festmahl, Gärten, Hexe, Holzfäller, Kapelle, Keller, Laboratorium, Werkstatt
		\\
		\smallskip
		\\
		\emph{Dorfplatz:}
		\\
		Bibliothek, Bürokrat, Dorf, Holzfäller, Jahrmarkt, Keller, Markt, Schmiede, Thronsaal, Umbau
		\\
	}
\end{tikzpicture}
\hspace{1cm}
	    % Basic settings for this card set
\renewcommand{\cardcolor}{basicgame}
\renewcommand{\cardextension}{Edition II}
\renewcommand{\cardextensiontitle}{Das Basisspiel}

\clearpage
\newpage
\section{\cardextension \ - \cardextensiontitle}

\begin{tikzpicture}
	\card
	\cardstrip
	\cardbanner{banner/blue.png}
	\cardicon{banner/coin.png}
	\cardprice{2}
	\cardtitle{Burggraben}
	\cardcontent{Spielt ein anderer Spieler eine Angriffskarte (mit der Aufschrift AKTION -- ANGRIFF), kannst du die Karte \emph{BURGGRABEN} vorzeigen, falls du sie in diesem Moment auf der Hand hast. In diesem Fall bist \emph{du} von der Auswirkung des Angriffst nicht betroffen, d.h. du musst mit der \emph{HEXE} keine Fluchkarte nehmen usw. Haben mehrer Spieler einen \emph{BURGGRABEN} auf der Hand, dürfen diese auch eingesetzt und vorgezeigt werden. Danach nehmen die Spieler ihre Karte zurück auf die Hand.
	\\
	\medskip
	\\
	Der Spieler, der den Angriff gespielt hat, darf unabhängig davon, ob ein oder mehrere \emph{BURGGRÄBEN} gespielt wurden, die weiteren Anweisungen seiner Aktionskarte ausführen. Der \emph{BURGGRABEN} darf auch in der eigenen Aktionsphase gespielt werden -- dann ziehst du 2 Karten nach.}
\end{tikzpicture}
\hspace{-1cm}
\begin{tikzpicture}
	\card
	\cardstrip
	\cardbanner{banner/white.png}
	\cardicon{banner/coin.png}
	\cardprice{2}
	\cardtitle{Kapelle}
	\cardcontent{Die ausgespielte \emph{KAPELLE} selber darf nicht entsorgt werden, da sie sich nicht mehr auf der Hand befindet. Weitere \emph{KAPELLEN} auf der Hand dürfen entsorgt werden.}
\end{tikzpicture}
\hspace{-1cm}
\begin{tikzpicture}
	\card
	\cardstrip
	\cardbanner{banner/white.png}
	\cardicon{banner/coin.png}
	\cardprice{2}
	\cardtitle{Keller}
	\cardcontent{Der ausgespielte \emph{KELLER} selber darf nicht abgelegt werden, da er sich nicht mehr in deiner Hand befindet. Sage an wie viele Karten du ablegst und lege diese auf deinen Ablagestapel. Danach ziehst du die gleiche Anzahl Karten vom Nachziehstapel. Sollte während dieses Vorgangs der Nachziehstapel aufgebraucht werden, wird dein Ablagestapel zusammen mit den soeben abgelegten Karten gemischt und als neuer Nachziehstapel bereitgelegt.}
\end{tikzpicture}
\hspace{-1cm}
\begin{tikzpicture}
	\card
	\cardstrip
	\cardbanner{banner/white.png}
	\cardicon{banner/coin.png}
	\cardprice{3}
	\cardtitle{Dorf}
	\cardcontent{Spielst du mehrere \emph{DÖRFER} hintereinander, zählst du am besten laut mit, wie viele Aktionen du noch spielen darfst, damit du den Überblick behältst.}
\end{tikzpicture}
\hspace{-1cm}
\begin{tikzpicture}
	\card
	\cardstrip
	\cardbanner{banner/white.png}
	\cardicon{banner/coin.png}
	\cardprice{3}
	\cardtitle{Händlerin}
	\cardcontent{DIe ziehst 1 Karte und erhältst + 1 Akton. Wenn du in diesen Zug vor dem Ausspeilen dieser \emph{HÄNDLERIN} noch kein Silber ausgespielt hast, erhältst du für das erste danach ausgespielte Silber +\coin{1}. Für jedes weitere ausgespielte Silber erhältst du keinen zusätzlichen Bonus. Hast du mherer \emph{HÄNDLERINNEN} ausgespielt, erhältst du pro \emph{HÄNDLERIN} +\coin{1}.}
\end{tikzpicture}
\hspace{-1cm}
\begin{tikzpicture}
	\card
	\cardstrip
	\cardbanner{banner/white.png}
	\cardicon{banner/coin.png}
	\cardprice{3}
	\cardtitle{Vasall}
	\cardcontent{Ist die aufgedeckte KArte eine Aktionskarte (auch ggf. kombineirte), \emph{darfst} du sie sofort ausspielen. Wenn du sie ausspielst, legst du sie in deinen SPielebreich und führst sofort die Anweisung darauf aus. Dafür benötigst du keine zusätzliche Aktion. Das Ausspielen der AKtionskarte verbraucht auch keien freie oder zusätzliche Aktion, die du durch das Ausspielen anderer Karten bereits gesammelt hast.}
\end{tikzpicture}
\hspace{-1cm}
\begin{tikzpicture}
	\card
	\cardstrip
	\cardbanner{banner/white.png}
	\cardicon{banner/coin.png}
	\cardprice{3}
	\cardtitle{Vorbotin}
	\cardcontent{Du ziehst 1 Karte und erhälts + 1 Aktion. Schau dir deinen Ablagestapel an. DU \emph{darfst} eine KArte daraus auswählen und oben auf deinen nachziehstapel legen. Die restlichen Karten (oder alle) legst du in beliebiger Reihenfolge zurück auf den Ablagestapel. Ist dein Ablagestapel leer, passiert nichts.}
\end{tikzpicture}
\hspace{-1cm}
\begin{tikzpicture}
	\card
	\cardstrip
	\cardbanner{banner/white.png}
	\cardicon{banner/coin.png}
	\cardprice{3}
	\cardtitle{Werkstatt}
	\cardcontent{Nimm dir eine Karte aus dem Vorrat und lege diese sofort auf deinen Ablagestapel. Du kansnt weder Geldkarten noch zusätzlich über Aktionskarten erhaltenes Geld oder Münzen (bei Erweiterungen mit Münzen) einsetzen, um den angegebenen Betrag auf der Karte zu erhöhen.}
\end{tikzpicture}
\hspace{-1cm}
\begin{tikzpicture}
	\card
	\cardstrip
	\cardbanner{banner/white.png}
	\cardicon{banner/coin.png}
	\cardprice{4}
	\cardtitle{Bürokrat}
	\cardcontent{Ist dein Nachziehstapel aufgebraucht, wenn du diese Karte spielst, legst du die Silberkarte verdeckt ab. Sie bildet dann deinen Nachziehstapel. Das Gleiche gilt für alle Mitspieler, die eine Punktekarte verdeckt auf den Nachziehstapel legen müssen.}
\end{tikzpicture}
\hspace{-1cm}
\begin{tikzpicture}
	\card
	\cardstrip
	\cardbanner{banner/white.png}
	\cardicon{banner/coin.png}
	\cardprice{4}
	\cardtitle{Geldverleiher}
	\cardcontent{\emph{Errata:} Der Kartentext ist falsch, es sollte \enquote{Du \emph{darfst} ein Kupfer aus der Hand entsorgen. Wenn du das tust: + \coin{3}.} heißen.
	\\
	\medskip
	\\
	Wenn du kein Kupfer zum entsorgen auf der Hand hast, erhältst du kein zusätzliches Geld für die Kaufphase.}
\end{tikzpicture}
\hspace{-1cm}
\begin{tikzpicture}
	\card
	\cardstrip
	\cardbanner{banner/white.png}
	\cardicon{banner/coin.png}
	\cardprice{4}
	\cardtitle{Miliz}
	\cardcontent{Deine Mitspieler müssen Karten aus ihrer Hand ablegen, bis sie nur noch 3 Karten auf der Hand haben. Spieler, die zum Zeitpunkt des Angriffs bereits 3 oder weniger Karten auf der Hand haben, müssen keine weiteren Karten ablegen.}
\end{tikzpicture}
\hspace{-1cm}
\begin{tikzpicture}
	\card
	\cardstrip
	\cardbanner{banner/white.png}
	\cardicon{banner/coin.png}
	\cardprice{4}
	\cardtitle{Schmiede}
	\cardcontent{Du musst 3 Karten vom Nachziehstapel ziehen und auf die Hand nehmen.}
\end{tikzpicture}
\hspace{-1cm}
\begin{tikzpicture}
	\card
	\cardstrip
	\cardbanner{banner/white.png}
	\cardicon{banner/coin.png}
	\cardprice{4}
	\cardtitle{Thronsaal}
	\cardcontent{\emph{Errata:} Der Kartentext ist falsch, es sollte \enquote{Du \emph{darfst} eine beliebige Aktionskarte aus der Hand zweimal ausspielen.} heißen.
	\\
	\medskip
	\\
	Wähle eine Aktionskarte aus deiner Hand und spiele sie zweimal aus, d.h. du legst die AKtionskarte aus, führst die Anweisungen der Karte komplett aus, nimmst die Karte zurück auf die Hand, legst sie noch einmal aus und führst die Anweisungen erneut aus. Für das doppelte ausspielen dieser Aktionskarte muss der Spieler keine zusätzliche Aktion (+1 Aktion) zur Verfügung haben -- sie ist sozusagen \enquote{kostenlos}. Legst du zwei \emph{THRONSAAL}-Karten aus, darfst du zuerst eine Aktion doppelt ausführen und dann eine andere Aktion ebenfalls doppelt ausführen. Du darfst aber nicht ein und und dieselbe Aktion vier mal ausführen. 
	\\
	\medskip
	\\
	Erlaubt die doppelt ausgespielte Karte +1 Aktion (z.B. der \emph{MARKT}), hast du nach der vollständigen Ausführung des \emph{THRONSAALS} zwei weitere Aktionen zur Verfügung. Hättest du zwei \emph{MARKT}-Karten regulär hintereinander ausgespielt, bleibt dir nur noch eine zusätzliche Aktion zur Verfügung, da das Ausspielen der zweiten Marktkarte schon die zusätzliche Aktion der ersten Karte aufgebraucht hätte. Beim \emph{THRONSAAL} ist es besonders wichtig, laut die verbleibenden Anzahlen an Aktionen mitzuzählen. Du darfst keine weitere Aktion ausspielen, bevor der \emph{THRONSAAL} komplett abgearbeitet ist.}
\end{tikzpicture}
\hspace{-1cm}
\begin{tikzpicture}
	\card
	\cardstrip
	\cardbanner{banner/white.png}
	\cardicon{banner/coin.png}
	\cardprice{4}
	\cardtitle{Umbau}
	\cardcontent{Der ausgespielte \emph{UMBAU} selber darf nicht entsorgt werden, da er sich nicht mehr in der Hand befindet. Weitere \emph{UMBAU}-Karten in deiner Hand dürfen entsorgt werden. Wenn du keine Karte zum Entsorgen auf der Hand hast, darfst du dir auch keine neue Karte nehmen. Die neue Karte, die du dir nimmst, darf maximal bis zu \coin{2} mehr als die entsorgte Karte kosten. Der Betrag darf weder durch weitere Geldkarten, Münzen oder zusätzliches Geld von anderen Aktionskarten erhöht werden. Die neue Karte kann die gleiche Karte sein, wie die, die du entsorgt hast. Lege die neue Karte auf deinen Ablagestapel.}
\end{tikzpicture}
\hspace{-1cm}
\begin{tikzpicture}
	\card
	\cardstrip
	\cardbanner{banner/white.png}
	\cardicon{banner/coin.png}
	\cardprice{4}
	\cardtitle{Wilddiebin}
	\cardcontent{Du ziehst 1 Karte, erhältst + 1 Aktion und +\coin{1}. Dann schaust du, wie viele Vorratsstapel (Fluch-, Geld-, Punkte- und Aktionskarten) bereits leer sind. Ist kein Stapel leer, musst du keine Handkarten ablegen. Ist ein Stapel leer, legst du 1 Handkarte ab usw. Wenn du nicht so viele Karten auf der Hand hast, wie Vorratsstapel leer sind, legst du so viele Karten ab, wie du kannst.}
\end{tikzpicture}
\hspace{-1cm}
\begin{tikzpicture}
	\card
	\cardstrip
	\cardbanner{banner/white.png}
	\cardicon{banner/coin.png}
	\cardprice{5}
	\cardtitle{Banditin}
	\cardcontent{Zuerst nimmst du ein Gold vom Vorrat und legst es auf deinen Ablagestapel. Dann deckt jeder Mitspieler -- beginnend bei deinem linken Mitspieler -- die obersten zwei Karten seine Nachziehstapels auf. Deckt ein SPieler zwei Geldkarten (auch ggf. kombiniert) außer Kupfer aus, muss er \emph{eine} davon entsorgen. Dabei darf er selbst entscheiden, welche Geldkarte er entsorgt. Die andere Geldkarte wird -- genauso wie alle anderen Karten -- abgelegt. Deckt ein Spieler eine Geldkarte außer Kupfer sowie eine andere Karte (z.B. ein Kupfer oder eine beliebige Aktionskarte) auf, wird diese Geldkarte entsorgt. Die andere aufgedeckte Karte wird abgelegt.}
\end{tikzpicture}
\hspace{-1cm}
\begin{tikzpicture}
	\card
	\cardstrip
	\cardbanner{banner/white.png}
	\cardicon{banner/coin.png}
	\cardprice{5}
	\cardtitle{Bibliothek}
	\cardcontent{Aktionskarten darfst du zur Seite legen, sobald du sie ziehst, musst dies aber nicht tun. Hast du bereits 7 oder mehr Karten auf der Hand, wenn du die \emph{BIBLIOTHEK} ausspielst, ziehst du keine Karte nach. Wenn dein Nachziehstapel während des Ziehens aufgebraucht ist, mischst du den Ablagestapel, mischst aber die zur Seite gelegten Aktionskarten nicht mit ein. Diese werden erst auf den Ablagestapel gelegt, sobald du 7 Karten auf der Hand hast. Sollten die Karten nicht reichen, ziehst du nur so viele Karten wie möglich.}
\end{tikzpicture}
\hspace{-1cm}
\begin{tikzpicture}
	\card
	\cardstrip
	\cardbanner{banner/white.png}
	\cardicon{banner/coin.png}
	\cardprice{5}
	\cardtitle{Hexe}
	\cardcontent{Wenn du die \emph{HEXE} spielst und nicht mehr genügend Fluchkarten vorrätig sind, werden diese im Uhrzeigersinn (beginnend bei deinem linken Nachbarn) verteilt. Die Mitspieler legen die Fluchkarten sofort auf ihren Ablagestapel. Du ziehst immer 2 Karten von deinem Nachziehstapel, auch wenn keine Fluchkarten mehr im Vorrat sind.}
\end{tikzpicture}
\hspace{-1cm}
\begin{tikzpicture}
	\card
	\cardstrip
	\cardbanner{banner/white.png}
	\cardicon{banner/coin.png}
	\cardprice{5}
	\cardtitle{Jahrmarkt}
	\cardcontent{Spielst du mehrer \emph{JAHRMÄRKTE} hintereinander, zählst du am besten laut mit, wie viele Aktionen du noch ausspielen darfst, damit du den Überblick behältst.}
\end{tikzpicture}
\hspace{-1cm}
\begin{tikzpicture}
	\card
	\cardstrip
	\cardbanner{banner/white.png}
	\cardicon{banner/coin.png}
	\cardprice{5}
	\cardtitle{Laboratorium}
	\cardcontent{Du \emph{musst} zuerst zwei Karten vom Nachziehstapel auf die Hand nehmen. Dann \emph{darfst} du eine weiter Aktionskarte ausspielen.}
\end{tikzpicture}
\hspace{-1cm}
\begin{tikzpicture}
	\card
	\cardstrip
	\cardbanner{banner/white.png}
	\cardicon{banner/coin.png}
	\cardprice{5}
	\cardtitle{Markt}
	\cardcontent{Du musst eine Karte vom Nachziehstapel auf die Hand nehmen. Du \emph{darfst} in der Aktionsphase eine weiter Aktionskarte ausspielen. Du \emph{darfst} in der Kaufphase einen zusätzlichen Kauf tätigen und hast dafür ein zusätzliches Geld zur Verfügung.}
\end{tikzpicture}
\hspace{-1cm}
\begin{tikzpicture}
	\card
	\cardstrip
	\cardbanner{banner/white.png}
	\cardicon{banner/coin.png}
	\cardprice{5}
	\cardtitle{Mine}
	\cardcontent{\emph{Errata:} Der Kartentext ist falsch, es sollte \enquote{Du \emph{darfst} eine beliebige Geldkarte aus der Hand entsorgen. Nimm eine Geldkarte vom Vorrat auf die Hand, die bis zu \coin{3} mehr kostet.} heißen.
	\\
	\medskip
	\\
	Normalerweise entsorgst du ein Kupfer und nimmst dir dafür ein Silber, oder du entsorgst ein Silber und nimmst dir ein Gold. Du kannst dir aber auch eine gleichwertige oder billigere Karte nehmen. Die neue Karte nimmst du sofort auf die Hand und darfst sie noch während deines Zuges einsetzen. Wer keine Geldkarte zum Entsorgen hat, erhält keine neue Karte.}
\end{tikzpicture}
\hspace{-1cm}
\begin{tikzpicture}
	\card
	\cardstrip
	\cardbanner{banner/white.png}
	\cardicon{banner/coin.png}
	\cardprice{5}
	\cardtitle{\scriptsize{Ratsversammlung}}
	\cardcontent{Jeder Spieler \emph{muss} eine Karte von seinem Nachziehstapel auf die Hand nehmen.}
\end{tikzpicture}
\hspace{-1cm}
\begin{tikzpicture}
	\card
	\cardstrip
	\cardbanner{banner/white.png}
	\cardicon{banner/coin.png}
	\cardprice{5}
	\cardtitle{Torwächterin}
	\cardcontent{Du ziehst 1 Karte und erhältst + 1 Aktion. Dann ziehst du dir die obersten 2 Karten deines Nachziehstapels an. Du kannst beide Karten entsorgen, beide Karten ablegen oder sie in beliebiger Reihenfolge zurück auf den Nachziehstapel legen. Du kannst aber auch eine entsorgten und eine ablegen, oder eine entsorgen und die anderen zurück auf den Nachziehstapel legen, oder eine ablegen und die anderen zurücklegen.}
\end{tikzpicture}
\hspace{-1cm}
\begin{tikzpicture}
	\card
	\cardstrip
	\cardbanner{banner/white.png}
	\cardicon{banner/coin.png}
	\cardprice{6}
	\cardtitle{Töpferei}
	\cardcontent{Nimm eine Karte vom Vorrat, die zu diesem Zeitpunkt maximal \coin{5} kostet. Du darfst kein zusätzliches \coin{\enspace} einsetzen, um dir eine teurere Karte zu nehmen. Außer \coin{\enspace} darf die Karte keine zusätzlichen Kosten enthalten.
	\\
	\medskip
	\\
	DU darfst dir zum Beispiel keine Karte mit \potion (aus Alchemie) oder \hex (aus Empires) in den Kosten nehmen. Die genommene Karte nimmst du direkt auf die Hand. Anschließend legst du eine beliebige Handkarte (das kann die gerade genommene oder eine andere sein) oben auf den Nachziehstapel.}
\end{tikzpicture}
\hspace{-1cm}
\begin{tikzpicture}
	\card
	\cardbanner{banner/green.png}
	\cardicon{banner/coin.png}
	\cardprice{4}
	\cardtitle{Gärten}
	\cardcontent{DIese Karte ist die einzige Punktekarte unter den Königreichkarten. SIe hat bis zum Ende des Spiels keine Funktion. Bei der Wertung des Spiels erhält der Spieler, der die Karte in seinem Kartensatz (Nachziehstapel, Handkarten und Ablagestapel) hat, für jeweils 10 Karten einen Siegpunkt. Es wird imemr abgerundet, d.h. 39 Karten ergeben 3 Siegpunkte, ebenso wie 31 Karten 3 Spiegpunkte ergeben. Wer mehrere \emph{GÄRTEN} besitzt, erhält für jeden \emph{GARTEN} die entsprechende Anzahl an Siegpunkten.}
\end{tikzpicture}
\hspace{-1cm}
\begin{tikzpicture}
	\card
	\cardstrip
	\cardbanner{banner/white.png}
	\cardicon{}
	\cardprice{}
	\cardtitle{\scriptsize{Empfohlene 10er Sätze\qquad}}
	\cardcontent{\emph{Erstes Spiel:}
		\\
		Burggraben, Dorf, Händlerin, Keller, Markt, Miliz, Mine, Schmiede, Umbau, Werkstatt
		\\
		\smallskip
		\\
		\emph{Verzerrte Größen:}
		\\
		Banditin, Bürokrat, Gärtner, Hexe, Jahrmarkt, Kapelle, Thronsaal, Töpferei, Torwächterin, Werkstatt
		\\
		\smallskip
		\\
		\emph{Schleichweg:}
		\\
		Bürokrat, Dorf, Geldverleiher, Laboratorium, Jahrmarkt, Ratsversammlung, Töpferei, Torwächterin, Vasall, Vorbotin
		\\
		\smallskip
		\\
		\emph{Kunststück:}
		\\
		Bibliothek, Gärten, Jahrmarkt, Keller, Miliz, Ratsversammlung, Schmiede, Thronsaal, Vorbotin, Wilddiebin
		\\
		\smallskip
		\\
		\emph{Verbesserungen:}
		\\
		Burggraben, Geldverleiher, Händlerin, Hexe, Keller, Markt, Mine, Töpferei, Umbau, Wilddiebin
		\\
		\smallskip
		\\
		\emph{Silber \& Gold:}
		\\
		Bandit, Bürokrat, Geldverleiher, Händlerin, Kapelle, Laboratorium, Mine, Thronsaal, Vasall, Vorbotin
		\\
	}
\end{tikzpicture}
\hspace{1cm}
	    % Basic settings for this card set
\renewcommand{\cardcolor}{intrigue}
\renewcommand{\cardextension}{Edition I}
\renewcommand{\cardextensiontitle}{Die Intrige}

\clearpage
\newpage
\section{\cardextension \ - \cardextensiontitle}

\begin{tikzpicture}
	\card
	\cardstrip
	\cardbanner{banner/white.png}
	\cardicon{banner/coin.png}
	\cardprice{2}
	\cardtitle{Burghof}
	\cardcontent{Du ziehst 3 Karten von deinem Nachziehstapel und nimmst sie auf die Hand. Dann wählst du eine beliebige Karte aus deiner Hand und legst sie verdeckt auf den Nachziehstapel.}
\end{tikzpicture}
\hspace{-1cm}
\begin{tikzpicture}
	\card
	\cardstrip
	\cardbanner{banner/blue.png}
	\cardicon{banner/coin.png}
	\cardprice{2}
	\cardtitle{Geheimkammer}
	\cardcontent{Wenn du diese Karte in deiner eigenen Aktionsphase ausspielst, legst du eine beliebige Anzahl Karten (auch 0 Karten) aus deiner Hand ab. \emph{Danach} erhältst du +\coin{1} pro abgelegte Karte. 
	\\
	\medskip
	\\
	Spielt ein anderer Spieler eine Angriffskarte aus, darfst du diese Karte vorzeigen, wenn du sie in diesem Moment auf der Hand hast, auch wenn dich der Angriff nicht betrifft. Wenn du das tust, ziehst du zuerst 2 Karten nach und legst dann 2 beliebige Karten aus deiner Hand verdeckt zurück auf den Nachziehstapel. Du darfst auch die \emph{GEHEIMKAMMER} selbst auf den Nachziehstapel legen, da sich die Karte auch nach dem Vorzeigen auf deiner Hand befindet. Du darfst pro Angriff so viele Reaktionskarten vorzeigen, wie du möchtest. So kannst du zum Beispiel zuerst die \emph{GEHEIMKAMMER} abwickeln und danach einen \emph{BURGGRABEN} vorzeigen, um den Angriff abzuwehren. Die \emph{GEHEIMKAMMER} selbst wehrt einen Angriff nicht ab.}
\end{tikzpicture}
\hspace{-1cm}
\begin{tikzpicture}
	\card
	\cardstrip
	\cardbanner{banner/white.png}
	\cardicon{banner/coin.png}
	\cardprice{2}
	\cardtitle{Handlanger}
	\cardcontent{Du darfst von den vier Anweisungen der Karte \emph{genau 2} auswählen und diese für deinen Zug nutzen. Du musst zwei verschiedene Anweisungen wählen und darfst nicht z. B. zwei Karten ziehen oder 2 zusätzliche Käufe tätigen. Du musst dich sofort entscheiden, welche zwei Anweisungen du nutzen möchtest. Du darfst nicht eine Karte nachziehen und dich erst dann entscheiden, welche zweite Anweisung du ausführen möchtest.}
\end{tikzpicture}
\hspace{-1cm}
\begin{tikzpicture}
	\card
	\cardstrip
	\cardbanner{banner/white.png}
	\cardicon{banner/coin.png}
	\cardprice{3}
	\cardtitle{Armenviertel}
	\cardcontent{Du darfst in deiner Aktionsphase zwei zusätzliche Aktionen ausführen. Zeige deine Kartenhand vor. Wenn du keine Aktionskarte (oder Aktions-/Punktekarte) auf der Hand hast, ziehst du zwei Karten nach. Sollten sich darunter Aktionskarten befinden, darfst du diese auch gleich nutzen.}
\end{tikzpicture}
\hspace{-1cm}
\begin{tikzpicture}
	\card
	\cardstrip
	\cardbanner{banner/whitegreen.png}
	\cardicon{banner/coin.png}
	\cardprice{3}
	\cardtitle{Große Halle}
	\cardcontent{Diese Karte ist eine kombinierte Aktions- und Punktekarte. Sie kann in der Aktionsphase eingesetzt werden und bringt \emph{zusätzlich} bei Spielende 1 Siegpunkt. Wenn du diese Karte ausspielst, ziehst du eine Karte und darfst eine zusätzliche Aktion ausspielen.}
\end{tikzpicture}
\hspace{-1cm}
\begin{tikzpicture}
	\card
	\cardstrip
	\cardbanner{banner/white.png}
	\cardicon{banner/coin.png}
	\cardprice{3}
	\cardtitle{Maskerade}
	\cardcontent{Ziehe zwei Karten. Anschließend wählen alle Spieler (einschließlich dir selbst) eine beliebige Karte aus ihrer Hand und legen sie verdeckt neben ihrem linken Nachbarn ab. Wenn alle Karten verteilt sind, nimmt jeder Spieler die erhaltene Karte auf die Hand. Da die \emph{MASKERADE} keine Angriffskarte ist, dürfen die anderen Spieler keine Reaktionskarten vorzeigen. Danach darfst du eine Karte aus deiner Hand entsorgen.}
\end{tikzpicture}
\hspace{-1cm}
\begin{tikzpicture}
	\card
	\cardstrip
	\cardbanner{banner/white.png}
	\cardicon{banner/coin.png}
	\cardprice{3}
	\cardtitle{Trickser}
	\cardcontent{Alle Mitspieler müssen die oberste Karte ihres Nachziehstapels aufdecken und diese entsorgen. Du wählst für jeden Mitspieler jeweils eine Karte aus dem Vorrat, die genauso viel kostet, wie die entsorgte Karte und gibst sie dem Mitspieler. Dieser legt die neue Karte auf seinen Ablagestapel. Befindet sich keine Karte mit den gleichen Kosten im Vorrat, muss der Mitspieler seine Karte trotzdem entsorgen, erhält dafür aber keine neue Karte. Du darfst dem Mitspieler auch eine gleiche Karte wie die, die er entsorgt hat, zurückgeben. }
\end{tikzpicture}
\hspace{-1cm}
\begin{tikzpicture}
	\card
	\cardstrip
	\cardbanner{banner/white.png}
	\cardicon{banner/coin.png}
	\cardprice{3}
	\cardtitle{Verwalter}
	\cardcontent{Du wählst von den drei Anweisungen der Karte \emph{genau 1} aus und führst diese dann komplett aus. Wenn du dich entscheidest, zwei Karten zu entsorgen, du aber nur eine Karte auf der Hand hast, musst du diese entsorgen.}
\end{tikzpicture}
\hspace{-1cm}
\begin{tikzpicture}
	\card
	\cardstrip
	\cardbanner{banner/white.png}
	\cardicon{banner/coin.png}
	\cardprice{3}
	\cardtitle{Wunschbrunnen}
	\cardcontent{Zuerst ziehst du eine Karte. Du darfst in der Aktionsphase eine zusätzliche Aktion ausführen. Nenne eine Karte (z.B. \emph{KUPFER}) und decke die oberste Karte deines Nachziehstapels auf. Handelt es sich dabei um die von dir genannte Karte, nimmst du sie auf die Hand. Ansonsten legst du sie zurück auf den Nachziehstapel.}
\end{tikzpicture}
\hspace{-1cm}
\begin{tikzpicture}
	\card
	\cardstrip
	\cardbanner{banner/white.png}
	\cardicon{banner/coin.png}
	\cardprice{4}
	\cardtitle{Baron}
	\cardcontent{Du \emph{darfst} ein Anwesen (sofern du gerade eins auf der Hand hast) ablegen und erhältst dafür +\coin{4} für die Kaufphase. Wenn du das nicht tun kannst (weil du kein Anwesen auf der Hand hast) oder willst, musst du dir ein Anwesen nehmen, solange noch welche im Vorrat sind.}
\end{tikzpicture}
\hspace{-1cm}
\begin{tikzpicture}
	\card
	\cardstrip
	\cardbanner{banner/white.png}
	\cardicon{banner/coin.png}
	\cardprice{4}
	\cardtitle{Bergwerk}
	\cardcontent{Du ziehst zuerst eine Karte nach und \emph{darfst} dann diese Karte entsorgen, bevor du ggf. weitere Aktionen ausspielst. Du erhältst dafür +\coin{2}. Wenn du das \emph{BERGWERK} auf einen \emph{THRONSAAL} spielst, erhältst du den Bonus für das Entsorgen nur einmal, da du die Karte nur einmal entsorgen kannst. Die anderen Anweisungen (+ 1 Karte sowie + 2 Aktionen) werden durch den \emph{THRONSAAL} dagegen verdoppelt.}
\end{tikzpicture}
\hspace{-1cm}
\begin{tikzpicture}
	\card
	\cardstrip
	\cardbanner{banner/white.png}
	\cardicon{banner/coin.png}
	\cardprice{4}
	\cardtitle{Brücke}
	\cardcontent{Die Kosten aller Karten (auch Handkarten, Karten aus den Nachzieh- und Ablagestapeln) werden in diesem Spielzug für alle Belange um \coin{1} reduziert (nicht aber unter \coin{0}). Dieser Effekt ist kumulativ, d. h. die Kosten pro Karte können durch das Ausspielen bestimmter Karten (z. B. den \emph{THRONSAAL}) auch um \coin{2} oder mehr reduziert werden.}
\end{tikzpicture}
\hspace{-1cm}
\begin{tikzpicture}
	\card
	\cardstrip
	\cardbanner{banner/white.png}
	\cardicon{banner/coin.png}
	\cardprice{4}
	\cardtitle{Eisenhütte}
	\cardcontent{Du nimmst dir eine beliebige Karte vom Vorrat, die maximal \coin{4} kostet. Durch das Ausspielen bestimmter Aktionskarten (z. B. die \emph{BRÜCKE}) können die Kosten der Karten reduziert werden.
	\\
	\medskip
	\\
	Je nachdem, ob du dich für eine Aktions-, Geld- oder Punktekarte entschieden hast, erhältst du einen anderen Bonus. Solltest du dich für eine kombinierte Karte entscheiden, erhältst du die Boni beider Kartentypen. }
\end{tikzpicture}
\hspace{-1cm}
\begin{tikzpicture}
	\card
	\cardstrip
	\cardbanner{banner/white.png}
	\cardicon{banner/coin.png}
	\cardprice{4}
	\cardtitle{Kupferschmied}
	\cardcontent{Mit dieser Karte erhöhst du den Wert aller in diesem Zug gespielten \emph{KUPFER} um +\coin{1}. Der Effekt ist kumulativ, d. h. durch das Ausspielen anderer Aktionskarten (z. B. den \emph{THRONSAAL} oder einen weiteren \emph{KUPFERSCHMIED}) kann der Wert weiter erhöht werden.}
\end{tikzpicture}
\hspace{-1cm}
\begin{tikzpicture}
	\card
	\cardstrip
	\cardbanner{banner/white.png}
	\cardicon{banner/coin.png}
	\cardprice{4}
	\cardtitle{Späher}
	\cardcontent{Sollten für das Aufdecken der vier Karten nicht genügend Karten im Nachziehstapel sein, ziehst du zunächst die restlichen Karten und mischst dann deinen Ablagestapel neu, ohne die bereits aufgedeckten Karten mit einzumischen. Sollten auch dann nicht genügend Karten zur Verfügung stehen, ziehst du nur so viele Karten wie möglich. Du musst alle Punktekarten auf die Hand nehmen, die restlichen Karten legst du in beliebiger Reihenfolge auf den Nachziehstapel. Diese musst du deinen Mitspielern nicht zeigen. Kombinierte Aktions-/Punktekarten sind auch Punktekarten. }
\end{tikzpicture}
\hspace{-1cm}
\begin{tikzpicture}
	\card
	\cardstrip
	\cardbanner{banner/white.png}
	\cardicon{banner/coin.png}
	\cardprice{4}
	\cardtitle{Verschwörer}
	\cardcontent{Wenn du zu dem Zeitpunkt an dem du den \emph{VERSCHWÖRER} spielst, bereits mindestens 3 Aktionskarten (inklusive diesem \emph{VERSCHWÖRER}) ausgespielt hast, erhältst du den Bonus. Wenn du erst im weiteren Verlauf deiner Aktionsphase diese Bedingung erfüllst, erhältst du den Bonus nicht. Aktionskarten, die z. B. durch den \emph{THRONSAAL} doppelt ausgespielt werden dürfen, gelten als 2 ausgespielte Karten.}
\end{tikzpicture}
\hspace{-1cm}
\begin{tikzpicture}
	\card
	\cardstrip
	\cardbanner{banner/white.png}
	\cardicon{banner/coin.png}
	\cardprice{5}
	\cardtitle{Anbau}
	\cardcontent{Nachdem du dir eine Karte genommen und eine zusätzliche Aktion erhalten hast, musst du eine Karte aus deiner Hand entsorgen sofern du noch Handkarten hast. Du nimmst dir dafür eine Karte vom Vorrat, die \emph{genau} \coin{1} mehr kostet als die entsorgte Karte. Wenn keine solche Karte vorhanden ist, musst du zwar eine Karte entsorgen, erhältst aber keine Karte vom Vorrat. }
\end{tikzpicture}
\hspace{-1cm}
\begin{tikzpicture}
	\card
	\cardstrip
	\cardbanner{banner/white.png}
	\cardicon{banner/coin.png}
	\cardprice{5}
	\cardtitle{Handelsposten}
	\cardcontent{Wenn du nur eine Karte auf der Hand hast, musst du sie entsorgen, erhältst dafür aber kein Silber. Wenn du zwei oder mehr Karten auf der Hand hast, musst du genau zwei Karten entsorgen und nimmst dir dafür ein Silber direkt auf die Hand. Sollte kein Silber mehr im Vorrat sein, musst du die Karten trotzdem entsorgen, erhältst aber kein Silber. }
\end{tikzpicture}
\hspace{-1cm}
\begin{tikzpicture}
	\card
	\cardstrip
	\cardbanner{banner/green.png}
	\cardicon{banner/coin.png}
	\cardprice{5}
	\cardtitle{Herzog}
	\cardcontent{Diese Karte ist die einzige reine Punktekarte unter den Königreichkarten. Sie hat bis zum Ende des Spiels keine Funktion. Bei Spielende erhält der Spieler, der diese Karte in seinem Kartensatz (Nachziehstapel, Handkarten und Ablagestapel) hat, für jedes \emph{HERZOGTUM} im Kartensatz 1 Siegpunkt. Wer mehrere \emph{HERZÖGE} besitzt, erhält für jeden \emph{HERZOG} die entsprechende Anzahl Siegpunkte.} 
\end{tikzpicture}
\hspace{-1cm}
\begin{tikzpicture}
	\card
	\cardstrip
	\cardbanner{banner/white.png}
	\cardicon{banner/coin.png}
	\cardprice{5}
	\cardtitle{Kerkermeister}
	\cardcontent{Jeder Mitspieler (beginnend bei deinem linken Nachbarn) muss entweder zwei Karten ablegen oder einen \emph{FLUCH} vom Stapel nehmen. Ein Spieler kann sich entscheiden, die Karten abzulegen, auch wenn er nur eine oder gar keine Karte auf der Hand hat. Er legt dann nur so viele Karten ab, wie er kann. Er kann sich auch entscheiden, einen \emph{FLUCH} zu nehmen, wenn es keine \emph{FLÜCHE} mehr im Vorrat gibt.}
\end{tikzpicture}
\hspace{-1cm}
\begin{tikzpicture}
	\card
	\cardstrip
	\cardbanner{banner/white.png}
	\cardicon{banner/coin.png}
	\cardprice{5}
	\cardtitle{Lakai}
	\cardcontent{Du entscheidest dich für eine der beiden Optionen: Entweder erhältst du +\coin{2} in diesem Zug oder du legst alle deine Handkarten ab und ziehst vier neue Karten nach. Wenn du die zweite Option wählst, müssen außerdem alle Mitspieler, die fünf oder mehr Karten auf der Hand haben (alle anderen sind nicht betroffen), diese ablegen und ebenfalls vier Karten nachziehen. Jeder Spieler (auch wenn er von dem Angriff nicht betroffen ist) kann eine oder mehrere Reaktionskarten vorzeigen, wenn du den \emph{LAKAIEN} spielst.}
\end{tikzpicture}
\hspace{-1cm}
\begin{tikzpicture}
	\card
	\cardstrip
	\cardbanner{banner/white.png}
	\cardicon{banner/coin.png}
	\cardprice{5}
	\cardtitle{Saboteur}
	\cardcontent{Jeder Mitspieler (beginnend bei deinem linken Nachbarn) muss solange Karten von seinem Nachziehstapel aufdecken, bis er eine Karte aufdeckt, die mindestens \emph{3} kostet (\emph{Achtung:} Die Kosten einer Karte können durch die \emph{BRÜCKE} verringert werden). Er muss diese Karte sofort entsorgen und darf sich dafür eine Karte aus dem Vorrat nehmen, die mindestens \emph{2} weniger kostet. Die restlichen aufgedeckten Karten werden abgelegt. Sollte im gesamten restlichen Nachziehstapel keine Karte mit passendem Wert vorhanden sein, wird der Ablagestapel ohne die bereits aufgedeckten Karten gemischt und zum neuen Nachziehstapel. Sollte dann immer noch keine passende Karte zu fnden sein, legt der Spieler alle Karten ab und nichts passiert.}
\end{tikzpicture}
\hspace{-1cm}
\begin{tikzpicture}
	\card
	\cardstrip
	\cardbanner{banner/white.png}
	\cardicon{banner/coin.png}
	\cardprice{5}
	\cardtitle{Tribut}
	\cardcontent{Dein linker Mitspieler muss die obersten beiden Karten seines Nachziehstapels aufdecken und ablegen. Da der \emph{TRIBUT} keine Angriffskarte ist, kann sich der Mitspieler nicht gegen diese Anweisung wehren. Der Spieler, der ein \emph{TRIBUT} ausgespielt hat, erhält für die erste Karte den genannten Bonus. Für die zweite Karte erhält der Spieler nur dann ein weiteres Mal den Bonus, wenn nicht die gleiche Karte wie zuvor aufgedeckt wurde. Kombinierte Karten bringen dem Spieler auch doppelte Boni.}
\end{tikzpicture}
\hspace{-1cm}
\begin{tikzpicture}
	\card
	\cardstrip
	\cardbanner{banner/whitegreen.png}
	\cardicon{banner/coin.png}
	\cardprice{6}
	\cardtitle{Adelige}
	\cardcontent{Diese Karte ist eine kombinierte Aktions- und Punktekarte. Sie kann in der Aktionsphase eingesetzt werden und bringt außerdem bei Spielende 2 Siegpunkte. Wenn du diese Karte ausspielst, musst du dich entscheiden, ob du entweder 3 Karten ziehst oder  2 weitere Aktionen ausspielen willst. Du darfst die Anweisungen aber nicht mischen oder aufteilen. Wenn du die Karte das nächste Mal auf der Hand hast und ausspielst, darfst du natürlich eine andere Wahl treffen.}
\end{tikzpicture}
\hspace{-1cm}
\begin{tikzpicture}
	\card
	\cardstrip
	\cardbanner{banner/goldgreen.png}
	\cardicon{banner/coin.png}
	\cardprice{6}
	\cardtitle{Harem}
	\cardcontent{Diese Karte ist eine kombinierte Geld- und Punktekarte. Sie wird während des Zugs wie eine normale Geldkarte eingesetzt und bringt außerdem bei Spielende 2 Siegpunkte.}
\end{tikzpicture}
\hspace{-1cm}
\begin{tikzpicture}
	\card
	\cardstrip
	\cardbanner{banner/white.png}
	\cardtitle{\scriptsize{Empfohlene 10er Sätze\qquad}}
	\cardcontent{\emph{Siegestanz:}
		\\
		Adlige, Anbau, Brücke, Eisenhütte, Große Halle, Handlanger, Harem, Herzog, Maskerade, Späher 
		\\
		\smallskip
		\\
		\emph{Geheime Pläne:}
		\\
		Armenviertel, Eisenhütte, Handelsposten, Handlanger, Harem, Saboteur, Tribut, Trickser, Verschwörer, Verwalter
		\\
		\smallskip
		\\
		\emph{Beste Wünsche:}
		\\
		Anbau, Armenviertel, Burghof, Handelsposten, Kerkermeister, Kupferschmied, Maskerade, Späher, Verwalter, Wunschbrunnen
		\\
		\smallskip
		\\
		\emph{Demontage} (Intrige + \textit{Basisspiel}:)
		\\
		Bergwerk, Brücke, Geheimkammer, Kerkermeister, Trickser, Saboteur, \textit{Dieb}, \textit{Spion}, \textit{Thronsaal}, \textit{Umbau}
		\\
		\smallskip
		\\
		\emph{Eine Hand voll} (Intrige + \textit{Basisspiel}:)
		\\
		Adlige, Burghof, Kerkermeister, Lakai, Verwalter, \textit{Bürokrat}, \textit{Kanzler}, \textit{Miliz}, \textit{Mine}, \textit{Ratsversammlung}
		\\
		\smallskip
		\\
		\emph{Untergebene} (Intrige + \textit{Basisspiel}:)
		\\
		Adlige, Baron, Lakai, Maskerade, Verwalter, Handlanger, \textit{Bibliothek}, \textit{Hexe}, \textit{Jahrmarkt}, \textit{Keller}
		\\
	}
\end{tikzpicture}
\hspace{-1cm}

	    % Basic settings for this card set
\renewcommand{\cardcolor}{intrigue}
\renewcommand{\cardextension}{Edition II}
\renewcommand{\cardextensiontitle}{Die Intrige}

\clearpage
\newpage
\section{\cardextension \ - \cardextensiontitle}

\begin{tikzpicture}
	\card
	\cardstrip
	\cardbanner{banner/whitegreen.png}
	\cardicon{banner/coin.png}
	\cardprice{6}
	\cardtitle{Adelige}
	\cardcontent{Diese Karte ist \emph{zugleich} eine Aktions- und eine Punktekarte. Wenn du sie ausspielst, kannst du wählen, 3 Karten nachzuziehen \emph{oder} 2 zusätzliche Aktionen zu erhalten. Die beiden Anweisungen können jedoch nicht geteilt und gemischt werden. Bei Spielende sind die Adeligen 2 Punkte wert. Im Spiel zu 3. und zu 4. werden 12 Karten verwendet, im Spiel zu 2. werden 8 Karten verwendet.}
\end{tikzpicture}
\hspace{-1cm}
\begin{tikzpicture}
	\card
	\cardstrip
	\cardbanner{banner/white.png}
	\cardicon{banner/coin.png}
	\cardprice{5}
	\cardtitle{Anbau}
	\cardcontent{Du ziehst zuerst eine Karte. Danach \emph{musst} du eine Karte aus deiner Hand entsorgen und dann eine Karte nehmen, die genau 1 Geld mehr kostet als die entsorgte Karte. Ist keine solche Karte im Vorrat, erhältst du keine Karte, musst jedoch trotzdem eine entsorgen. Wenn du keine Karte zum Entsorgen hast, entsorgst du keine und nimmst dir keine Karte.}
\end{tikzpicture}
\hspace{-1cm}
\begin{tikzpicture}
	\card
	\cardstrip
	\cardbanner{banner/white.png}
	\cardicon{banner/coin.png}
	\cardprice{3}
	\cardtitle{Armenviertel}
	\cardcontent{Du erhältst 2 zusätzliche Aktionen. Dann \emph{musst} du deine Kartenhand vorzeigen. Wenn du keine Aktionskarten auf der Hand hast (auch kombinierte Aktions-/Punktekarten sind Aktionskarten), musst du 2 Karten nachziehen. Sollte die erste der gezogenen Karten eine Aktionskarte sein, ziehst du trotzdem eine zweite Karte.}
\end{tikzpicture}
\hspace{-1cm}
\begin{tikzpicture}
	\card
	\cardstrip
	\cardbanner{banner/white.png}
	\cardicon{banner/coin.png}
	\cardprice{5}
	\cardtitle{Austausch}
	\cardcontent{Entsorge zuerst eine Karte aus deiner Hand. Dann nimmst du dir eine Karte vom Vorrat, die maximal 2 mehr kostet als die entsorgte Karte. Eine Karte kostet nur dann maximal 2 mehr, wenn die restlichen Kosten (z. B. TRANK aus Empires oder SCHULDEN aus Alchemie) gleich oder niedriger sind. Wenn die genommene Karte eine Aktions- und/oder Geldkarte ist, legst du die Karte oben auf deinen Nachziehstapel. Ansonsten legst du die Karte auf den Ablagestapel. Ist die genommene Karte eine Punktekarte, nimmt sich jeder Mitspieler – beginnend bei deinem linken Mitspieler – einen Fluch. Ist die genommene Karte eine Punktekarte sowie eine Aktions- oder Geldkarte (z. B. MÜHLE), legst du die Karte oben auf deinen Nachziehstapel und jeder Mitspieler muss sich einen Fluch nehmen.}
\end{tikzpicture}
\hspace{-1cm}
\begin{tikzpicture}
	\card
	\cardstrip
	\cardbanner{banner/white.png}
	\cardicon{banner/coin.png}
	\cardprice{4}
	\cardtitle{Baron}
	\cardcontent{Du musst kein Anwesen ablegen, auch wenn du eines auf der Hand hast. Wenn du jedoch keines ablegst, musst du dir ein Anwesen nehmen, so lange noch welche im Vorrat sind. Du kannst nicht nur den +1 Kauf nutzen und die übrigen Anweisungen ignorieren.}
\end{tikzpicture}
\hspace{-1cm}
\begin{tikzpicture}
	\card
	\cardstrip
	\cardbanner{banner/white.png}
	\cardicon{banner/coin.png}
	\cardprice{4}
	\cardtitle{Bergwerk}
	\cardcontent{Du ziehst immer eine Karte nach und erhältst 2 zusätzliche Aktionen. Dann musst du entscheiden, ob du das Bergwerk entsorgst, bevor du weitere Aktionen ausspielst oder in die anderen Phasen übergehst. Wenn du das Bergwerk auf einen Thronsaal spielst, kannst du die Karte nur einmal entsorgen (d. h., du erhältst insgesamt +2 Karten, +4 Aktionen aber nur +2 Geld ).}
\end{tikzpicture}
\hspace{-1cm}
\begin{tikzpicture}
	\card
	\cardstrip
	\cardbanner{banner/white.png}
	\cardicon{banner/coin.png}
	\cardprice{4}
	\cardtitle{Brücke}
	\cardcontent{Die Kosten sind für alle Belange um 1 Geld reduziert. Wenn du z. B. ein Bergwerk ausspielst, danach eine Brücke, dann eine Eisenhütte, könntest du dir für die Eisen- hütte ein Herzogtum nehmen (kostet durch die Brücke nur noch 4 Geld). Die Karten der Spieler (Handkarten, Nachziehstapel und Ablagestapel) sind auch betroffen. Der Effekt ist kumulativ. Wenn du die Brücke auf einen Thronsaal spielst, sind die Kosten der Karten für diesen Zug um 2 Geld reduziert. Die Kosten sinken niemals unter 0 Geld. Wenn du eine Brücke und dann einen Anbau ausspielst, kannst du ein Kupfer entsorgen (das immer noch 0 Geld kostet) und dir einen Handlanger dafür nehmen (kostet durch die Brücke nur noch 1 Geld).}
\end{tikzpicture}
\hspace{-1cm}
\begin{tikzpicture}
	\card
	\cardstrip
	\cardbanner{banner/white.png}
	\cardicon{banner/coin.png}
	\cardprice{2}
	\cardtitle{Burghof}
	\cardcontent{Du ziehst 3 Karten und nimmst diese auf die Hand bevor du eine Karte auf den Nachziehstapel legst. Die Karte, die du auf den Nachziehstapel legst, muss keine der 3 gerade gezogenen Karten sein.}
\end{tikzpicture}
\hspace{-1cm}
\begin{tikzpicture}
	\card
	\cardstrip
	\cardbanner{banner/blue.png}
	\cardicon{banner/coin.png}
	\cardprice{4}
	\cardtitle{Diplomatin}
	\cardcontent{Diese Karte ist eine Aktions- und Reaktionskarte. Wird sie als Aktion in der Aktionsphase ausgespielt, nimmst du 2 Karten. Hast du dann 5 oder weniger Karten auf der Hand, erhältst du außerdem + 2 Aktionen. Spielt ein Mitspieler eine Angriffskarte aus und du hast zu diesem Zeitpunkt 5 oder mehr Karten auf der Hand, darfst du diese Karte – bevor der ausgespielte Angriff ausgeführt wird – aus der Hand aufdecken. Wenn du das tust, nimmst du diese DIPLOMATIN wieder auf die Hand, ziehst 2 Karten und legst dann 3 Karten (auch möglich inklusive dieser DIPLOMATIN) ab. Hast du dann immer noch 5 oder mehr Karten sowie eine DIPLOMATIN auf der Hand, darfst du die DIPLOMATIN noch einmal aufdecken – und dies so oft wiederholen wie du möchtest und die Bedingung der 5 oder mehr Karten auf der Hand erfüllt ist. Erst dann wird der Angriff ausgeführt. Hast du mehrere Reaktionskarten auf der Hand, mit denen du auf das Ausspielen einer Angriffskarte reagieren kannst, darfst du diese nacheinander in beliebiger Reihenfolge aufdecken.}
\end{tikzpicture}
\hspace{-1cm}
\begin{tikzpicture}
	\card
	\cardstrip
	\cardbanner{banner/white.png}
	\cardicon{banner/coin.png}
	\cardprice{4}
	\cardtitle{Eisenhütte}
	\cardcontent{Du nimmst dir eine Karte vom Vorrat und legst sie auf deinen Ablagestapel. Je nach Kartentyp der genommenen Karte erhältst du einen Bonus. Nimmst du eine Karte mit kombiniertem Kartentyp, z. B. Große Halle erhältst du +1 Aktion (weil die Große Halle eine Aktionskarte ist) und +1 Karte (weil die Große Halle auch eine Punktekarte ist).}
\end{tikzpicture}
\hspace{-1cm}
\begin{tikzpicture}
	\card
	\cardstrip
	\cardbanner{banner/white.png}
	\cardicon{banner/coin.png}
	\cardprice{4}
	\cardtitle{Geheimgang}
	\cardcontent{Du ziehst 2 Karten und erhältst + 1 Aktion. Dann nimmst du eine beliebige Karte aus deiner Hand (auch ggf. eine, die du gerade gezogen hast) und legst sie an eine beliebige Stelle in deinen Nachziehstapel. Du darfst sie oben drauf, unten drunter oder irgendwo in die Mitte legen. Du darfst dabei die Karten deines Nachziehstapels zählen, aber nicht ansehen. Befinden sich keine Karten in deinem Nachziehstapel, wird die zurückgelegte Karte zur einzigen Karte in deinem Nachziehstapel.}
\end{tikzpicture}
\hspace{-1cm}
\begin{tikzpicture}
	\card
	\cardstrip
	\cardbanner{banner/white.png}
	\cardicon{banner/coin.png}
	\cardprice{5}
	\cardtitle{Handelsposten}
	\cardcontent{Wenn du 2 oder mehr Karten auf der Hand hast, \emph{musst} du 2 Karten entsorgen und dir dafür ein Silber nehmen. Du nimmst das Silber direkt auf die Hand und kannst es auch in der Kaufphase verwenden. Wenn kein Silber mehr im Vorrat ist, erhältst du kein Silber, musst jedoch trotzdem 2 Karten entsorgen. Wenn du nur 1 Karte auf der Hand hast, musst du diese entsorgen, erhältst jedoch kein Silber. Wenn du keine Karte mehr auf der Hand hast, kannst du nichts entsorgen und erhältst auch kein Silber.}
\end{tikzpicture}
\hspace{-1cm}
\begin{tikzpicture}
	\card
	\cardstrip
	\cardbanner{banner/white.png}
	\cardicon{banner/coin.png}
	\cardprice{2}
	\cardtitle{Handlanger}
	\cardcontent{Wähle 2 verschiedene Anweisungen. Du darfst nicht eine Anweisung zweimal wählen. Du musst zuerst beide Anweisungen auswählen und sie dann erst (in jeder möglichen Reihenfolge) ausführen. Du kannst nicht eine Karte nachziehen und dann erst die zweite Anweisung wählen.}
\end{tikzpicture}
\hspace{-1cm}
\begin{tikzpicture}
	\card
	\cardstrip
	\cardbanner{banner/goldgreen.png}
	\cardicon{banner/coin.png}
	\cardprice{6}
	\cardtitle{Harem}
	\cardcontent{Diese Karte ist zugleich eine Geld- und eine Punktekarte. Du kannst sie in der Kaufphase spielen, genau wie ein Silber. Bei Spielende ist der Harem 2 Punkte wert. Im Spiel zu 3. und zu 4. werden 12 Karten verwendet, im Spiel zu 2. werden 8 Karten verwendet.}
\end{tikzpicture}
\hspace{-1cm}
\begin{tikzpicture}
	\card
	\cardstrip
	\cardbanner{banner/white.png}
	\cardicon{banner/coin.png}
	\cardprice{2}
	\cardtitle{Herumtreiberin}
	\cardcontent{Die Karte, die du entsorgst oder vom Müllstapel nimmst, muss den Typ AKTION beinhalten, d. h. sie kann auch eine kombinierte Aktionskarte (z. B. MÜHLE) sein. Genommene Karten werden auf den Ablagestapel gelegt – es sei denn, auf der Karte steht etwas anderes. Wird eine Karte entsorgt, die einen speziellen Effekt beim Entsorgen hat, tritt dieser ein.}
\end{tikzpicture}
\hspace{-1cm}
\begin{tikzpicture}
	\card
	\cardstrip
	\cardbanner{banner/green.png}
	\cardicon{banner/coin.png}
	\cardprice{5}
	\cardtitle{Herzog}
	\cardcontent{Diese Karte hat bis zum Ende des Spiels keine Funktion. Bei Spielende ist der Herzog 1 Punkt pro Herzogtum in Handkarten, Nachziehstapel und Ablagestapel wert. Im Spiel zu 3. und zu 4. werden 12 Karten verwendet, im Spiel zu 2. werden 8 Karten verwendet.}
\end{tikzpicture}
\hspace{-1cm}
\begin{tikzpicture}
	\card
	\cardstrip
	\cardbanner{banner/white.png}
	\cardicon{banner/coin.png}
	\cardprice{5}
	\cardtitle{Höflinge}
	\cardcontent{Decke eine Karte aus deiner Hand auf. Zähle dann die Typen, denen diese Karte angehört – also AKTION, GELD, REAKTION, ANGRIFF, PUNKTE, FLUCH etc. Pro Typ, dem die Karte angehört, entscheidest du dich für eine der vier angegebenen Optionen. Dabei darfst du keine der Optionen doppelt auswählen. Wenn du zum Beispiel eine PATROUILLE (AKTION) aufdeckst, darfst du eine Option auswählen, deckst du einen KARAWANENWÄCHTER aus Abenteuer (AKTION – DAUER – REAKTION) auf, darfst du 3 unterschiedliche Optionen wählen. Entscheidest du dich für das Gold, legst du dieses auf den Ablagestapel. Kannst du keine Handkarte aufdecken, erhältst du nichts.}
\end{tikzpicture}
\hspace{-1cm}
\begin{tikzpicture}
	\card
	\cardstrip
	\cardbanner{banner/white.png}
	\cardicon{banner/coin.png}
	\cardprice{5}
	\cardtitle{Kerkermeister}
	\cardcontent{Jeder Mitspieler, beginnend mit dem Spieler links vom Angreifer, muss sich eine der beiden Anweisungen wählen und diese dann ausführen. Ein Spieler kann wählen, 2 Karten abzulegen, auch wenn er weniger als 2 Karten auf der Hand hat. Hat er nur eine Karte auf der Hand legt er diese ab. Hat er keine Karte mehr auf der Hand, muss er auch keine ablegen. Ein Spieler kann wählen einen Fluch zu nehmen, auch wenn keine Fluchkarten mehr im Vorrat sind. In diesem Fall nimmt er keinen Fluch. Fluchkarten nehmen die Spieler direkt auf die Hand.}
\end{tikzpicture}
\hspace{-1cm}
\begin{tikzpicture}
	\card
	\cardstrip
	\cardbanner{banner/white.png}
	\cardicon{banner/coin.png}
	\cardprice{5}
	\cardtitle{Lakai}
	\cardcontent{Zunächst entscheidest du dich für eine der der beiden Anweisungen. Entweder erhältst du +2 virtuelles Geld \emph{oder} du wählst die zweite Anweisung für den Angriff. In diesem Fall sind nur Spieler mit 5 oder mehr Karten auf der Hand betroffen. Wehrt ein Spieler den Angriff mit einem Burggraben ab, darf er weder Karten nachziehen, noch muss er Karten ablegen. Ein Spieler kann auf den Angriff mit der Geheimkammer reagieren, auch wenn er weniger als 5 Karten auf der Hand hat. Anschließend hast du +1 Aktion, unabhängig davon, welche der beiden Anweisungen du gewählt hast.}
\end{tikzpicture}
\hspace{-1cm}
\begin{tikzpicture}
	\card
	\cardstrip
	\cardbanner{banner/white.png}
	\cardicon{banner/coin.png}
	\cardprice{3}
	\cardtitle{Maskerade}
	\cardcontent{Du ziehst zuerst 2 Karten. Dann wählen alle Spieler gleichzeitig eine Karte aus ihrer Hand und legen diese verdeckt zwischen sich und den Spieler zu ihrer Linken. Erst dann nehmen alle Spieler die Karten, die sie vom Spieler rechts bekommen haben auf. Die Spieler wählen also zuerst, welche Karte sie weiter geben, bevor sie sehen, welche Karte sie bekommen. Am Ende darfst nur du eine Karte aus deiner Hand entsorgen. Die Maskerade ist kein Angriff. Die übrigen Spieler dürfen also keine Reaktionskarten aus ihrer Hand vorzeigen um sich zu schützen.}
\end{tikzpicture}
\hspace{-1cm}
\begin{tikzpicture}
	\card
	\cardstrip
	\cardbanner{banner/whitegreen.png}
	\cardicon{banner/coin.png}
	\cardprice{4}
	\cardtitle{Mühle}
	\cardcontent{Diese Karte ist eine kombinierte Aktions- und Punktekarte. Als Punktekarte bringt sie beim Zählen der Punkte 1. Spielst du die MÜHLE als Aktionskarte aus, ziehst du 1 Karte und erhältst + 1 Aktion. Du darfst 2 Karten aus deiner Hand ablegen. Wenn du das tust, erhältst du + 2 . Tust du das nicht (weil du zum Beispiel nicht genügend Karten auf der Hand hast), erhältst du nichts. Nur, wenn du nicht mehr als eine Karte auf der Hand hast, darfst du genau eine Karte ablegen, erhältst dafür aber kein GELD.}
\end{tikzpicture}
\hspace{-1cm}
\begin{tikzpicture}
	\card
	\cardstrip
	\cardbanner{banner/white.png}
	\cardicon{banner/coin.png}
	\cardprice{5}
	\cardtitle{Patrouille}
	\cardcontent{Ziehe zuerst 3 Karten. Decke dann die obersten 4 Karten deines Nachziehstapels auf. So aufgedeckte Punktekarten (auch ggf. kombinierte) und Flüche nimmst du alle auf die Hand. Die restlichen Karten legst du in beliebiger Reihenfolge zurück auf den Nachziehstapel.}
\end{tikzpicture}
\hspace{-1cm}
\begin{tikzpicture}
	\card
	\cardstrip
	\cardbanner{banner/white.png}
	\cardicon{banner/coin.png}
	\cardprice{3}
	\cardtitle{Trickser}
	\cardcontent{Jeder Mitspieler, beginnend mit dem Spieler links vom Angreifer, muss die oberste Karte von seinem Nachziehstapel aufdecken. Er muss diese Karte entsorgen und du wählst eine Karte aus dem Vorrat, die das gleiche kostet. Diese Karte nimmt sich der Spieler und legt sie bei sich ab. Ist im Vorrat keine Karte, mit den gleichen Kosten, erhält der Spieler nichts, muss jedoch die Karte trotzdem entsorgen. Entsorgt er z. B. ein Kupfer, kannst du einen Fluch auswählen, den er nehmen muss. Du kannst auch die selbe Karte wählen, die er entsorgt hat. Die gewählte Karte muss im Vorrat zur Verfügung stehen. Du kannst also keine Karte aus einem leeren Stapel oder vom Müll wählen. Sind keine Karten mehr im Vorrat, die genau so viel kosten, wie die entsorgte Karte, erhält der Spieler nichts. Deckt ein Spieler den Burggraben aus seiner Hand auf, muss er keine Karte vom Nachziehstapel aufdecken und entsorgen und erhält auch keine Karte.}
\end{tikzpicture}
\hspace{-1cm}
\begin{tikzpicture}
	\card
	\cardstrip
	\cardbanner{banner/white.png}
	\cardicon{banner/coin.png}
	\cardprice{4}
	\cardtitle{Verschwörer}
	\cardcontent{Du überprüfst die Bedingung ob du +1 Karte und +1 Aktion erhältst wenn du den Verschwörer ausgespielt hast. Wenn die Bedingung später im Zug erfüllt wird, überprüfst du die Bedingung nicht rückwirkend. Wird eine Karte auf den Thronsaal gespielt, zählt der Thronsaal selbst als gespielte Aktionskarte und die darauf gespielte Aktionskarte zusätzlich zweimal als gespielte Aktionskarte. Wenn du z. B. den Verschwörer auf den Thronsaal spielst, ist der Thronsaal die erste Aktionskarte, der zuerst ausgespielte Verschwörer ist die zweite Aktionskarte (du erhältst also keine +1 Karte und keine +1 Aktion). Wenn du den Verschwörer zum zweiten mal ausspielst, hast du 3 Aktionskarten ausgespielt und erhältst +1 Karte und +1 Aktion.}
\end{tikzpicture}
\hspace{-1cm}
\begin{tikzpicture}
	\card
	\cardstrip
	\cardbanner{banner/white.png}
	\cardicon{banner/coin.png}
	\cardprice{3}
	\cardtitle{Verwalter}
	\cardcontent{Wenn du dich entscheidest, 2 Karten zu entsorgen und 2 oder mehr Karten auf der Hand hast, musst du genau 2 Karten entsorgen. Wenn du dich entscheidest, 2 Karten zu entsorgen, aber aber nur 1 Karte auf der Hand hast musst du diese Karte entsorgen. Du kannst die verschiedenen Anweisungen nicht mischen, du musst wählen: \emph{entweder} +2 Karten \emph{oder} +2 Geld \emph{oder} 2 Karten entsorgen.}
\end{tikzpicture}
\hspace{-1cm}
\begin{tikzpicture}
	\card
	\cardstrip
	\cardbanner{banner/white.png}
	\cardicon{banner/coin.png}
	\cardprice{3}
	\cardtitle{Wunschbrunnen}
	\cardcontent{Du ziehst zuerst eine Karte nach. Dann benennst du eine Karte (z. B. „Kupfer“, nicht „Geld“) und deckst die oberste Karte von deinem Nachziehstapel auf. Wenn es sich um die benannte Karte handelt, nimmst du sie auf die Hand. Wenn nicht, legst du sie zurück auf den Nachziehstapel.}
\end{tikzpicture}
\hspace{-1cm}
\begin{tikzpicture}
	\card
	\cardstrip
	\cardbanner{banner/white.png}
	\cardtitle{\scriptsize{Empfohlene 10er Sätze\qquad}}
	\cardcontent{\emph{Siegestanz:}
		\\
		Adlige, Austausch, Baron, Eisenhütte, Harem, Herzog, Höflinge, Maskerade, Mühle, Patrouille
		\\
		\smallskip
		\\
		\emph{Verschwörung:}
		\\
		Bergwerk, Eisenhütte, Geheimgang, Handelsposten, Handlanger, Herumtreiberin, Kerkermeister, Trickser, Verschwörer, Verwalter 
		\\
		\smallskip
		\\
		\emph{Beste Wünsche:}
		\\
		Anbau, Armenviertel, Baron, Burghof, Diplomatin, Geheimgang, Herzog, Kerkermeister, Verschwörer, Wunschbrunnen
		\\
	}
\end{tikzpicture}
\hspace{-1cm}
	    % Basic settings for this card set
\renewcommand{\cardcolor}{seaside}
\renewcommand{\cardextension}{Erweiterung I}
\renewcommand{\cardextensiontitle}{Seaside}

\clearpage
\newpage
\section{\cardextension \ - \cardextensiontitle}

\begin{tikzpicture}
	\card
	\cardstrip
	\cardbanner{banner/white.png}
	\cardicon{banner/coin.png}
	\cardprice{}
	\cardtitle{}
	\cardcontent{}
\end{tikzpicture}

\hspace{-1cm}
\begin{tikzpicture}
\card
\cardstrip
\cardbanner{banner/white.png}
\cardicon{}
\cardprice{}
\cardtitle{\scriptsize{Empfohlene 10er Sätze\qquad}}
\cardcontent{\emph{Name:}
	\\
	Karten ...
	\\
	\smallskip
	\\
	\emph{Name:}
	\\
	Karten ...
	\\
	\smallskip
	\\
	\emph{Name:}
	\\
	Karten ...
	\\
	\smallskip
	\\
	\emph{Name:}
	\\
	Karten ...
	\\
	\smallskip
	\\
	\emph{Name:}
	\\
	Karten ...
	\\
	\smallskip
	\\
	\emph{Name:}
	\\
	Karten ...
	\\
}
\end{tikzpicture}
\hspace{1cm}
	    % Basic settings for this card set
\renewcommand{\cardcolor}{alchemy}
\renewcommand{\cardextension}{Erweiterung II}
\renewcommand{\cardextensiontitle}{Die Alchemisten}

\clearpage
\newpage
\section{\cardextension \ - \cardextensiontitle}

\begin{tikzpicture}
	\card
	\cardstrip
	\cardbanner{banner/green.png}
	\cardicon{banner/potion.png}
	\cardtitle{Weinberg}
	\cardcontent{Diese Königreichkarte ist eine Punktekarte, keine Aktionskarte. Sie hat bis zum Ende des Spiels keine Funktion. Bei der Wertung zählt sie 1 Punkt pro volle 3 Aktionskarten im gesamten Kartensatz (Nachziehstapel, Ablagestapel und Handkarten) des Spielers. Du zählst alle deine Aktionskarten bei Spielende, teilst die Anzahl durch 3 und rundest ab. Kombinierte Aktionskarten sind auch Aktionskarten. Für 11 Aktionskarten erhältst du beispielsweise für jeden deiner Weinberge 3 Punkte. Im Spiel zu 3. und 4. werden 12 Karten verwendet, im Spiel zu 2. werden 8 Karten verwendet.}
\end{tikzpicture}
\hspace{-1cm}
\begin{tikzpicture}
	\card
	\cardstrip
	\cardbanner{banner/white.png}
	\cardicon{banner/potion.png}
	\cardtitle{Verwandlung}
	\cardcontent{Hast du keine Karte mehr auf der Hand, die du entsorgen könntest, erhältst du nichts. Wenn du einen Fluch entsorgst erhältst du nichts. Der Fluch ist keine Punkte-, keine Aktions- und keine Geldkarte. Entsorgst du eine Karte mit kombiniertem Kartentyp, erhältst du den Bonus für beide Kartentypen. Für die Adeligen (Dominion – Die Intrige) nimmst du dir z. B. ein Herzogtum und ein Gold. Die Karten nimmst du dir vom Vorrat. Ist keine entsprechende Karte mehr im Vorrat, erhältst du nichts.
	}
\end{tikzpicture}
\hspace{-1cm}
\begin{tikzpicture}
	\card
	\cardstrip
	\cardbanner{banner/white.png}
	\cardicon{banner/coin.png}
	\cardprice{2}
	\cardtitle{\scriptsize{Kräuterkundiger}}
	\cardcontent{Wenn du diese Karte ausspielst erhältst du +1 Geld für die Kaufphase und du darfst eine weitere Karte kaufen. Wenn du den Kräuterkundigen ablegst (normalerweise in der Aufräumphase), darfst du eine Geldkarte, die vor dir ausliegt, zurück auf deinen Nachziehstapel legen, statt sie abzulegen. Wenn dein Nachziehstapel leer ist, legst du nur die Geldkarte als neuen Nachziehstapel bereit. Du entscheidest, in welcher Reihenfolge du die vor dir ausliegenden Karten ablegst. Hast du z. B. einen Kräuterkundigen, einen Alchemisten und einen Trank ausliegen, darfst du zuerst den Alchemisten zurück auf deinen Nachziehstapel legen, dann den Kräuterkundigen ablegen und dafür den Trank zurück auf deinen Nachziehstapel legen. Wenn du mehrere Kräuterkundige im Spiel hast, darfst du für jeden davon eine Geldkarte zurück auf deinen Nachziehstapel legen.
	}
\end{tikzpicture}
\hspace{-1cm}
\begin{tikzpicture}
	\card
	\cardstrip
	\cardbanner{banner/white.png}
	\cardicon{banner/coin.png}
	\cardprice{2}
	\cardiconaddition{banner/potion.png}
	\cardtitle{\quad Apotheker}
	\cardcontent{Wenn du diese Karte ausspielst musst du sofort eine Karte nachziehen. Dann deckst du die obersten 4 Karten von deinem Nachziehstapel auf. Sollte dein Nachziehstapel beim Aufdecken zu Ende gehen, mischt du deinen Ablagestapel. Hast du nach dem Mischen immernoch weniger als 4 Karten zum Aufdecken, deckst du nur soviele Karten auf wie möglich. Danach musst du alle aufgedeckten Kupfer- und Trankkarten auf die Hand nehmen. Die übrigen aufgedeckten Karten legst du in beliebiger Reihenfolge zurück auf deinen Nachziehstapel. Wenn dein Nachziehstapel nach dem aufdecken der Karten leer ist, werden die zurück gelegten Karten dein neuer Nachziehstapel.
	}
\end{tikzpicture}
\hspace{-1cm}
\begin{tikzpicture}
	\card
	\cardstrip
	\cardbanner{banner/white.png}
	\cardicon{banner/coin.png}
	\cardprice{2}
	\cardiconaddition{banner/potion.png}
	\cardtitle{\quad Universität}
	\cardcontent{Du darfst die eine Aktionskarte, die bis zu 5 Geld kostet, vom Vorrat nehmen. Du musst jedoch keine Karte nehmen. Kombinierte Aktionskarten sind auch Aktionskarten. Du darfst keine Karte mit Trank in den Kosten nehmen.
	}
\end{tikzpicture}
\hspace{-1cm}
\begin{tikzpicture}
	\card
	\cardstrip
	\cardbanner{banner/white.png}
	\cardicon{banner/coin.png}
	\cardprice{2}
	\cardiconaddition{banner/potion.png}
	\cardtitle{\quad Vision}
	\cardcontent{Zuerst decken alle Spieler die oberste Karte ihres Nachziehstapels auf. Du entscheidest dann Spieler für Spieler extra (auch bei dir selbst), ob er die aufgedeckte Karte auf seinen Ablagestapel oder zurück auf seinen Nachziehstapel legt. Danach deckst du solange Karten von deinem Nachziehstapel auf, bis du eine Karte aufgedeckt hast, die keine Aktionskarte ist. Nun nimmst du alle gerade aufgedeckten Karten auf die Hand (auch die Nicht-Aktionskarte). Ist bereits die erste Karte, die du dabei aufgedeckt hast, keine Aktionskarte, so nimmst du nur diese eine Karte auf die Hand. Wenn du nach dem Mischen deines Ablagestapels alle Karten aufgedeckt hast und nur Aktionskarten o en liegen, so nimmst du alle aufgedeckten Karten auf die Hand. Kombinierte Aktionskarten sind auch Aktionskarten. Die durch die erste Anweisung aufgedeckten Karten aller Spieler werden hierbei nicht mehr beachtet. Du kannst auch keine dieser Karten auf die Hand nehmen.
	}
\end{tikzpicture}
\hspace{-1cm}
\begin{tikzpicture}
	\card
	\cardstrip
	\cardbanner{banner/white.png}
	\cardicon{banner/coin.png}
	\cardprice{3}
	\cardiconaddition{banner/potion.png}
	\cardtitle{\quad Alchemist}
	\cardcontent{Wenn du diese Karte ausspielst musst du sofort 2 Karten nachziehen und darfst dann eine weitere Aktionskarte ausspielen. Wenn mindestens ein Trank im Spiel ist, darfst du in der Aufräumphase alle vor dir ausliegenden Alchemisten zurück auf deinen Nachziehstapel legen, statt sie abzulegen. Hast du mehrere Alchemisten im Spiel, darfst du für jeden einzeln entscheiden, ob du ihn normal ablegst oder zurück auf deinen Nachziehstapel legst. Ist dein Nachziehstapel leer, legst du den Alchemisten alleine als Nachziehstapel bereit. Du legst den oder die Alchemisten ab oder zurück auf deinen Nachziehstapel, bevor du die Karten am Ende deines Zuges nachziehst. Du darfst in der Kaufphase einen Trank ausspielen, auch wenn du keine Karte damit kaufst.
	}
\end{tikzpicture}
\hspace{-1cm}
\begin{tikzpicture}
	\card
	\cardstrip
	\cardbanner{banner/gold.png}
	\cardicon{banner/coin.png}
	\cardprice{3}
	\cardiconaddition{banner/potion.png}
	\cardtitle{\quad \scriptsize{Stein\thinspace der\thinspace Weisen}}
	\cardcontent{Diese Karte ist eine Geldkarte und eine Königreichkarte. Sie ist nur im Spiel, wenn sie als eine der 10 ausliegenden Königreichkarten für dieses Spiel ausgewühlt wurde. Der Stein der Weisen darf, wie andere Geldkarten, in der Kaufphase ausgespielt werden. Wenn du den Stein der Weisen ausspielst zählst du zunächst die momentane Anzahl der Karten in deinem Nachziehstapel und in deinem Ablagestapel. Dann zählst du die Anzahl beider Stapel zusammen, teilst die Summe durch 5 und rundest ab. Das Ergebnis gibt den Wert der Karte für diese Kaufphase an. Der Wert gilt für diese gesamte Kaufphase, auch wenn sich die Anzahl der Karten noch verändert. Spielt ein Spieler mehrere dieser Karten, so hat jede den ermittelten Wert. Wird die Karte in einem späteren Zug erneut ausgespielt, wird auch der Wert neu ermittelt. Du darfst beim Durchzählen deines Nachziehstapels weder die Karten ansehen, noch deren Reihenfolge verändern. Beim Durchzählen deines Ablagestapels darfst du die Karten ansehen und deren Reihenfolge verändern. Du zählst nur die Karten der beiden Stapel, ausgespielte, beiseite gelegte und Handkarten zählst du nicht mit. Du darfst keine Geldkarten mehr ausspielen, nachdem du eine Karte gekauft hast. Du darfst also weder den Stein der Weisen noch andere Geldkarten ausspielen, wenn du in dieser Kaufphase bereits eine Karte gekauft hast.
	}
\end{tikzpicture}
\hspace{-1cm}
\begin{tikzpicture}
	\card
	\cardstrip
	\cardbanner{banner/white.png}
	\cardicon{banner/coin.png}
	\cardprice{3}
	\cardiconaddition{banner/potion.png}
	\cardtitle{\quad Vertrauter}
	\cardcontent{Sind nicht mehr genügend Fluchkarten im Vorrat, wenn du den Vertrauten aus- spielst, werden die restlichen Fluchkarten, beginnend beim Spieler links von dir, in Spielerreihenfolge verteilt. Du ziehst immer eine Karte nach und darfst eine weitere Aktionskarte ausspielen, auch wenn keine Fluchkarten mehr im allgemeinen Vorrat sind. Die Fluchkarten legen die Spieler sofort auf ihren Ablagestapel.
	}
\end{tikzpicture}
\hspace{-1cm}
\begin{tikzpicture}
	\card
	\cardstrip
	\cardbanner{banner/white.png}
	\cardicon{banner/coin.png}
	\cardprice{4}
	\cardiconaddition{banner/potion.png}
	\cardtitle{\quad Golem}
	\cardcontent{Du deckst solange Karten von deinem Nachziehstapel auf, bis 2 Aktionskarten offen liegen, die keine Golemkarten sind. Dann legst du alle aufgedeckten Golemkarten und alle Karten, die keine Aktionskarten sind, ab. Wenn du auch nach dem Mischen deines Ablagestapels keine oder nur 1 Aktionskarte (außer dem Golem) aufdecken kannst, so führst du die Anweisung mit weniger als 2 Karten aus. Hast du auf diese Weise 1 oder 2 Aktionskarten aufgedeckt, musst du diese in beliebiger Reihenfolge ausspielen. Du darfst auf diese Weise aufgedeckte Aktionskarten nicht auf die Hand nehmen. Alle Anweisungen, die sich auf deine Handkarten beziehen, haben keinen Effekt auf die beiden aufgedeckten Karten. Ist z. B. eine der aufgedeckten Karten ein Thronsaal, so kannst du die andere der beiden Karten nicht für diesen Thronsaal auswühlen.
	}
\end{tikzpicture}
\hspace{-1cm}
\begin{tikzpicture}
	\card
	\cardstrip
	\cardbanner{banner/white.png}
	\cardicon{banner/coin.png}
	\cardprice{5}
	\cardtitle{Lehrling}
	\cardcontent{Wenn du eine Karte auf deiner Hand hast, musst du eine Karte entsorgen. Wenn du eine Karte, die 0 Geld kostet (z. B. Kupfer oder Fluch) entsorgst oder wenn du keine Karte mehr auf der Hand hast, ziehst du keine Karte nach. Ansonsten ziehst du für eine Karte, die x Geld kostet, X Karten nach. Zusätzlich dazu ziehst du 2 Karten nach, wenn die entsorgte Karte kostet. Wenn du z. B. die Karte Golem (4 Geld, Trank) entsorgst, ziehst du 6 Karten nach.
	}
\end{tikzpicture}
\hspace{-1cm}
\begin{tikzpicture}
	\card
	\cardstrip
	\cardbanner{banner/white.png}
	\cardicon{banner/coin.png}
	\cardprice{6}
	\cardiconaddition{banner/potion.png}
	\cardtitle{\quad \footnotesize{Besessenheit}}
	\cardcontent{\tiny{(Teil 1:)
	\\
	\medskip
	\\
	Der Spieler links von dir ist der aktive Spieler bei dem Extra-Zug durch die Besessenheit. Kartenanweisungen, die sich auf den aktiven Spieler beziehen, betreffen also den Spieler links von dir und seinen Kartensatz. In den Kartenanweisungen wird der aktive Spieler meist mit \enquote{du} angesprochen. Du darfst alle Karten sehen, die der Spieler links von dir in seinem Extra-Zug sieht. Das betrifft auch die Karten, die er in der Aufräumphase für seinen nächsten Zug nachzieht. Du darfst auch alle Karten ansehen, die er ansehen darf, z. B. zur Seite gelegte Karten auf dem Eingeborenendorf (Dominion – Seaside). Und du darfst alle Kartenstapel durchzählen, die er durchzählen darf. Du füllst alle Entscheidungen für den Spieler links von dir, welche Karten er ausspielt und in welcher Reihenfolge, alle Entscheidungen, die durch Kartenanweisungen erlaubt werden und welche Karten er kauft. Alle Karten, die der Spieler links von dir in seinem Extra-Zug nimmt oder kauft, legst du auf deinen Ablagestapel. Dies gilt auch, wenn er die Karte aufgrund einer bestimmten Anweisung auf die Hand oder nehmen oder anderswo ablegen müsste. Du erhältst nur Karten, die er nimmt oder kauft, keine Marker, z. B. Piratenschiff (Dominion – Seaside). Wenn der Spieler links von dir in seinem Extra-Zug Karten entsorgt, werden diese zunächst zur Seite gelegt und am Ende des Zuges (nach der Aufräumphase) auf seinen Ablagestapel gelegt. Für die Kartenanweisung, die das Entsorgen der Karte fordert, gilt diese als entsorgt, z. B. könntest du das Bergwerk (Dominion – Die Intrige) entsorgen und + erhalten. Das Bergwerk wird dabei nicht auf den Müllstapel gelegt, der Spieler verliert die Karte also nicht. Karten anderer Spieler, die während dieses Zuges entsorgt werden (z. B. durch Angriffskarten, wie Trickser oder Saboteur, Dominion – Die Intrige), werden dauerhaft entsorgt. Karten, die weitergegeben werden (z. B. durch die Maskerade, Dominion – Die Intrige), erhält der Spieler am Ende des Zuges nicht zurück.\\}}
\end{tikzpicture}
\hspace{-1cm}
\begin{tikzpicture}
	\card
	\cardstrip
	\cardbanner{banner/white.png}
	\cardicon{banner/coin.png}
	\cardprice{6}
	\cardiconaddition{banner/potion.png}
	\cardtitle{\quad \footnotesize{Besessenheit}}
	\cardcontent{\tiny{(Teil 2:)
	\\
	\medskip
	\\
	Karten, die in den Vorrat zurückgelegt werden (z. B. Botschafter, Dominion – Seaside), erhält der Spieler auch nicht zurück. Spielt der Spieler links von dir (auf deinen Wunsch hin) eine Angriffskarte, so bist du, wie alle übrigen Spieler, normal betroffen. Du kannst, wie üblich, mit Reaktionskarten aus deiner eigenen Hand auf den Angriff reagieren. Du kannst keine Reaktionskarten des Spielers links von dir verwenden um auf den Angriff zu reagieren. Besessenheit bewirkt einen Extra-Zug, wie z. B. der Außenposten (Dominion – Seaside). Dieser Extra-Zug  findet erst nach deinem Zug statt. Du hast also alle Karten abgelegt und die Handkarten für deinen nächsten Zug nachgezogen. Der Außenposten verhindert nur einen weiteren Extra-Zug durch einen weiteren Außenposten. Du kannst weitere Extrazüge durch andere Karten, wie z. B. die Besessenheit erhalten. Wenn du in deinem Zug den Außenposten (Dominion – Seaside) und die Besessenheit spielst, führst du zuerst den Extra-Zug für den Außenposten aus, danach den Extra-Zug für die Besessenheit. Wenn der Spieler links von dir (auf deinen Wunsch hin) einen Außenposten spielt, erhält er dadurch einen Extra-Zug. In diesem Extra-Zug füllt er selbst seine Entscheidungen und erhält auch Karten, die er nimmt oder kauft selbst. Wenn der Spieler links von dir (auf deinen Wunsch hin) eine weitere Besessenheit ausspielt, so wird ein weiterer Extra-Zug gespielt, in dem der Spieler links von dir (nicht du) die Entscheidungen für den Spieler links von ihm füllt. Extra-Züge (z. B. durch Besessenheit oder Außenposten) werden für die Siegbedingung nicht beachtet. Im Gegensatz zum Außenposten ist die Besessenheit keine Dauer-Karte und wird in der Aufräumphase abgelegt. Die Wirkung der Karte Besessenheit ist kumulativ, spielst du z. B. die Karte in deinem Zug zweimal, so werden auch 2 Extra-Züge durchgeführt. Wichtig: Der Extrazug durch die Besessenheit ist nicht dein Extrazug, sondern der Extrazug des Spielers links von dir. Nach dem Extrazug (oder den Extrazügen) führt der Spieler seinen normalen Zug aus.\\}}
\end{tikzpicture}
\hspace{-1cm}
\begin{tikzpicture}
	\card
	\cardbanner{banner/gold.png}
	\cardicon{banner/coin.png}
	\cardprice{4}
	\cardtitle{Trank}
\end{tikzpicture}
\hspace{-1cm}
\begin{tikzpicture}
	\card
	\cardstrip
	\cardbanner{banner/white.png}
	\cardicon{}
	\cardprice{}
	\cardtitle{\scriptsize{Empfohlene 10er Sätze\qquad}}
	\cardcontent{\emph{Verbotene Künste} (Alchemisten + \textit{Basisspiel}:)
	\\
	Besessenheit, Lehrling, Universität, Vertrauter, \textit{Dieb}, \textit{Gärten}, \textit{Keller}, \textit{Laboratorium}, \textit{Ratsversammlung}, \textit{Thronsaal}
	\\
	\smallskip
	\\
	\emph{Quacksalber:} (Alchemisten + \textit{Basisspiel}:)
	\\
	Alchemist, Apotheker, Golem, Kräuterkundiger, Verwandlung, \textit{Jahrmarkt}, \textit{Kanzler}, \textit{Keller}, \textit{Miliz}, \textit{Schmiede}
	\\
	\smallskip
	\\
	\emph{Chemiestunde:} (Alchemisten + \textit{Basisspiel}:)
	\\
	Alchemist, Golem, Stein der Weisen, Universität, \textit{Burggraben}, \textit{Bürokrat}, \textit{Hexe}, \textit{Holzfäller}, \textit{Markt}, \textit{Umbau}
	\\
	\smallskip
	\\
	\emph{Diener:} (Alchemisten + \textit{Die Intrige}:)
	\\
	Besessenheit, Golem, Verwandlung, Vision, Weinberg, \textit{Große Halle}, \textit{Handlanger}, \textit{Lakai}, \textit{Verschwörer}, \textit{Verwalter}
	\\
	\smallskip
	\\
	\emph{Geheime Forschungen:} (Alchemisten + \textit{Die Intrige}:)
	\\
	Kräuterkundiger, Stein der Weisen, Universität, Vertrauter, \textit{Adlige}, \textit{Armenviertel}, \textit{Brücke}, \textit{Kerkermeister}, \textit{Lakai}, \textit{Maskerade}
	\\
	\smallskip
	\\
	\emph{Tröpfe, Tränke, Trottel:} (Alchemisten + \textit{Die Intrige}:)
	\\
	Apotheker, Golem, Lehrling, Vision, \textit{Adlige}, \textit{Baron}, \textit{Eisenhütte}, \textit{Handelsposten}, \textit{Kupferschmied}, \textit{Wunschbrunnen}
	\\
	}
\end{tikzpicture}
\hspace{1cm}

	    % Basic settings for this card set
\renewcommand{\cardcolor}{prosperity}
\renewcommand{\cardextension}{Erweiterung III}
\renewcommand{\cardextensiontitle}{Blütezeit}

\clearpage
\newpage
\section{\cardextension \ - \cardextensiontitle}

\begin{tikzpicture}
	\card
	\cardstrip
	\cardbanner{banner/gold.png}
	\cardicon{banner/coin.png}
	\cardprice{5}
	\cardtitle{Abenteuer}
	\cardcontent{\emph{Errata:} Die Reihenfolge der beiden letzten Anweisungen sollte lauten: \enquote{Lege die übrigen aufgedeckten Karten ab. Lege diese Geldkarte aus.} 

	\smallskip

	Diese Geldkarte hat den Wert \coin{1}, wie ein Kupfer. Wenn du diese Karte auslegst, deckst du solange Karten von deinem Nachziehstapel auf, bis du eine Geldkarte aufgedeckt hast. Die auf diese Weise aufgedeckte Geldkarte legst du sofort aus. Die übrigen gerade aufgedeckten Karten legst du auf deinen Ablagestapel. Wenn du auch nach dem Mischen keine Geldkarte aufdecken kannst, legst du alle aufgedeckten Karten ab. Enthält die gerade ausgelegte Geldkarte zusätzliche Anweisungen, so führst du diese Anweisungen nun aus. Wenn du z. B. durch das Abenteuer ein weiteres Abenteuer aufdeckst, deckst du erneut Karten auf, bis du eine weitere Geldkarte \enquote{findest}. Denke daran, dass du die Geldkarten aus deiner Hand in beliebiger Reihenfolge auslegen kannst. Hast du z. B. das Abenteuer und den Lohn auf der Hand, darfst du wählen, welche der Karten du zuerst auslegst.}
\end{tikzpicture}
\hspace{-0.6cm}
\begin{tikzpicture}
	\card
	\cardstrip
	\cardbanner{banner/white.png}
	\cardicon{banner/coin.png}
	\cardprice{4}
	\cardtitle{Arbeiterdorf}
	\cardcontent{Du ziehst zuerst eine Karte nach, dann darfst du 2 weitere Aktionskarten auslegen und in der Kaufphase eine weitere Karte kaufen.}
\end{tikzpicture}
\hspace{-0.6cm}
\begin{tikzpicture}
	\card
	\cardstrip
	\cardbanner{banner/white.png}
	\cardicon{banner/coin.png}
	\cardprice{7}
	\cardtitle{Ausbau}
	\cardcontent{Du kannst den ausgelegten Ausbau selbst nicht entsorgen, da du die Karte nicht mehr auf der Hand hast, wenn du die Anweisung ausführst. Hast du einen weiteren Ausbau auf der Hand, kannst du diesen jedoch entsorgen. Hast du keine Karte auf der Hand, die du entsorgen könntest, darfst du dir auch keine neue Karte nehmen. Die neue Karte darf bis zu \coin{3} mehr kosten, als die entsorgte Karte. Du darfst keine Geldkarten oder virtuelles Geld verwenden, um den Betrag zu erhöhen. Du darfst auch eine Karte nehmen, die weniger kostet oder eine der entsorgten identische Karte nehmen. Die Karte, die du nimmst, legst du auf deinen Ablagestapel.}
\end{tikzpicture}
\hspace{-0.6cm}
\begin{tikzpicture}
	\card
	\cardstrip
	\cardbanner{banner/gold.png}
	\cardicon{banner/coin.png}
	\cardprice{7}
	\cardtitle{Bank}
	\cardcontent{Diese Geldkarte hat einen variablen Wert. Die Bank hat einen Wert von \coin{1} für jede Geldkarte, die zu diesem Zeitpunkt im Spiel ist. Die Bank selbst wird dabei mitgezählt. Hierfür ist es wichtig, dass du Geldkarten in beliebiger Reihenfolge auslegen kannst. Wenn du die Bank alleine ausliegen hast, hat sie einen Wert von \coin{1}. Wenn du 2 Banken hintereinander auslegst, hat die erste einen Wert von \coin{1}, die zweite Bank einen Wert von \coin{2}. Der Wert jeder Bank wird sofort nach dem Auslegen bestimmt und ändert sich nicht mehr z. B. durch das Auslegen weiterer Geldkarten.}
\end{tikzpicture}
\hspace{-0.6cm}
\begin{tikzpicture}
	\card
	\cardstrip
	\cardbanner{banner/white.png}
	\cardicon{banner/coin.png}
	\cardprice{4}
	\cardtitle{Bischof}
	\cardcontent{Für diese Karte werden die Punkte-Marker benötigt (siehe neue Regeln). Zuerst erhältst du +\coin{1} und 1 Punkte-Marker auf dein Tableau. Dann musst du eine Karte entsorgen, wenn du noch eine Karte auf der Hand hast. Du darfst dir eine Anzahl Punkte-Marker nehmen, die der Hälfte der Geld-Kosten der entsorgten Karte entspricht. Es werden dabei nur die Geld-Kosten beachtet. Trank-Kosten werden ignoriert. Wenn du z. B. den Golem (\coin{4}, \potion) (Dominion – Die Alchemisten) entsorgst, nimmst du dir 2 zusätzliche Punkte-Marker und legst diese auf dein Tableau. Danach dürfen deine Mitspieler reihum jeweils eine Karte entsorgen. Sie erhalten jedoch keine Punkte-Marker.}
\end{tikzpicture}
\hspace{-0.6cm}
\begin{tikzpicture}
	\card
	\cardstrip
	\cardbanner{banner/white.png}
	\cardicon{banner/coin.png}
	\cardprice{4}
	\cardtitle{Denkmal}
	\cardcontent{Für diese Karte werden die Punkte-Marker benötigt (siehe neue Regeln).}
\end{tikzpicture}
\hspace{-0.6cm}
\begin{tikzpicture}
	\card
	\cardstrip
	\cardbanner{banner/white.png}
	\cardicon{banner/coin.png}
	\cardprice{5}
	\cardtitle{Gesindel}
	\cardcontent{Du ziehst zunächst 3 Karten nach. Dann muss, beginnend mit dem Spieler links von dir, jeder Mitspieler 3 Karten von seinem Nachziehstapel aufdecken. Wenn er (auch nach dem eventuell nötigen Mischen seines Ablagestapels) nur weniger als 3 Karten aufdecken kann, deckt er nur so viele auf, wie möglich. Dann legt er alle auf diese Weise aufgedeckten Aktions- und Geldkarten auf seinen Ablagestapel. Karten mit kombinierten Kartentypen (z. B. Harem, Dominion – Die Intrige) sind davon auch betroffen. Die übrigen der gerade aufgedeckten Karten legt er in beliebiger Reihenfolge zurück auf seinen Nachziehstapel.}
\end{tikzpicture}
\hspace{-0.6cm}
\begin{tikzpicture}
	\card
	\cardstrip
	\cardbanner{banner/white.png}
	\cardicon{banner/coin.png}
	\cardprice{5}
	\cardtitle{Gewölbe}
	\cardcontent{Zuerst ziehst du 2 Karten nach. Dann legst du \enquote{eine beliebige Anzahl} Handkarten ab. Du kannst auch 0 Karten ablegen. Du kannst auch Karten ablegen, die du gerade nachgezogen hast. Beginnend mit dem Spieler links von dir, entscheidet sich dann jeder Mitspieler, ob er Handkarten ablegen möchte oder nicht. Wenn er sich dafür entscheidet Karten abzulegen, muss er genau 2 Karten ablegen und zieht dann eine Karte nach. Der Spieler kann sich auch für das Ablegen entscheiden, wenn er nur 1 Karte auf der Hand hat. In diesem Fall legt er diese ab, zieht dann allerdings keine Karte nach. Der Spieler darf sich nicht dafür entscheiden, nur 1 Karte abzulegen, wenn er 2 oder mehr Karten auf der Hand hat.}
\end{tikzpicture}
\hspace{-0.6cm}
\begin{tikzpicture}
	\card
	\cardstrip
	\cardbanner{banner/white.png}
	\cardicon{banner/coin.png}
	\cardprice{6}
	\cardtitle{\scriptsize{Großer Markt}}
	\cardcontent{Du musst nicht alle Geldkarten aus deiner Hand auslegen. Kupfer, die du auf der Hand hast, verbieten nicht, den Großen Markt zu kaufen. Kupfer, die in diesem Zug im Spiel waren, es nun aber nicht mehr sind, verbieten dir auch nicht, den Großen Markt zu kaufen. Hast zu z. B. 2 Käufe und 11 Kupfer auslegen und kaufst zuerst den Münzer, so entsorgst du nach diesem Kauf alle Geldkarten im Spiel. \emph{Nun hast du keine Kupfer mehr im Spiel und kannst den Großen Markt kaufen} (siehe auch: Münzer). Wenn du den Großen Markt auf eine andere Art nimmst, z. B. durch den Ausbau, so verhindern Kupfer, die im Spiel sind, dies auch nicht. Andere Geldkarten als Kupfer verbieten dir nicht, den Großen Markt zu kaufen, selbst wenn diese auch \coin{1} wert sind, wie z. B. Lohn.}
\end{tikzpicture}
\hspace{-0.6cm}
\begin{tikzpicture}
	\card
	\cardstrip
	\cardbanner{banner/white.png}
	\cardicon{banner/coin.png}
	\cardprice{6}
	\cardtitle{\scriptsize{Halsabschneider}}
	\cardcontent{Für diese Karte werden die Punkte-Marker benötigt (siehe neue Regeln). Zunächst erhältst du +1 Kauf und +\coin{2}. Dann müssen deine Mitspieler reihum ihre Handkarten auf 3 reduzieren. Für jede Karte, die du in dieser Runde kaufst, nimmst du dir einen Punkte-Marker und legst diesen auf dein Tableau. Hast du mehrere Handlanger im Spiel, erhältst du für jeden Handlanger einen Punkte-Marker pro gekaufter Karte. Hast du jedoch den Handlanger auf den Königshof gespielt, erhältst du dafür nur einen Punkte-Marker, weil die Karte nur einmal im Spiel ist.}
\end{tikzpicture}
\hspace{-0.6cm}
\begin{tikzpicture}
	\card
	\cardstrip
	\cardbanner{banner/white.png}
	\cardicon{banner/coin.png}
	\cardprice{3}
	\cardtitle{\footnotesize{Handelsroute}}
	\cardcontent{\tiny{\begin{Spacing}{1}
	\vspace{1em}
	Du erhältst einen zusätzlichen Kauf in deiner Kaufphase und +\coin{1} für jeden Geld-Marker, der auf dem Tableau Handelsroute liegt, wenn du die Karte auslegst. Dann musst du eine Karte aus deiner Hand entsorgen (wenn du noch eine Karte auf der Hand hast). 
 
	\emph{Spielvorbereitung:} Wenn die Handelsroute im Spiel verwendet wird (entweder als eine der 10 Königreichkarten oder im Schwarzmarktstapel), wird bei der Spielvorbereitung das Tableau Handelsroute bereit gelegt. Zusätzlich wird auf jeden Punktekartenstapel im gesamten Vorrat ein Geld-Marker gelegt. Punktekarten sind: die Basiskarten: Anwesen, Herzogtum, Provinz und Kolonie, die Königreichkarten: z. B. Gärten (Dominion – Basisspiel) und Herzog (Dominion – Die Intrige), kombinierte Karten: z. B. Harem und Adelige (Dominion – Die Intrige). Punktekarten, die nicht im Vorrat sind, werden nicht beachtet, z. B. aus dem Schwarzmarktstapel (Promokarte Schwarzmarkt). Auf den Ritter-Stapel (Dominion – Dark Ages) wird kein Marker gelegt, auch wenn die oberste Karte eine Punktekarte ist.
 
	\emph{Tableau Handelsroute:} Wird im Spielverlauf jeweils die erste Karte eines Punktekartenstapels genommen oder gekauft, legt der Spieler den Geld-Marker von diesem Stapel auf das Tableau Handelsroute. Legt ein Spieler die Handelsroute aus, so erhält er +1 virtuelles Geld für jeden Geld-Marker, der zu diesem Zeitpunkt auf dem Tableau liegt. Dabei ist egal, welcher Spieler den Geld-Marker auf das Tableau gelegt hat. Es wird nur beim Kaufen oder Nehmen der ersten Karte jedes Punktekartenstapels ein Geld-Marker auf das Tableau Handelsroute gelegt, also nur wenn der Geld-Marker, der bei Spielaufbau auf den Stapel gelegt wurde, noch dort liegt. Die Geld-Marker bleiben für den Rest des Spiels auf dem Tableau. Geld-Marker werden nicht vom Tableau entfernt und während des Spiels werden keine neuen Geld-Marker auf die Punktekartenstapel gelegt. Z. B. wenn ein Punktekartenstapel im Vorrat wieder durch den Botschafter (Dominion – Seaside) aufgefüllt wird.
	\end{Spacing}}}
\end{tikzpicture}
\hspace{-0.6cm}
\begin{tikzpicture}
	\card
	\cardstrip
	\cardbanner{banner/white.png}
	\cardicon{banner/coin.png}
	\cardprice{8*}
	\cardtitle{Hausierer}
	\cardcontent{Normalerweise kostet diese Karte \coin{8}. In der Kaufphase kostet der Hausierer pro Aktionskarte, die du selbst im Spiel hast, \coin{2} weniger. Das betrifft auch Hausierer auf der Hand oder in den Stapeln aller Spieler. Die Kosten sinken niemals unter \coin{0}. Wenn du z. B. das Arbeiterdorf auf den Königshof spielst, hast du 2 Aktionskarten im Spiel, obwohl du das Arbeiterdorf dreimal ausgespielt hast. Wenn du den Hausierer ausserhalb der üblichen Kaufphase (z. B. durch den Schwarzmarkt) kaufst oder nimmst, kostet der Hausierer \coin{8}.}
\end{tikzpicture}
\hspace{-0.6cm}
\begin{tikzpicture}
	\card
	\cardstrip
	\cardbanner{banner/gold.png}
	\cardicon{banner/coin.png}
	\cardprice{6}
	\cardtitle{Hort}
	\cardcontent{Diese Geldkarte hat einen Wert von \coin{2}, wie ein Silber. Wenn diese Karte im Spiel ist, und du eine Punktekarte kaufst, nimmst du dir zusätzlich ein Gold und legst es auf deinen Ablagestapel. Wenn kein Gold mehr im Vorrat ist, nimmst du nichts. Wenn du mehrere Karten Hort im Spiel hast, erhältst du mehrere Gold für den Kauf einer Punktekarte. Wenn du z. B. +1 Kauf hast und 2 Karten Hort auslegst, könntest du 2 Anwesen kaufen und dazu 4 Gold bekommen. Auch kombinierte Karten, wie z. B. Adlige und Harem (Dominion – Die Intrige) sind Punktekarten. Du nimmst dir ein Gold, auch wenn du z. B. den Wachturm vorzeigst und die gekaufte Punktekarte sofort entsorgst. Du nimmst dir nur ein Gold, wenn du eine Punktekarte tatsächlich kaufst, wenn du sie auf eine andere Weise nimmst, bekommst du kein Gold.}
\end{tikzpicture}
\hspace{-0.6cm}
\begin{tikzpicture}
	\card
	\cardstrip
	\cardbanner{banner/gold.png}
	\cardicon{banner/coin.png}
	\cardprice{5}
	\cardtitle{\tiny{Königliches Siegel}}
	\cardcontent{Diese Geldkarte hat den Wert \coin{2}, wie ein Silber. Wenn du mehrere Karten nimmst oder kaufst, kannst du für jede Karte getrennt entscheiden, ob du sie auf den Ablagestapel oder auf den Nachziehstapel legst. Ist dein Nachziehstapel leer und du entscheidest dich, eine Karte auf den Nachziehstapel zu legen, so wird dies die einzige Karte deines Nachziehstapels. Karten, die du durch die Besessenheit (Dominion – Die Alchemisten) im Zug eines anderen Spielers erhältst, darfst du nicht auf den Nachziehstapel legen, da das Königliche Siegel nur für den besessenen Spieler gilt.}
\end{tikzpicture}
\hspace{-0.6cm}
\begin{tikzpicture}
	\card
	\cardstrip
	\cardbanner{banner/white.png}
	\cardicon{banner/coin.png}
	\cardprice{7}
	\cardtitle{Königshof}
	\cardcontent{\emph{Errata:} Der Text auf der Karte sollte heißen: \enquote{Du darfst eine Aktionskarte aus deiner Hand wählen. Spiele diese Aktionskarte dreimal aus.} 

	\smallskip
	
	Diese Karte funktioniert ähnlich wie der Thronsaal (Dominion – Basisspiel), mit dem Unterschied, dass du die gewählte Karte 3mal spielst. Du wählst also eine Aktionskarte aus deiner Hand, legst sie aus, führst sie komplett aus, nimmst sie zurück auf die Hand, legst sie erneut aus, führst die Anweisungen nochmals aus, nimmst sie wieder zurück auf die Hand und legst sie dann ein drittes Mal aus. Dieses dreimalige Auslegen verbraucht keine Aktionen. (Das Auslegen des Königshofs verbraucht eine Aktion.) Du darfst keine anderen Aktionskarten auslegen, bis du die gewählte Karte 3mal ausgelegt und ausgeführt hast, ausser die Anweisung auf der gewählte Karte erlaubt es explizit, wie es z. B. der Königshof selbst tut. Wenn du eine Karte wählst, die +1 Aktion gibt, hast du am Ende +3 Aktionen. Legst du einen Königshof auf einen anderen Königshof aus, so wählst du nacheinander 3 Karten aus und legst jede gewählte Karte 3mal aus. Du legst also nicht eine Karte 9mal aus.}
\end{tikzpicture}
\hspace{-0.6cm}
\begin{tikzpicture}
	\card
	\cardstrip
	\cardbanner{banner/white.png}
	\cardicon{banner/coin.png}
	\cardprice{7}
	\cardtitle{\footnotesize{Kunstschmiede}}
	\cardcontent{Du darfst eine beliebige Anzahl Karten aus deiner Hand entsorgen. Das bedeutet, du darfst auch 0 Karten entsorgen. Dafür musst du dir eine Karte nehmen, die 0 Geld kostet. Dies Anweisung unterscheidet sich von Karten wie z. B. Ausbau, weil hier die Summe der Kosten aller entsorgten Karten betrachtet wird, nicht die Anzahl oder der Wert der einzelnen Karten. Ist im Vorrat keine Karte, die exakt soviel kostet, wie die entsorgten Karten, darfst du dir keine Karte nehmen. Es werden dabei nur die Geld-Kosten beachtet. Trank-Kosten (Dominion – Die Alchemisten) werden ignoriert. Du zählst weder Trank zur Summe, noch darfst du eine Karte mit Trank-Kosten nehmen.}
\end{tikzpicture}
\hspace{-0.6cm}
\begin{tikzpicture}
	\card
	\cardstrip
	\cardbanner{banner/white.png}
	\cardicon{banner/coin.png}
	\cardprice{5}
	\cardtitle{Leihaus}
	\cardcontent{Diese Karte erlaubt dir, deinen Ablagestapel durchzusehen (was normalerweise nicht erlaubt ist). Du darfst deinen Ablagestapel nur durchsehen, wenn du das Leihhaus gerade ausgelegt hast. Du musst deinen Mitspielern nicht die gesamten Karten deines Ablagestapels zeigen, nur die Kupfer, die du auf die Hand nimmst. Nachdem du die Kupfer auf die Hand genommen hast, legst du die übrigen Karten in beliebiger Reihenfolge auf den Ablagestapel zurück.}
\end{tikzpicture}
\hspace{-0.6cm}
\begin{tikzpicture}
	\card
	\cardstrip
	\cardbanner{banner/gold.png}
	\cardicon{banner/coin.png}
	\cardprice{3}
	\cardtitle{Lohn}
	\cardcontent{(Wir haben bei der Übersetzung von \enquote{Loan} aus dem englischen Original bewusst \enquote{Lohn} gewählt, da uns dieser Name als passender erscheint, als eine sprachlich korrekte Übersetzung.) Diese Geldkarte hat einen Wert von \coin{1}, wie Kupfer. Wenn du sie auslegst, deckst du solange Karten von deinem Nachziehstapel auf, bis du eine Geldkarte aufgedeckt hast. Dann entscheidest du dich, ob du die aufgedeckte Geldkarte entsorgst oder ablegst. Danach legst du alle übrigen aufgedeckten Karten ab. Wenn du (auch nach dem Mischen deines Ablagestapels) keine Karten mehr aufdecken kannst, legst du alle aufgedeckten Karten ab. Beachte, dass du deine Geldkarten in beliebiger Reihenfolge auslegen kannst und nicht alle Geldkarten aus deiner Hand auslegen musst.}
\end{tikzpicture}
\hspace{-0.6cm}
\begin{tikzpicture}
	\card
	\cardstrip
	\cardbanner{banner/white.png}
	\cardicon{banner/coin.png}
	\cardprice{5}
	\cardtitle{Münzer}
	\cardcontent{Wenn du diese Karte kaufst, entsorgst du alle Geldkarten, die zu diesem Zeitpunkt im Spiel sind und nur diese (nicht aus deiner Hand oder sonst woher). Bedenke, dass du nicht alle Geldkarten auslegen musst. Wenn du in dieser Runde mehrere Karten kaufst, entsorgst du die Geldkarten im Spiel, direkt nachdem du den Münzer gekauft hast. Geldkarten, die du in dieser Runde ausgelegt hast, haben aber bereits Geld \enquote{produziert}, auch wenn du sie entsorgst. Du kannst also für den gesamten Wert der ausgelegten Geldkarten neue Karten kaufen. Du kannst allerdings keine weiteren Geldkarten mehr auslegen, sobald du eine Karte gekauft hast. Wenn du den Münzer auslegst, darfst du eine Geldkarte aus deiner Hand aufdecken. Dann nimmst du dir sofort eine identische Karte aus dem Vorrat und legst diese auf deinen Ablagestapel. Die aufgedeckte Geldkarte nimmst du zurück auf die Hand. Die aufgedeckte Geldkarte kann auch eine kombinierte Karte sein, z. B. Harem (Dominion – Die Intrige). Wenn du den Münzer kaufst und den Wachturm aus deiner Hand aufdeckst, darfst du den Münzer sofort auf deinen Nachziehstapel legen. Die ausgelegten Geldkarten müssen jedoch entsorgt werden. Wenn du den Münzer kaufst, während das Königliche Siegel im Spiel ist, wird das Königliche Siegel entsorgt, bevor du den Münzer zurück auf deinen Nachziehstapel legen dürftest.}
\end{tikzpicture}
\hspace{-0.6cm}
\begin{tikzpicture}
	\card
	\cardstrip
	\cardbanner{banner/white.png}
	\cardicon{banner/coin.png}
	\cardprice{5}
	\cardtitle{Quacksalber}
	\cardcontent{Beginnend mit dem Spieler links von dir muss jeder Mitspieler entscheiden, ob er einen Fluch aus seiner Hand ablegt oder sich ein Kupfer und einen Fluch vom Vorrat nimmt und auf seinen Ablagestapel legt. Er darf sich auch für die zweite Möglichkeit entscheiden, wenn einer oder beide Stapel leer sind. In diesem Fall nimmt er sich nur eine der noch vorhandenen Karten oder, wenn beide Stapel leer sind, keine Karte. Deckt ein Spieler einen Burggraben (Dominion – Basisspiel) aus seiner Hand auf, darf er weder eine Karte nehmen noch eine Karte ablegen. Er kann nicht nur einen Teil des Angriffs abwehren. Entscheidet sich ein Spieler für die zweite Möglichkeit und deckt einen Wachturm aus seiner Hand auf, darf er sofort eine oder beide Karten entsorgen.}
\end{tikzpicture}
\hspace{-0.6cm}
\begin{tikzpicture}
	\card
	\cardstrip
	\cardbanner{banner/gold.png}
	\cardicon{banner/coin.png}
	\cardprice{5}
	\cardtitle{\scriptsize{Schmuggelware}}
	\cardcontent{Diese Geldkarte hat einen Wert von \coin{3}, wie ein Gold. Wenn du sie ausspielst, erhältst du zunächst +1 Kauf. Dann benennt der Spieler links von dir eine Karte, die du in dieser Runde nicht kaufen darfst. Er kann auch eine Karte benennen, die nicht im Vorrat ist, aber z. B. im Schwarzmarkt-Stapel. Wenn du mehrere Karten Schmuggelware auslegst, benennt der Spieler links von dir jedesmal eine Karte. Du darfst in diesem Zug keine der benannten Karten kaufen. Hierfür ist es wichtig, dass du Geldkarten in beliebiger Reihenfolge auslegen kannst. Du kannst also z. B. zuerst eine Schmuggelware auslegen, dann benennt der Spieler eine Karte, danach kannst du weitere Geldkarten auslegen. Die Anzahl der Karten, die ein Spieler noch auf der Hand hält, ist für die übrigen Spieler sichtbar. Du darfst die benannte Karte nicht kaufen, du darfst sie aber nehmen, wenn dir dies eine Anweisung einer anderen Karte erlaubt, z. B. Hort. Beachte, dass du in dieser Runde keine weiteren Geldkarten mehr auslegen darfst, sobald du eine Karte gekauft hast.}
\end{tikzpicture}
\hspace{-0.6cm}
\begin{tikzpicture}
	\card
	\cardstrip
	\cardbanner{banner/white.png}
	\cardicon{banner/coin.png}
	\cardprice{5}
	\cardtitle{Stadt}
	\cardcontent{\emph{Errata:} Auf der Karte Stadt hat sich ein Tippfehler eingeschlichen. Statt \enquote{1+ Kauf} sollte es heißen \enquote{+1 Kauf}.

	\smallskip
	
	Du ziehst zuerst eine Karte nach und erhältst +2 Aktionen. Wenn mindestens ein Stapel im Vorrat leer ist, ziehst du eine weitere Karte nach. Wenn genau 1 Stapel im Vorrat leer ist, erhältst du nur +1 Karte. Oder wenn 2 oder mehr Stapel im Vorrat leer sind, erhältst du +1 Karte, +\coin{1} und +1 Kauf. Es gibt keine weiteren Boni, wenn 3 oder mehr Stapel im Vorrat leer sind. Die jeweiligen Bedingungen müssen beim Ausspielen der Karte erfüllt sein. 

	\smallskip
	 
	Der Effekt der Karte wird nicht rückwirkend verändert, wenn ein Stapel im Vorrat später leer wird (z. B. durch den Ausbau) oder auch wieder aufgefüllt wird (z. B. durch den Botschafter, Dominion-Seaside). 

	\smallskip
	
	Der Müllstapel ist nicht Teil des Vorrats und wird somit nicht beachtet. Sind beim Ausspielen der Karte also z. B. 2 Stapel leer, so erhältst du insgesamt +2 Karten, +2 Aktionen, +\coin{1} und +1 Kauf.}
\end{tikzpicture}
\hspace{-0.6cm}
\begin{tikzpicture}
	\card
	\cardstrip
	\cardbanner{banner/gold.png}
	\cardicon{banner/coin.png}
	\cardprice{4}
	\cardtitle{Steinbruch}
	\cardcontent{Diese Geldkarte hat der Wert 1 Geld, wie ein Kupfer. Wenn diese Karte im Spiel ist, kosten Aktionskarten 2 Geld weniger. Dieser Effekt ist kumulativ. Wenn du z. B. in der Kaufphase 2 Steinbrüche ausspielst, kosten Aktionskarten um 4 Geld weniger. Der Effekt kann auch durch andere Karten ergänzt werde. Spielst du z. B. in der Aktionsphase ein Arbeiterdorf und in der Kaufphase 2 Steinbrüche, kostet der Hausierer nur noch 2 Geld. Handkarten, Karten im Nachziehstapel und im Ablagestapel sind auch durch diesen Effekt betroffen. Kombinierte Karten, wie z. B. Adelige (Dominion – Die Intrige) sind auch Aktionskarten.}
\end{tikzpicture}
\hspace{-0.6cm}
\begin{tikzpicture}
	\card
	\cardstrip
	\cardbanner{banner/gold.png}
	\cardicon{banner/coin.png}
	\cardprice{4}
	\cardtitle{Talisman}
	\cardcontent{Diese Geldkarte hat den Wert 1 Geld, wie ein Kupfer. Wenn diese Karte im Spiel ist und du eine Karte kaufst, die 4 Geld oder weniger kostet und die keine Punktekarte ist, nimmst du dir zusätzlich eine weitere identische Karte. Du nimmst diese Karte vom Vorrat und legst sie auf deinen Ablagestapel. Gibt es keine weitere Karte mit diesem Namen im Vorrat, nimmst du dir nichts. Hast du mehrere Talismane im Spiel, nimmst du dir für jeden Talisman eine zusätzliche Karte. Wenn du mehrere Karten kaufst, auf diese Bedingungen zutreffen (4 Geld oder weniger, keine Punktekarte), nimmst du dir von jeder eine weitere Karte. Wenn du z. B. 2 Talismane, 4 Kupfer und 2 Käufe hast und dir ein Silber und eine Handelsroute kaufst, nimmst du dir zusätzlich 2 weitere Silber und 2 weitere Handelsrouten. Der Talisman wirkt nur für Karten, die du kaufst. Nimmst du dir eine Karte auf eine andere Weise, z. B. durch den Ausbau, erhältst du keine weitere. Kombinierte Karten, wie z. B. die Große Halle (Dominion – Die Intrige) sind auch Punktekarten. Bei den Kosten von 4 Geld oder weniger werden die aktuell geltenden Kosten betrachtet. Das sind nicht unbedingt die aufgedruckten Kosten. Wenn du z. B. in deiner Aktionsphase 2 Aktionskarten ausgelegt hast, kostet der Hausierer nur noch 4 Geld und du darfst dir somit einen weiteren Hausierer nehmen.}
\end{tikzpicture}
\hspace{-0.6cm}
\begin{tikzpicture}
	\card
	\cardstrip
	\cardbanner{banner/blue.png}
	\cardicon{banner/coin.png}
	\cardprice{3}
	\cardtitle{Wachturm}
	\cardcontent{\tiny{\begin{Spacing}{1}
	Wenn du diese Karte in deinem Zug auslegst, ziehst du solange Karten von deinem Nachziehstapel, bis du 6 Karten auf der Hand hast. Hast du nach dem Ausspielen des Wachturmes bereits 6 oder mehr Karten auf der Hand, ziehst du keine Karte nach. Immer wenn du eine Karte nimmst oder kaufst, egal ob in deinem eigenen Zug oder im Zug eines Mitspielers, darfst du den Wachturm aus deinem Hand aufdecken und dann entscheiden, ob du die neue Karte entsorgst oder oben auf deinen Nachziehstapel legst. Du darfst den Wachturm jedesmal, wenn du eine Karte nimmst oder kaufst aus deiner Hand aufdecken. Wie üblich bei Reaktionskarten, deckst du den Wachturm nur auf und nimmst ihn dann zurück auf deine Hand. Spielt ein Mitspieler z. B. den Quacksalber, kannst du den Wachturm nur bei einer oder bei beiden Karten (Kupfer und Fluch) aufdecken und für jede Karte getrennt entscheiden. Du kannst mit dem Wachturm auch nacheinander auf mehrere Angriffe unterschiedlicher Mitspieler reagieren und ihn dann in deinem eigenen Zug nochmals einsetzen. Auch wenn du dich dafür entscheidest eine Karte, die du gerade genommen hast oder nehmen musstest, zu entsorgen, musst du diese Karte zuerst nehmen. Die Karte ist also nicht mehr im Vorrat, und andere Karten, die auf genommene Karten Bezug nehmen, wie z. B. die Schmuggler (Dominion – Seaside), können auch darauf angewandt werden. Wenn während des Extrazuges durch die Besessenheit (Dominion – Die Alchemisten) Karten genommen oder gekauft werden, kannst du den Wachturm nicht aus deiner Hand aufdecken, da der besessene Spieler die Karte nimmt. Du kannst den Wachturm auch aus deiner Hand aufdecken, wenn du eine Karte nimmst, die du nicht wie üblich auf deinen Ablagestapel legst, wie z. B. die Mine (Dominion – Basisspiel).
	\end{Spacing}}}
\end{tikzpicture}
\hspace{-0.6cm}
\begin{tikzpicture}
	\card
	\cardstrip
	\cardbanner{banner/gold.png}
	\cardicon{banner/coin.png}
	\cardprice{9}
	\cardtitle{Platin}
	\cardcontent{}
\end{tikzpicture}
\hspace{-0.6cm}
	\begin{tikzpicture}
	\card
	\cardstrip
	\cardbanner{banner/green.png}
	\cardicon{banner/coin.png}
	\cardprice{11}
	\cardtitle{Kolonien}
	\cardcontent{}
\end{tikzpicture}
\hspace{-0.6cm}
\begin{tikzpicture}
	\card
	\cardstrip
	\cardbanner{banner/white.png}
	\cardtitle{\scriptsize{Empfohlene 10er Sätze\qquad}}
	\cardcontent{\emph{Blütezeit und Basisspiel:}

	\smallskip 
	
	\emph{Haufenweise Geld:} \\ 
	Abenteuer, Bank, Großer Markt, Königliches Siegel, Münzer, Abenteurer, Geldverleiher, Laboratorium, Mine, Spion 

	\smallskip 
	
	\emph{Die Armee des Königs:} \\ 
	Ausbau, Gesindel, Gewölbe, Handlanger, Königshof, Bürokrat, Burggraben, Dorf, Ratsversammlung, Spion 

	\smallskip 
	
	\emph{Ein gutes Leben:} \\ 
	Denkmal, Hort, Leihaus, Quacksalber, Schmuggelware, Bürokrat, Dorf, Gärten, Kanzler, Keller}
\end{tikzpicture}
\hspace{-0.6cm}
\begin{tikzpicture}
	\card
	\cardstrip
	\cardbanner{banner/white.png}
	\cardtitle{\scriptsize{Empfohlene 10er Sätze\qquad}}
	\cardcontent{\emph{Blütezeit und Die Intrige:}

	\smallskip 
	
	\emph{Pfade zum Sieg:} \\ 
	Bischof, Denkmal, Halsabschneider, Hausierer, Leihaus, Anbau, Armenviertel, Baron, Handlanger, Harem
	
	\smallskip 
	
	\emph{All along the watchtower:} \\ 
	Gewölbe, Handelsroute, Hort, Talisman, Wachturm, Bergwerk, Brücke, Große Halle, Handlanger, Kerkermeister

	\smallskip 
	
	\emph{Glücksritter:} \\ 
	Ausbau, Bank, Gewölbe, Königshof, Kunstschmiede, Brücke, Kupferschmied, Tribut, Trickser, Wunschbrunnen}
\end{tikzpicture}
\hspace{-0.6cm}
\begin{tikzpicture}
	\card
	\cardstrip
	\cardbanner{banner/white.png}
	\cardtitle{\scriptsize{Empfohlene 10er Sätze\qquad}}
	\cardcontent{\emph{Blütezeit:}

	\smallskip 
	
	\emph{Anfänger:} \\ 
	Abenteuer, Arbeiterdorf, Ausbau, Bank, Denkmal, Gesindel, Halsabschneider, Königliches Siegel, Leihaus, Wachturm 

	\smallskip 
	
	\emph{Freundliche Interaktion:} \\ 
	Arbeiterdorf, Bischof, Gewölbe, Handelsroute, Hausierer, Hort, Königliches Siegel, Kunstschmiede, Schmuggelware, Stadt 

	\smallskip 
	
	\emph{Große Aktionen:} \\ 
	Ausbau, Gesindel, Gewölbe, Großer Markt, Königshof, Lohn, Münzer, Stadt, Steinbruch, Talisman}
\end{tikzpicture}
\hspace{0.6cm}

	    % Basic settings for this card set
\renewcommand{\cardcolor}{cornucopia}
\renewcommand{\cardextension}{Erweiterung IV}
\renewcommand{\cardextensiontitle}{Reiche Ernte}

\clearpage
\newpage
\section{\cardextension \ - \cardextensiontitle}

\begin{tikzpicture}
	\card
	\cardstrip
	\cardbanner{banner/white.png}
	\cardicon{banner/coin.png}
	\cardprice{4}
	\cardtitle{Bauerndorf}
	\cardcontent{Du erhältst zunächst +2 Aktionen. Dann deckst du solange Karten von deinem Nachziehstapel auf, bis entweder eine Geldkarte oder eine Aktionskarte offen liegt. Nimm diese Geld- oder Aktionskarte auf die Hand. Lege die übrigen aufgedeckten Karten auf deinen Ablagestapel. Du darfst nicht wählen, ob du eine Geldkarte oder eine Aktionskarte auf die Hand nehmen möchtest. Du musst die erste aufgedeckte Geld- oder Aktionskarte auf die Hand nehmen. Kombinierte Kartentypen sind gleichzeitig alle angegebenen Kartentypen. Du darfst die aufgenommene Karte auch in diesem Zug (nach den üblichen Regeln) ausspielen. Wenn du auch nach dem Mischen deines Ablagestapels keine Geld- oder Aktionskarte aufdecken kannst, nimmst du keine Karte auf die Hand.}
\end{tikzpicture}
\hspace{-1cm}
\begin{tikzpicture}
	\card
	\cardstrip
	\cardbanner{banner/white.png}
	\cardicon{banner/coin.png}
	\cardprice{5}
	\cardtitle{Ernte}
	\cardcontent{Decke die obersten 4 Karten von deinem Nachziehstapel auf. Du erhältst +\coin{1} für jede aufgedeckte Karte mit unterschiedlichem Namen. Deckst du z. B. 2 Kupfer, 1 Silber und 1 Anwesen auf, erhältst du +\coin{3}. Kannst du (auch nach dem Mischen des Ablagestapels) nur weniger als 4 Karten aufdecken, deckst du nur so viele auf, wie möglich.}
\end{tikzpicture}
\hspace{-1cm}
\begin{tikzpicture}
	\card
	\cardstrip
	\cardbanner{banner/green.png}
	\cardicon{banner/coin.png}
	\cardprice{6}
	\cardtitle{Festplatz}
	\cardcontent{Diese Königreichkarte ist eine Punktekarte, keine Aktionskarte. Sie hat bis zum Ende des Spiels keine Funktion. Bei der Wertung zählt sie 2 Siegpunkte für je volle 5 Karten mit unterschiedlichem Namen im gesamten Kartensatz (Nachziehstapel, Ablagestapel und Handkarten) des Spielers. Bei Spielende suchst du aus deinem gesamten Kartensatz je eine Karte jedes Namens heraus. Diese Karten zählst du und teilst die Anzahl durch 5. Das Ergebnis (abgerundet) multipliziert mit 2 ergibt die Punkte. Hast du 0-4 Karten mit unterschiedlichem Namen erhältst du keine Punkte, für 5-9 Karten 2 Punkte , für 10-14 Karten 4 Punkte usw.\\ \smallskip
	Im Spiel mit 2 Spielern werden 8 Karten Festplatz verwendet, bei 3 oder mehr Spielern 12 Karten.}
\end{tikzpicture}
\hspace{-1cm}
\begin{tikzpicture}
	\card
	\cardstrip
	\cardbanner{banner/gold.png}
	\cardicon{banner/coin.png}
	\cardprice{5}
	\cardtitle{Füllhorn}
	\cardcontent{Du spielst diese Karte, wie jede andere Geldkarte, in der Kaufphase aus. Das Füllhorn hat den Wert \coin{0}. Es bringt also kein Geld für den Kauf. Du darfst dir jedoch sofort, wenn du die Karte ausspielst, eine Karte aus dem Vorrat nehmen. Diese Karte darf bis zu \coin{1} pro Karte mit unterschiedlichem Namen, die du im Spiel hast, kosten. Karten, die du im Spiel hast sind: in diesem Zug ausgespielte Aktions- und Geldkarten (das Füllhorn selbst eingeschlossen) und bei dir ausliegende Dauerkarten (Dominion – Seaside). Nicht im Spiel sind Karten, die du nach dem ausspielen entsorgt hast (z. B. Festmahl, Dominion – Basisspiel). Nimmst du dir mit dem Füllhorn eine Punktekarte (auch eine kombinierte), entsorgst du das Füllhorn. Du entsorgst das Füllhorn nicht, wenn du dir eine Punktekarte auf eine andere Art nimmst oder kaufst.}
\end{tikzpicture}
\hspace{-1cm}
\begin{tikzpicture}
	\card
	\cardstrip
	\cardbanner{banner/white.png}
	\cardicon{banner/coin.png}
	\cardprice{5}
	\cardtitle{Harlekin}
	\cardcontent{Du erhältst zunächst +\coinj{2}. Dann muss, beginnend mit dem Spieler links von dir, reihum jeder Mitspieler die oberste Karte von seinem Nachziehstapel ablegen. Ist es eine Punktekarte (auch eine kombinierte), so muss er sich einen Fluch nehmen. Ist kein Fluch mehr im Vorrat, nimmt der Spieler keinen. Ist es keine Punktekarte, kannst du wählen: Entweder der Spieler muss sich eine Karte mit gleichem Namen aus dem Vorrat nehmen und auf seinen Ablagestapel legen oder du nimmst dir eine Karte mit gleichem Namen und legst sie auf deinen Ablagestapel. Ist keine Karte mit diesem Namen mehr im Vorrat nimmt keiner von beiden eine solche Karte.}
\end{tikzpicture}
\hspace{-1cm}
\begin{tikzpicture}
	\card
	\cardstrip
	\cardbanner{banner/white.png}
	\cardicon{banner/coin.png}
	\cardprice{4}
	\cardtitle{Junge Hexe}
	\cardcontent{
		\begin{singlespace}\tiny{
			\emph{Spielvorbereitung:} Wird die Junge Hexe im Spiel verwendet, so wird zusätzlich ein Bannstapel benötigt. (Dies gilt auch, wenn sich die Junge Hexe im Schwarzmarktstapel befindet – Promokarte: Schwarzmarkt.) Für den Bannstapel wählen die Spieler eine nicht verwendete Königreichkarte, die \coin{2} oder \coin{3} kostet, aus einer beliebigen Edition oder Erweiterung aus. Der komplette Stapel dieser Karte wird als zusätzlicher Stapel in den Vorrat gelegt. Die Platzhalterkarte Junge Hexe wird quer unter den Stapel geschoben, um diesen Stapel als Bannstapel zu kennzeichnen. In diesem Spiel werden also 11 unterschiedliche Königreichskarten verwendet. Zusätzlich zur üblichen (aufgedruckten) Funktion sind alle diese Karten Bannkarten. Der Bannstapel wird behandelt wie die anderen Königreichstapel. Er ist Teil des Vorrats, die Karten können wie üblich gekauft oder genommen werden und der Stapel wird für die Spielendebedingung berücksichtigt.\\ \smallskip
			Wenn du die Junge Hexe ausspielst, ziehst du zuerst 2 Karten nach, dann musst du 2 Karten aus deiner Hand ablegen. Beginnend mit dem Spieler links von dir darf nun reihum jeder Mitspieler eine Bannkarte aus seiner Hand aufdecken. Wenn er das nicht macht, muss er sich einen Fluch vom Vorrat nehmen. Ist kein Fluch mehr im Vorrat, muss der Spieler keinen nehmen. Die Spieler dürfen wie üblich auch Reaktionskarten (z. B. Pferdehändler) aus ihrer Hand aufdecken. Dies tun die Spieler, bevor sie eine Bannkarte aufdecken. Die Anweisungen auf der Reaktionskarte werden also ausgeführt, bevor der Angriff durch die Bannkarte abgewehrt wird. Ist die Bannkarte eine Reaktionskarte, kann sie zuerst aufgedeckt werden um die übliche Reaktion auszuführen und dann (wenn der Spieler sie dann noch auf der Hand hat) nochmals in ihrer Funktion als Bannkarte.
		} \end{singlespace}	
	}
\end{tikzpicture}
\hspace{-1cm}
\begin{tikzpicture}
	\card
	\cardstrip
	\cardbanner{banner/white.png}
	\cardicon{banner/coin.png}
	\cardprice{3}
	\cardtitle{Menagerie}
	\cardcontent{Zuerst erhältst du +1 Aktion. Dann deckst du deine Handkarten auf. Hast du nur Karten mit unterschiedlichem Namen auf der Hand, ziehst du 3 Karten nach. Hast du wenigstens eine Karte mit gleichem Namen auf der Hand (eine Karte doppelt), ziehst du 1 Karte nach. Hast du z. B. nur ein Kupfer und ein Silber auf der Hand, hast du nur Karten mit unterschiedlichem Namen (obwohl beide Geldkarten sind). Du darfst also 3 Karten nachziehen.}
\end{tikzpicture}
\hspace{-1cm}
\begin{tikzpicture}
	\card
	\cardstrip
	\cardbanner{banner/white.png}
	\cardicon{banner/coin.png}
	\cardprice{4}
	\cardtitle{Nachbau}
	\cardcontent{Zuerst entsorgst du eine Karte aus deiner Hand und nimmst dir dafür eine Karte, die genau \coin{1} mehr kostet. Dann entsorgst du eine andere Karte aus deiner Hand und nimmst dir wieder eine Karte, die genau \coin{1} mehr kostet. Du nimmst die Karten vom Vorrat und legst sie auf deinen Ablagestapel. Ist im Vorrat keine Karte, die genau 1 Geld mehr kostet als die entsorgte, nimmst du dir keine Karte.}
\end{tikzpicture}
\hspace{-1cm}
\begin{tikzpicture}
	\card
	\cardstrip
	\cardbanner{banner/blue.png}
	\cardicon{banner/coin.png}
	\cardprice{4}
	\cardtitle{Pferdehändler}
	\cardcontent{Wenn du diese Karte in deinem Zug ausspielst, erhältst du zuerst +1 Kauf und +\coin{3}. Dann musst du 2 Karten aus deiner Hand ablegen. Hast du weniger als 2 Karten auf der Hand, legst du nur so viele wie möglich ab.
	\\
	\smallskip
	\\				
	Spielt ein Mitspieler eine Angriffskarte, darfst du diese Karte aus deiner Hand vorzeigen und dann zur Seite legen. Dann wird der Angriff normal ausgeführt. Du kannst jeden Pferdehändler nur gegen einen Angriff verwenden, da du ihn danach nicht mehr auf der Hand hast. Du darfst mehrere Pferdehändler gegen den selben Angriff zur Seite legen. Zu Beginn deines nächsten Zuges nimmst du alle zur Seite gelegten Pferdehändler wieder auf die Hand und ziehst zusätzlich für jeden dieser Pferdehändler 1 Karte nach.}
\end{tikzpicture}
\hspace{-1cm}
\begin{tikzpicture}
	\card
	\cardstrip
	\cardbanner{banner/white.png}
	\cardicon{banner/coin.png}
	\cardprice{5}
	\cardtitle{Treibjagd}
	\cardcontent{Zuerst ziehst du eine Karte nach und erhältst +1 Aktion. Dann deckst du deine gesamten Handkarten auf und legst diese offen aus. Danach deckst du solange Karten von deinem Nachziehstapel auf, bis du eine Karte aufdeckst, die du nicht auf der (momentan offen ausliegenden) Hand hast. Nimm diese gerade aufgedeckte Karte und deine offen ausliegenden Handkarten zurück auf die Hand. Lege die übrigen aufgedeckten Karten auf deinen Ablagestapel. Kannst du (auch nach dem Mischen deines Ablagestapels) keine Karte aufdecken, von der du nicht bereits eine auf der Hand hast, legst du die aufgedeckten Karten ab und nimmst nur deine Handkarten zurück auf die Hand.}
\end{tikzpicture}
\hspace{-1cm}
\begin{tikzpicture}
	\card
	\cardstrip
	\cardbanner{banner/white.png}
	\cardicon{banner/coin.png}
	\cardprice{4}
	\cardtitle{Turnier}
	\cardcontent{
		\begin{singlespace}\tiny{
			\emph{Spielvorbereitung:} Wird das Turnier für das Spiel verwendet, so wird zusätzlich der Preisstapel benötigt. (Dies gilt auch, wenn sich das Turnier im Schwarzmarktstapel befindet – Promokarte: Schwarzmarkt.) Für den Preisstapel werden die 5 einmaligen Preiskarten benötigt. Alle Preiskarten werden als ein Stapel offen bereit gelegt. Der Preisstapel besteht nur aus diesen 5 Karten und ist nicht Teil des Vorrats. \\ \smallskip
			\emph{Errata:} Der Kartentext is falsch, es sollte „Jeder Spieler darf eine Provinz aus seiner Hand aufdecken. Wenn du das machst, lege die Provinz ab und nimm dir einen Preis (vom Preisstapel) oder ein Herzogtum und lege die neue Karte auf deinen Nachziehstapel. Wenn es kein Mitspieler macht: +1 Karte, +\coin{1}.“ statt „Du darfst eine Provinz aus deiner Hand ablegen. Wenn du das machst [...] Jeder Mitspieler darf...“ heißen. \\ \smallskip
			Zuerst erhältst du +1 Aktion. Dann darf jeder Spieler eine Provinz aus seiner Hand aufdecken. Wenn du das machst, lege die Provinz ab und nehme dir eine frei wählbare Preiskarte vom Preisstapel oder ein Herzogtum vom Vorrat und lege die neue Karte auf deinen Nachziehstapel. Du kannst dich auch für einen der beiden Stapel entscheiden, wenn dieser leer ist. Ist der Stapel für den du dich entscheidest leer, nimmst du dir keine Karte. Wenn keiner deiner Mitspieler eine Provinz aufgedeckt hat, erhältst du +1 Karte und +\coin{1}. Es gibt also 4 mögliche Ergebnisse: \\	\smallskip
	\begin{itemize}	
	\item Du legst keine Provinz ab und keiner deiner Mitspieler deckt eine Provinz auf: Du erhältst +1 Karte, +1 Aktion und +\coin{1}.
	\item Du legst keine Provinz ab und mindestens einer deiner Mitspieler deckt eine Provinz auf: Du erhältst +1 Aktion.
	\item Du legst eine Provinz ab aber keiner deiner Mitspieler deckt eine Provinz auf: Du nimmst dir eine Preiskarte oder ein Herzogtum und erhältst +1 Karte, +1 Aktion und +\coin{1}.
	\item Du legst eine Provinz ab und mindestens einer deiner Mitspieler deckt eine Provinz auf: Du nimmst dir eine Preiskarte oder ein Herzogtum und erhältst +1 Aktion.
	\end{itemize}
	\\
	\smallskip
	\\	
	Du darfst den Preisstapel jederzeit durchsehen.
	} \end{singlespace} }
\end{tikzpicture}
\hspace{-1cm}
\begin{tikzpicture}
	\card
	\cardstrip
	\cardbanner{banner/white.png}
	\cardicon{banner/coin.png}
	\cardprice{3}
	\cardtitle{Wahrsagerin}
	\cardcontent{Du erhältst zunächst +\coin{2}. Dann muss, beginnend mit dem Spieler links von dir, reihum jeder Mitspieler solange Karten von seinem Nachziehstapel aufdecken, bis entweder eine Punktekarte oder ein Fluch offen liegt. Er legt diese Punkte- oder Fluchkarte verdeckt auf seinen Nachziehstapel. Der Spieler darf nicht wählen, ob er eine Punktekarte oder einen Fluch auf den Nachziehstapel legt, er muss die erste aufgedeckte Punkte- oder Fluchkarte zurücklegen. Die übrigen aufgedeckten Karten legt er auf seinen Ablagestapel. Wenn der Spieler auch nach dem Mischen seines Ablagestapels keine Punkte- oder Fluchkarte aufdecken kann, legt er alle aufgedeckten Karten ab. Karten mit kombinierten Kartentypen sind gleichzeitig alle angegebenen Kartentypen.}
\end{tikzpicture}
\hspace{-1cm}
\begin{tikzpicture}
	\card
	\cardstrip
	\cardbanner{banner/white.png}
	\cardicon{banner/coin.png}
	\cardprice{2}
	\cardtitle{Weiler}
	\cardcontent{Zuerst ziehst du immer eine Karte nach und erhältst +1 Aktion. Nun darfst du eine Karte aus deiner Hand ablegen um entweder zusätzlich +1 Aktion oder +1 Kauf zu erhalten. Oder du legst 2 Karten aus deiner Hand ab und erhältst zusätzlich +1 Aktion und +1 Kauf. Du kannst jedoch nicht +2 Aktionen oder +2 Käufe zusätzlich erhalten. Du kannst auch darauf verzichten, Karten abzulegen und erhältst dann nichts weiter.}
\end{tikzpicture}
\hspace{-1cm}
\begin{tikzpicture}
	\card
	\cardstrip
	\cardbanner{banner/white.png}
	\cardicon{}
	\cardprice{}
	\cardtitle{\scriptsize{Empfohlene 10er Sätze\qquad}}
	\cardcontent{\emph{Reiche Ernte	und	Basisspiel:}
	\\
	\smallskip
	\\
	\emph{Kopfgeld:}
	\\
	Ernte, Füllhorn, Menagerie, Treibjagd, Turnier (+ Preiskarten) / Geldverleiher, Jahrmarkt, Keller, Miliz, Schmiede
	\\
	\smallskip
	\\
	\emph{Böses Omen:}
	\\
	Füllhorn, Harlekin, Nachbau, Wahrsagerin, Weiler / Abenteurer, Bürokrat, Laboratorium, Spion, Thronsaal
	\\
	\smallskip
	\\
	\emph{Wanderzirkus:}
	\\
	Bauerndorf, Festplatz, Harlekin, Junge Hexe (+ Kanzler als Bannstapel), Pferdehändler
	/ Festmahl, Laboratorium, Markt, Umbau, Werkstatt
	\\
	}
\end{tikzpicture}
\hspace{-1cm}
\begin{tikzpicture}
	\card
	\cardstrip
	\cardbanner{banner/white.png}
	\cardicon{}
	\cardprice{}
	\cardtitle{\scriptsize{Empfohlene 10er Sätze\qquad}}
	\cardcontent{\emph{Reiche Ernte und	Die	Intrige}
	\\
	\smallskip
	\\
	\emph{Wer zuletzt lacht:}
	\\
	Bauerndorf, Ernte, Harlekin, Pferdehändler, Treibjagd / Adlige, Handlanger, Lakai, Trickser, Verwalter
	\\
	\smallskip
	\\
	\emph{Würze des Lebens:}	
	\\
	Festplatz, Füllhorn, Junge Hexe (+ Wunschbrunnen als Bannstapel), Nachbau, Turnier (+ Preiskarten) / 	Bergwerk, Burghof, Große Halle, Kupferschmied, Tribut 
	\\
	\smallskip
	\\
	\emph{Kleine  Siege:}
	\\
	Nachbau, Treibjagd, Turnier (+ Preiskarten), Wahrsagerin, Weiler / Große Halle, Handlanger, Harem, Herzog, Verschwörer
	\\
	}
\end{tikzpicture}
\hspace{1cm}
	    % Basic settings for this card set
\renewcommand{\cardcolor}{hinterlands}
\renewcommand{\cardextension}{Erweiterung V}
\renewcommand{\cardextensiontitle}{Hinterland}

\clearpage
\newpage
\section{\cardextension \ - \cardextensiontitle}

\begin{tikzpicture}
	\card
	\cardstrip
	\cardbanner{banner/white.png}
	\cardicon{banner/coin.png}
	\cardprice{3}
	\cardtitle{Aufbau}
	\cardcontent{Zuerst musst du eine Karte aus deiner Hand entsorgen. Den gerade ausgespielten Ausbau kannst du nicht entsorgen, da du ihn nicht mehr auf der Hand hast. Hast du keine Karte mehr auf der Hand, kannst du keine Karte entsorgen und darfst dir auch keine Karten nehmen. Wenn du eine Karte entsorgt hast nimmst du dir 2 Karten, eine Karte, die genau \coin{1} mehr kostet, als die entsorgte Karte und eine Karte, die genau \coin{1} weniger kostet, als die entsorgte Karte. Du nimmst beide Karten aus dem Vorrat. Kannst du eine der beiden Karten nicht nehmen, weil keine Karte mit genau den geforderten Kosten im Vorrat ist, nimmst du dir trotzdem die andere Karte (wenn eine im Vorrat ist). Du legst die beiden Karten sofort auf deinen Nachziehstapel, anstatt auf deinen Ablagestapel. Du darfst die Karten in beliebiger Reihenfolge nehmen. Entsorgst du z. B. eine Herzogin (\coin{2}), so nimmst du dir eine Karte, die genau \coin{3} kostet (z. B. die Oase) und eine Karte, die genau \coin{1} kostet. Da es (üblicherweise) keine Karte gibt, die \coin{1} kostet, nimmst du dir keine weitere Karte. Entsorgst du z. B. ein Kupfer, so müsstest du dir eine Karte mit den Kosten \coin{1} und eine Karte mit den Kosten \coin{-1} nehmen, da es (üblicherweise) beide Karten nicht gibt, nimmst du dir keine Karte.}
\end{tikzpicture}
\hspace{-1cm}
\begin{tikzpicture}
\card
\cardstrip
\cardbanner{banner/gold.png}
\cardicon{banner/coin.png}
\cardprice{5}
\cardtitle{Blutzoll}
\cardcontent{Diese Karte ist eine Geldkarte mit dem Wert \coin{1}, wie ein Kupfer. Wenn du den Blutzoll ausspielst, darfst du dir ein Kupfer direkt auf die Hand nehmen. Du kannst dieses Kupfer also noch in diesem Zug ausspielen. Ist kein Kupfer mehr im Vorrat, so nimmst du dir kein Kupfer. Wenn du den Blutzoll nimmst oder kaufst, muss sich, beginnend mit dem Spieler links von dir, reihum jeder Mitspieler einen Fluch nehmen. Sind nicht mehr genügend Flüche im Vorrat, werden nur so viele verteilt wie möglich. Blutzoll ist keine Angriffskarte und kann z. B. nicht mit dem Burggraben abgewehrt werden.}
\end{tikzpicture}
\hspace{-1cm}
\begin{tikzpicture}
\card
\cardstrip
\cardbanner{banner/white.png}
\cardicon{banner/coin.png}
\cardprice{5}
\cardtitle{Botschaft}
\cardcontent{Wenn du die Botschaft ausspielst, ziehst du zunächst 5 Karten auf deine Hand nach. Wenn du, auch nach dem Mischen deines Ablagestapels, nur weniger als 5 Karten nachziehen kannst, ziehst du nur soviele Karten nach wie möglich. Dann legst du 3 Karten aus deiner Hand ab. Du kannst also auch Karten ablegen, die du gerade nachgezogen hast. Wenn du die Botschaft nimmst oder kaufst, muss sich, beginnend mit dem Spieler links von dir, reihum jeder Spieler ein Silber nehmen. Ist kein Silber mehr im Vorrat, nimmt der Spieler kein Silber.}
\end{tikzpicture}
\hspace{-1cm}
\begin{tikzpicture}
\card
\cardstrip
\cardbanner{banner/white.png}
\cardicon{banner/coin.png}
\cardprice{4}
\cardtitle{Edler Räuber}
\cardcontent{Wenn du diese Karte ausspielst, erhältst du +\coin{1}. Wenn du diese Karte ausspielst \emph{und auch} wenn du diese Karte kaufst, muss beginnend bei dem Spieler links von dir, reihum jeder Mitspieler die obersten beiden Karten von seinem Nachziehstapel aufdecken. Hat der Spieler 1 Silber oder 1 Gold aufgedeckt, so muss er dieses entsorgen. Hat der Spieler ein Silber und ein Gold (oder auch 2 Silber bzw. 2 Gold) aufgedeckt, entscheidest du welches davon er entsorgen muss. Die andere aufgedeckte Karte legt er jeweils auf seinen Ablagestapel. Hat der Spieler weder Silber noch Gold aufgedeckt, so legt er beide Karten ab. Hat der Spieler überhaupt keine Geldkarte aufgedeckt, so legt er beide Karten ab \emph{und} nimmt sich ein Kupfer. Du musst alle entsorgten Silber und Gold nehmen. Wenn du den Edlen Räuber ausspielst, können deine Mitspieler mit Reaktionskarten, wie z. B. Burggraben (Dominion – Basisspiel), auf den Angriff reagieren. Wenn du den Edlen Räuber jedoch kaufst, dürfen deine Mitspieler keine Reaktionskarten aufdecken, da du den Edlen Räuber nicht ausgespielt hast.}
\end{tikzpicture}
\hspace{-1cm}
\begin{tikzpicture}
\card
\cardstrip
\cardbanner{banner/blue.png}
\cardicon{banner/coin.png}
\cardprice{4}
\cardtitle{\scriptsize{Fahrender Händler}}
\cardcontent{\begin{spacing}{0.1cm}\tiny{Wenn du diese Karte ausspielst, musst du eine Karte aus deiner Hand entsorgen. Wenn du das machst, nimmst du dir so viele Silber vom Vorrat, wie die Geld-Kosten der entsorgten Karte. Entsorgst du z. B. eine Karte, die 3 Geld kostet, so nimmst du dir 3 Silber, entsorgst du z. B. ein Kupfer (0 Geld), nimmst du dir kein Silber. Sind im Vorrat nicht mehr genügend Silber, nimmst du dir nur so viele wie möglich. Kannst du keine Karte aus der Hand entsorgen, nimmst du dir auch kein Silber. Sind die Kosten der Karten verändert (z. B. durch die Fernstraße), nimmst du dir nur soviele Silber, wie die entsorgte Karte in diesem Moment kostet. Entsorgst du z. B. ein Anwesen, während die Fernstraße im Spiel ist, nimmst du dir nur 1 Silber. Andere Kosten außer Geld, z. B. Tränke aus Dominion – Die Alchemisten, haben keine Auswirkungen darauf, wieviele Silber du dir nimmst.\\
Der Fahrende Händler ist auch eine Reaktionskarte. Immer wenn du eine Karte nimmst, kannst du die Karte aus deiner Hand aufdecken. Du kannst den Fahrenden Händler aufdecken, wenn du von einem Angriff betroffen bist oder auch, wenn du in deinem eigenen Zug freiwillig oder unfreiwillig eine Karte nimmst oder kaufst. Wenn du das machst, nimmst du dir anstatt der Karte, die du dir eigentlich nehmen müsstest, ein Silber. Die andere Karte nimmst du nicht. Spielt z. B. ein Mitspieler die Hexe (Dominion – Basisspiel) aus, so kannst du den Fahrenden Händler aus deiner Hand aufdecken. Wenn du das machst, nimmst du dir keinen Fluch, sondern stattdessen ein Silber. Wenn du dir auf diese Art ein Silber nimmst, legst du es immer auf deinen Ablagestapel, auch wenn du die andere Karte z. B. auf die Hand genommen oder auf deinen Nachziehstapel gelegt hättest. Wenn du eine Karte kaufst, dir jedoch stattdessen ein Silber nimmst, musst du die ursprünglich gekaufte Karte trotzdem bezahlen. Wenn die Karte, die du ursprünglich nehmen müsstest, beim Nehmen eine Anweisung auslöst („Wenn du diese Karte nimmst ...“), wird diese nicht ausgelöst, weil du die Karte nicht nimmst. Wenn du z. B. einen Blutzoll kaufst, kannst du den Fahrenden Händler aus deiner Hand aufdecken und statt dem Blutzoll ein Silber nehmen. Kein Mitspieler muss sich einen Fluch nehmen, weil du den Blutzoll nicht genommen hast. Wenn die Karte, die du ursprünglich nehmen müsstest, beim Kaufen eine Anweisung auslöst („Wenn du diese Karte kaufst ...“), wird diese ausgelöst, weil der Kauf durchgeführt ist, obwohl du die Karte nicht nimmst. Wenn du z. B. ein Fruchtbares Land kaufst, kannst du den Fahrenden Händler aus deiner Hand aufdecken und statt dem Fruchtbaren Land ein Silber nehmen. Die zusätzliche Anweisung auf dem Fruchtbaren Land wird ausgelöst, weil du die Karte gekauft hast (selbst wenn du sie nicht wie üblich genommen hast). Du entsorgst also eine Karte aus deiner Hand und nimmst dir eine andere Karte, die 2 Geld kostet, genauso als hättest du das Fruchtbare Land nach deinem Kauf auch wirklich genommen.\\}
	\end{spacing}
	}
\end{tikzpicture}
\hspace{-1cm}
\begin{tikzpicture}
\card
\cardstrip
\cardbanner{banner/white.png}
\cardicon{banner/coin.png}
\cardprice{5}
\cardtitle{Feilscher}
\cardcontent{Wenn du diese Karte ausspielst erhältst du +\coin{2}. Wenn diese Karte im Spiel ist und du eine Karte kaufst, musst du dir immer eine zusätzliche Karte nehmen, die billiger ist als die gerade gekaufte und die keine Punktekarte ist. Die zusätzliche Karte nimmst du aus dem Vorrat und legst sie auf deinen Ablagestapel. Hast du einen Feilscher im Spiel und kaufst eine Provinz, so kannst du dir z. B. ein Gold nehmen. Du nimmst dir nur eine zusätzliche Karte, wenn du eine Karte kaufst. Wenn du eine Karte auf eine andere Art nimmst, nimmst du dir keine weitere Karte. Ist im Vorrat keine billigere Karte, als die gerade gekaufte, so nimmst du dir keine zusätzliche Karte. Diese Anweisung gilt nur solange der Feilscher im Spiel ist. Spielst du den Feilscher z. B. auf einen Thronsaal, spielst du die Karte zwar zweimal aus, die Karte ist jedoch nur einmal im Spiel. Du darfst dir in diesem Fall also nicht 2 zusätzliche Karten nehmen. Hast du 2 Feilscher im Spiel, so darfst (und musst) du dir auch 2 zusätzliche Karten nehmen. Kombinierte Punktekarten sind auch Punktekarten und dürfen nicht mit dem Feilscher genommen werden.}
\end{tikzpicture}
\hspace{-1cm}
\begin{tikzpicture}
\card
\cardstrip
\cardbanner{banner/white.png}
\cardicon{banner/coin.png}
\cardprice{5}
\cardtitle{Fernstrasse}
\cardcontent{Wenn du diese Karte ausspielst, ziehst du zuerst eine Karte nach und erhältst +1 Aktion. Solange diese Karte im Spiel ist, kosten alle Karten \coin{1} weniger (niemals jedoch weniger als \coin{0}). Das gilt für alle Karten im Vorrat, auf den Händen der Spieler und in deren Kartenstapeln. Spielst du z. B. die Fernstraße, danach den Aufbau und entsorgst ein Kupfer aus deiner Hand, so kannst du dir z. B. ein Anwesen nehmen. Das Anwesen kostet nur noch 1 Geld, während das Kupfer immer noch 0 Geld kostet. Diese Anweisung gilt nur, solange die Fernstraße im Spiel ist. Spielst du die Fernstraße z. B. auf einen Thronsaal, spielst du die Karte zwar zweimal aus, die Karte ist jedoch nur einmal im Spiel. In diesem Fall kosten also alle Karten immer noch nur \coin{1} weniger. Hast du 2 Fernstraßen im Spiel, so kosten alle Karten \coin{2} weniger (immer noch bis zu einem Minimum von \coin{0}).}
\end{tikzpicture}
\hspace{-1cm}
\begin{tikzpicture}
\card
\cardstrip
\cardbanner{banner/green.png}
\cardicon{banner/coin.png}
\cardprice{6}
\cardtitle{\scriptsize{Fruchtbares Land}}
\cardcontent{Diese Königreichkarte ist eine Punktekarte, keine Aktionskarte. Sie hat bis zum Ende des Spiels keine Funktion. Bei der Wertung zählt sie 2 Siegpunkte. Wenn du das Fruchtbare Land kaufst (nicht, wenn du die Karte auf eine andere Art nimmst), musst du eine Karte aus deiner Hand entsorgen und dir eine Karte nehmen, die genau \coin{2} mehr kostet als die entsorgte Karte. Wenn im Vorrat keine Karte ist, die genau \coin {2} mehr kostet als die entsorgte Karte, oder du keine Karte entsorgen kannst, nimmst du dir keine Karte (außer dem gerade gekauften Fruchtbaren Land). \\ \smallskip
	Im Spiel zu 3. und 4. werden 12 Karten Fruchtbares Land verwendet, im Spiel zu 2. werden 8 Karten Fruchtbares Land verwendet.}
\end{tikzpicture}
\hspace{-1cm}
\begin{tikzpicture}
\card
\cardstrip
\cardbanner{banner/white.png}
\cardicon{banner/coin.png}
\cardprice{5}
\cardtitle{Gasthaus}
\cardcontent{Wenn du diese Karte ausspielst, ziehst du zuerst 2 Karten nach und erhältst +2 Aktionen. Dann legst du 2 Karten aus deiner Hand ab. Du kannst auch Karten ablegen, die du gerade nachgezogen hast. \\ \smallskip 
	Wenn du diese Karte nimmst oder kaufst, darfst du dir sofort deinen Nachziehstapel durchsehen (was üblicherweise nicht erlaubt ist) und beliebig viele Aktionskarten daraus in deinen Nachziehstapel einmischen. Du darfst auch das gerade genommene Gasthaus selbst einmischen, da es, wenn die Anweisung ausgeführt wird, bereits auf dem Ablagestapel liegt. Kombinierte Aktionskarten sind auch Aktionskarten. Du musst die Aktionskarten, die du in deinen Nachziehstapel einmischst, deinen Mitspielern vorzeigen.}
\end{tikzpicture}
\hspace{-1cm}
\begin{tikzpicture}
\card
\cardstrip
\cardbanner{banner/white.png}
\cardicon{banner/coin.png}
\cardprice{4}
\cardtitle{\footnotesize{Gewürzhändler}}
\cardcontent{Du \emph{darfst} eine Karte aus deiner Hand entsorgen. Wenn du das machst, wählst du eine der beiden angegebenen Kombinationen. \emph{Entweder} ziehst du 2 Karten nach und erhältst +1 Aktion oder du erhältst +\coin{2} und +1 Kauf. Wenn du keine Karte aus deiner Hand entsorgen kannst oder willst, darfst du auch keine der angegebenen Anweisungen ausführen.}
\end{tikzpicture}
\hspace{-1cm}
\begin{tikzpicture}
\card
\cardstrip
\cardbanner{banner/white.png}
\cardicon{banner/coin.png}
\cardprice{6}
\cardtitle{Grenzdorf}
\cardcontent{Wenn du das Grenzdorf ausspielst ziehst du 1 Karte von deinem Nachziehstapel und erhältst +2 Aktionen.\\ \smallskip Wenn du das Grenzdorf nimmst oder kaufst, musst du dir zusätzlich eine weitere Karte nehmen, die weniger kostet als das Grenzdorf. Üblicherweise nimmst du dir eine Karte, die bis zu \coin{5} kostet. Sollte das Grenzdorf jedoch weniger als \coin{6} gekostet haben, z. B. weil die Karte Fernstraße im Spiel ist, darfst du dir nur eine Karte nehmen, die weniger kostet als das Grenzdorf momentan. Du nimmst dir die zusätzliche Karte nur, wenn du das Grenzdorf nimmst oder kaufst, nicht jedesmal, wenn du es ausspielst.}
\end{tikzpicture}
\hspace{-1cm}
\begin{tikzpicture}
\card
\cardstrip
\cardbanner{banner/white.png}
\cardicon{banner/coin.png}
\cardprice{2}
\cardtitle{Herzogin}
\cardcontent{Wenn du die Herzogin ausspielst, erhältst du zunächst +\coin{2}. Dann sieht sich, beginnend mit dem Spieler links von dir, reihum jeder Spieler (auch du selbst) die oberste Karte von seinem Nachziehstapel an und entscheidet selbst, ob er sie ablegt oder zurück auf seinen Nachziehstapel legt. Kann ein Spieler, auch nach dem Mischen seines Ablagestapels, keine Karte ansehen, so sieht er keine Karte an. 
	\\
	\smallskip
	\\
	Wenn die Herzogin als eine der 10 Königreichkarten verwendet wird, darf sich jeder Spieler, immer wenn er ein Herzogtum nimmt oder kauft, zusätzlich eine Karte Herzogin nehmen.}
\end{tikzpicture}
\hspace{-1cm}
\begin{tikzpicture}
\card
\cardstrip
\cardbanner{banner/white.png}
\cardicon{banner/coin.png}
\cardprice{5}
\cardtitle{Kartograph}
\cardcontent{Du ziehst zunächst 1 Karte von deinem Nachziehstapel und erhältst +1 Aktion. Dann siehst du dir die obersten 4 Karten von deinem Nachziehstapel an. (Wenn du auch nach dem Mischen deines Ablagestapels nur weniger als 4 Karten ansehen kannst, siehst du nur soviele Karten an wie möglich.) Lege beliebig viele (also 0-4) dieser Karten ab. Lege die übrigen dieser angesehenen Karten in beliebiger Reihenfolge zurück auf deinen Nachziehstapel. Du musst die Karten, die du zurück legst, deinen Mitspielern nicht zeigen.}
\end{tikzpicture}
\hspace{-1cm}
\begin{tikzpicture}
\card
\cardstrip
\cardbanner{banner/goldblue.png}
\cardicon{banner/coin.png}
\cardprice{2}
\cardtitle{Katzengold}
\cardcontent{Diese Karte ist gleichzeitig eine Geld- und eine Reaktionskarte. Katzengold kann, wie jede andere Geldkarte, in der Kaufphase ausgespielt werden. Wenn du in diesem Zug zum ersten Mal ein Katzengold ausspielst, ist es \coin{1} wert, wie ein Kupfer. Jedes weitere Katzengold, das du in diesem Zug ausspielst, ist \coin{4} wert. Spielst du z. B. 3 Karten Katzengold in deiner Kaufphase aus, so hast du \coin{9} (\coin{1} + \coin{4} + \coin{4}) für deinen Kauf bzw. deine Käufe zur Verfügung. Katzengold ist auch eine Reaktionskarte. Immer wenn ein anderer Spieler eine Provinz nimmt oder kauft, darfst du Katzengold aus deiner Hand entsorgen. Wenn du das machst, nimmst du dir ein Gold und legst es sofort auf deinen Nachziehstapel. Ist kein Gold mehr im Vorrat, erhältst du kein Gold, du kannst das Katzengold jedoch trotzdem entsorgen.}
\end{tikzpicture}
\hspace{-1cm}
\begin{tikzpicture}
\card
\cardstrip
\cardbanner{banner/white.png}
\cardicon{banner/coin.png}
\cardprice{3}
\cardtitle{Komplott}
\cardcontent{\emph{Errata:} Der Kartentext ist falsch, es sollte „Zu Beginn deiner Aufräumphase darfst du eine Aktionskarte wählen, die du im Spiel hast. Wenn du die gewählte Karte in diesem Zug ablegst, lege sie auf deinen Nachziehstapel.“ statt „Zu Beginn deiner Aufräumphase darfst du eine deiner ausgespielten Aktionskarten wählen. Wenn du die gewählte Karte in dieser Aufräumphase ablegen würdest, darfst du sie stattdessen auf deinen Nachziehstapel legen.“ heißen. 
	\\
	\smallskip
	\\
	Wenn du diese Karte ausspielst, ziehst du zuerst 1 Karte nach und erhältst +1 Aktion. Zu Beginn deiner Aufräumphase darfst du eine Aktionskarte wählen, die du im Spiel hast. Du kannst das auch gerade ausgespielte Komplott selbst wählen. Statt die gewählte Aktionskarte in dieser Aufräumphase auf deinen Ablagestapel zu legen, darfst du sie auf deinen Nachziehstapel legen. Das machst du, bevor du Karten für die nächste Runde nachziehst. Wenn du die gewählte Aktionskarte nicht in dieser Aufräumphase ablegst (z. B. eine Dauerkarte aus Dominion – Seaside), kannst du sie auch nicht auf deinen Nachziehstapel legen.
	\\
	\smallskip
	\\	
	\emph{Anmerkung:} Das „Ablegen“ wird dennoch ausgelöst, bevor die Karte auf den Nachziehstapel gelegt wird.}
\end{tikzpicture}
\hspace{-1cm}
\begin{tikzpicture}
\card
\cardstrip
\cardbanner{banner/white.png}
\cardicon{banner/coin.png}
\cardprice{4}
\cardtitle{\footnotesize{Lebenskünstler}}
\cardcontent{Diese Karte hat 4 verschiedene Anweisungen, die du nacheinander ausführen musst (nur die letzte ist optional). Zuerst nimmst du dir ein Silber vom Vorrat und legst es auf deinen Ablagestapel. Ist kein Silber mehr im Vorrat, nimmst du dir kein Silber. Danach siehst du dir die oberste Karte deines Nachziehstapels an und entscheidest, ob du sie zurück auf deinen Nachziehstapel legst oder ablegst. Wenn dein Nachzieh- und dein Ablagestapel leer sind, siehst du dir keine Karte an. Dann ziehst du solange Karten nach, bis du 5 Karten auf der Hand hast. Hast du bereits 5 oder mehr Karten auf der Hand, ziehst du keine Karten nach. Kannst du, auch nach dem Mischen deines Ablagestapels, nicht genügend Karten nachziehen, so ziehst du nur soviele Karten nach wie möglich. Zuletzt darfst du eine Karte aus deiner Hand entsorgen, die keine Geldkarte ist. Kombinierte Geldkarten sind auch Geldkarten.}
\end{tikzpicture}
\hspace{-1cm}
\begin{tikzpicture}
\card
\cardstrip
\cardbanner{banner/white.png}
\cardicon{banner/coin.png}
\cardprice{5}
\cardtitle{Mandarin}
\cardcontent{Wenn du diese Karte ausspielst erhältst du zuerst +\coin{3}. Dann musst du eine Karte aus deiner Hand auf deinen Nachziehstapel legen. Wenn du keine Karte mehr auf der Hand hast, legst du keine Karte auf deinen Nachziehstapel.\\ \smallskip
	Wenn du diese Karte nimmst oder kaufst, musst du alle Geldkarten, die du im Spiel hast, in beliebiger Reihenfolge zurück auf deinen Nachziehstapel legen. Du musst deinen Mitspielern nicht zeigen, in welcher Reihenfolge du die Karten zurück legst. Du musst alle Geldkarten, die du im Spiel hast, zurück legen. Geldkarten auf deiner Hand sind nicht im Spiel. Du legst sie also nicht auf den Nachziehstapel. Du hast jedoch alle Münzen, der ausgespielten Geldkarten in dieser Kaufphase zur Verfügung. Hast du z. B. +1 Kauf und 4 Gold ausgespielt und kaufst den Mandarin, so legst du die 4 Gold zurück auf deinen Nachziehstapel und hast trotzdem noch \coin{7} für den zweiten Kauf zur Verfügung. Du legst deine ausgespielten Geldkarten auch auf den Nachziehstapel zurück, wenn du den Mandarin nimmst, wobei du üblicherweise nur in der Kaufphase Geldkarten im Spiel hast.}
\end{tikzpicture}
\hspace{-1cm}
\begin{tikzpicture}
\card
\cardstrip
\cardbanner{banner/white.png}
\cardicon{banner/coin.png}
\cardprice{5}
\cardtitle{Markgraf}
\cardcontent{Zuerst ziehst du 3 Karten nach und erhältst +1 Kauf. Dann muss, beginnend mit dem Spieler links von dir, reihum jeder Mitspieler eine Karte nachziehen und dann solange Karten aus seiner Hand ablegen, bis er nur noch 3 Karten auf der Hand hat. Hat ein Spieler, auch nachdem er eine Karte nachgezogen hat, nur 3 oder weniger Karten auf der Hand, so legt er keine Karte ab.}
\end{tikzpicture}
\hspace{-1cm}
\begin{tikzpicture}
\card
\cardstrip
\cardbanner{banner/white.png}
\cardicon{banner/coin.png}
\cardprice{4}
\cardtitle{Nomadencamp}
\cardcontent{Wenn du diese Karte ausspielst, erhältst du +1 Kauf und +\coin{2}. Wenn du das Nomadencamp nimmst, musst du es sofort auf deinen Nachziehstapel legen, statt wie üblich auf den Ablagestapel.}
\end{tikzpicture}
\hspace{-1cm}
\begin{tikzpicture}
\card
\cardstrip
\cardbanner{banner/white.png}
\cardicon{banner/coin.png}
\cardprice{3}
\cardtitle{Oase}
\cardcontent{Du ziehst zuerst 1 Karte nach und erhältst +1 Aktion und +\coin{1}. Dann legst du eine Karte aus deiner Hand ab. Du darfst auch die gerade nachgezogene Karte ablegen. Wenn du keine Karte nachziehen kannst, weil dein Nachzieh- und Ablagestapel leer sind, musst du trotzdem eine Karte aus deiner Hand ablegen.}
\end{tikzpicture}
\hspace{-1cm}
\begin{tikzpicture}
\card
\cardstrip
\cardbanner{banner/white.png}
\cardicon{banner/coin.png}
\cardprice{3}
\cardtitle{Orakel}
\cardcontent{Reihum, beginnend mit dem Spieler links von dir, muss jeder Spieler (auch du selbst) die obersten beiden Karten von seinem Nachziehstapel aufdecken. Du entscheidest, ob er beide Karten ablegen oder zurück auf seinen Nachziehstapel legen muss. Der Spieler entscheidet selbst, in welcher Reihenfolge er die Karten auf seinen Nachziehstapel legt. Dann erst ziehst du 2 Karten nach. Hast du also deine eigenen Karten zurück auf deinen Nachziehstapel gelegt, so ziehst du diese Karten nach.}
\end{tikzpicture}
\hspace{-1cm}
\begin{tikzpicture}
\card
\cardstrip
\cardbanner{banner/gold.png}
\cardicon{banner/coin.png}
\cardprice{5}
\cardtitle{Schatztruhe}
\cardcontent{Diese Geldkarte ist \coin{3} wert, wie ein Gold. \\ \smallskip
	Wenn du die Schatztruhe nimmst oder kaufst, musst du dir zusätzlich 2 Kupfer nehmen. Sind nur weniger als 2 Kupfer im Vorrat, so nimmst du dir nur so viele wie möglich. Du nimmst dir die 2 Kupfer nur, wenn du die Schatztruhe nimmst oder kaufst, nicht jedesmal, wenn du sie ausspielst.}
\end{tikzpicture}
\hspace{-1cm}
\begin{tikzpicture}
\card
\cardstrip
\cardbanner{banner/green.png}
\cardicon{banner/coin.png}
\cardprice{4}
\cardtitle{Seidenstrasse}
\cardcontent{Diese Königreichkarte ist eine Punktekarte, keine Aktionskarte. Sie hat bis zum Ende des Spiels keine Funktion. Bei der Wertung zählt sie 1 Siegpunkt für je volle 4 Punktekarten in deinem gesamten Kartensatz (Nachziehstapel, Ablagestapel und Handkarten). Bei Spielende suchst du aus deinem gesamten Kartensatz alle Punktekarten heraus und zählst sie. (Die Seidenstraße selbst ist auch eine Punktekarte und zählt zur Gesamtzahl der Punktekarten.) Die Ergebnis durch 4 geteilt und abgerundet ergibt die Punkte.\\ \smallskip
	 Im Spiel zu 3. und 4. werden 12 Karten Seidenstraße verwendet, im Spiel zu 2. werden 8 Karten Seidenstraße verwendet.}
\end{tikzpicture}
\hspace{-1cm}
\begin{tikzpicture}
\card
\cardstrip
\cardbanner{banner/white.png}
\cardicon{banner/coin.png}
\cardprice{5}
\cardtitle{Stallungen}
\cardcontent{Du darfst eine Karte aus deiner Hand ablegen. Wenn du das machst, ziehst du 3 Karten nach und erhältst +1 Aktion. Wenn du keine Karte aus deiner Hand ablegen kannst oder willst, darfst du auch keine der angegebenen Anweisungen ausführen.}
\end{tikzpicture}
\hspace{-1cm}
\begin{tikzpicture}
\card
\cardstrip
\cardbanner{banner/greenblue.png}
\cardicon{banner/coin.png}
\cardprice{3}
\cardtitle{Tunnel}
\cardcontent{\begin{singlespace}\tiny{
	Diese Karte ist gleichzeitig eine Punktekarte und eine Reaktionskarte. Bei der Wertung zählt sie 2 Siegpunkte. 
	\\
	\smallskip
	\\
	Die Reaktion kann eingesetzt werden, wenn du den Tunnel außerhalb einer Aufräumphase ablegen musst. Du darfst dann den Tunnel aufdecken bevor du ihn ablegst und nimmst dir ein Gold. Du darfst den Tunnel nicht „freiwillig“ ablegen, eine Kartenanweisung muss dich dazu zwingen. Dies kann in deinem eigenen Zug (z. B. durch die Oase) geschehen, oder im Zug eines Mitspielers (z. B. durch den Markgrafen). Die Reaktion setzt auch ein, wenn du den Tunnel nicht wie üblich aus deiner Hand ablegst, sondern von deinem Nachziehstapel oder von aufgedeckten Karten ablegst (z. B. durch das Orakel). Wenn der Tunnel normalerweise nicht aufgedeckt würde (z. B. durch den Kartographen), musst du ihn aufdecken, um dir ein Gold zu nehmen. Du musst den Tunnel nicht aufdecken (und dementsprechend auch kein Gold nehmen), wenn du ihn ablegst. Die Reaktion wird nicht ausgelöst, wenn du die Karte auf deinen Ablagestapel legst, ohne sie aufzudecken. Du kannst die Karte nicht aufdecken, wenn du sie gerade genommen oder gekauft hast. Du kannst den Tunnel auch nicht aufdecken, wenn du z. B. durch den Kanzler (Dominion – Basisspiel) deinen Nachziehstapel ablegst oder wenn durch die Besessenheit (Dominion – Die Alchemisten) entsorgte Karten auf deinen Ablagestapel gelegt werden. Die Reaktion setzt auch nicht ein, wenn du den Tunnel regulär in der Aufräumphase ablegst. Es muss eine Anweisung geben, die dich zwingt, eine Karte abzulegen. 
	\\
	\smallskip
	\\	
	Das Gold nimmst du dir aus dem Vorrat. Ist dort kein Gold mehr, nimmst du dir keines.}
	\end{singlespace}
	}
\end{tikzpicture}
\hspace{-1cm}
\begin{tikzpicture}
	\card
	\cardstrip
	\cardbanner{banner/white.png}
	\cardicon{banner/coin.png}
	\cardprice{2}
	\cardtitle{Wegkreuzung}
	\cardcontent{Decke deine gesamte Kartenhand auf. Dann ziehe soviele Karten nach, wie du Punktekarten auf deiner Hand hast. Kombinierte Punktekarten sind auch Punktekarten. Hast du keine Punktekarten auf der Hand, ziehst du auch keine Karten nach. Wenn du in diesem Zug zum ersten Mal eine Wegkreuzung ausspielst, erhältst du +3 Aktionen. Für alle Wegkreuzungen, die du danach in diesem Zug ausspielst, erhältst nur noch +1 Aktion. Spielst du die Karte z. B. durch den Thronsaal (Dominion – Basisspiel) zweimal aus (und hast in diesem Zug noch keine Wegkreuzung ausgespielt), so erhältst du beim ersten Mal +3 Aktionen, beim zweiten Ausspielen nochmals +1 Aktion, also insgesamt +4 Aktionen.}
\end{tikzpicture}
\hspace{-1cm}
\begin{tikzpicture}
	\card
	\cardstrip
	\cardbanner{banner/white.png}
	\cardtitle{\scriptsize{Empfohlene 10er Sätze\qquad}}
	\cardcontent{\emph{Hinterland:}\\ \smallskip
	\emph{Einführung:} \\ 
	Aufbau, Feilscher, Gewürzhändler, Lebenskünstler, Markgraf, Nomadencamp, Oase, Schatztruhe, Stallungen, Wegkreuzung \\ 
	\smallskip \emph{Lauterer Wettbewerb:} \\ 
	Aufbau, Blutzoll, Edler Räuber, Fahrender Händler, Fruchtbares Land, Grenzdorf, Herzogin, Kartograph, Seidenstraße, Stallungen \\
	\smallskip \emph{Gelegenheiten:} \\ 
	Fahrender Händler, Feilscher, Fernstraße, Gewürzhändler, Grenzdorf, Herzogin, Katzengold, Komplott, Nomadencamp, Schatztruhe \\
	\smallskip \emph{Eröffnungen:} \\ 
	Botschaft, Gasthaus, Kartograph, Lebenskünstler, Mandarin, Nomadencamp, Oase, Orakel, Tunnel, Wegkreuzung \\
}
\end{tikzpicture}
\hspace{-1cm}
\begin{tikzpicture}
	\card
	\cardstrip
	\cardbanner{banner/white.png}
	\cardtitle{\scriptsize{Empfohlene 10er Sätze\qquad}}
	\cardcontent{\emph{Hinterland und Basisspiel:}\\ \smallskip
	\emph{Straßenräuber:} \\ 
	Bibliothek, Geldverleiher, Keller, Thronsaal, Werkstatt / Edler Räuber, Fernstraße, Gasthaus, Markgraf, Oase \\
	\smallskip \emph{Abenteuerfahrt:} \\ 
	Abenteurer, Jahrmarkt, Kanzler, Laboratorium, Umbau/ Fruchtbares Land, Gewürzhändler, Katzengold, Orakel, Wegkreuzung
	\smallskip \emph{Hinterland und die Intrige:} \\ 
	\smallskip \emph{Geld aus Nichts:} \\ 
	Armenviertel, Große Halle, Handlanger, Kerkermeister, Kupferschmied / Kartograph, Lebenskünstler, Schatztruhe, Seidenstraße, Tunnel \\
	\smallskip \emph{Am Hof des Herzogs:} \\ 
	Anbau, Harem, Herzog, Maskerade, Verschwörer / Edler Räuber, Feilscher, Gasthaus, Herzogin, Komplott \\
}
\end{tikzpicture}
\hspace{-1cm}
\begin{tikzpicture}
	\card
	\cardstrip
	\cardbanner{banner/white.png}
	\cardtitle{\scriptsize{Empfohlene 10er Sätze\qquad}}
	\cardcontent{\emph{Hinterland und Seaside:}\\ \smallskip
	\emph{Reisende:} \\ 
	Ausguck, Beutelschneider, Handelsschiff, Insel, Lagerhaus / Fruchtbares Land, Kartograph, Seidenstraße, Stallungen, Wegkreuzung \\
	\smallskip \emph{Diplomatie:} \\ 
	Bazar, Botschafter, Embargo, Karawane, Schmuggler / Blutzoll, Botschaft, Edler Räuber, Fahrender Händler, Fruchtbares Land \\
	\smallskip \emph{Hinterland und die Alchemisten:} \\ 
	\smallskip \emph{Träume sind Schäume:} \\ 
	Apotheker, Kräuterkundiger, Lehrling, Stein der Wei- sen, Verwandlung / Blutzoll, Herzogin, Katzengold, Komplott, Lebenskünstler \\
	\smallskip \emph{Weinviertel:} \\ 
	Golem, Lehrling, Universität, Vertrauter, Weinberg / Feilscher, Fernstraße, Fruchtbares Land, Nomadencamp, Wegkreuzung \\
}
\end{tikzpicture}
\hspace{-1cm}
\begin{tikzpicture}
	\card
	\cardstrip
	\cardbanner{banner/white.png}
	\cardtitle{\scriptsize{Empfohlene 10er Sätze\qquad}}
	\cardcontent{\emph{Hinterland und Blütezeit:}\\ \smallskip
	\emph{Karriereleiter:} \\ 
	Ausbau, Bischof, Hort, Münzer, Wachturm / Blutzoll, Edler Räuber, Fahrender Händler, Feilscher, Fruchtbares Land \\
	\smallskip \emph{Schatzfund:} \\ 
	Abenteuer, Bank, Denkmal, Handelsroute, Königliches Siegel/ Aufbau, Blutzoll, Katzengold, Mandarin, Schatztruhe \\
	\smallskip \emph{Hinterland und Reiche Ernte:} \\ 
	\smallskip \emph{Schmalhans:} \\ 
	Füllhorn, Harlekin, Pferdehändler, Turnier, Weiler / Edler Räuber, Fahrender Händler, Katzengold, Mandarin, Tunnel \\
	\smallskip \emph{Wanderzirkus:} \\ 
	Bauerndorf, Festplatz, Harlekin, Menagerie, Treibjagd / Botschaft, Grenzdorf, Katzengold, Nomadencamp, Oase \\
}
\end{tikzpicture}

\hspace{1cm}

	    % Basic settings for this card set
\renewcommand{\cardcolor}{darkages}
\renewcommand{\cardextension}{Erweiterung VI}
\renewcommand{\cardextensiontitle}{Dark Ages}

\clearpage
\newpage
\section{\cardextension \ - \cardextensiontitle}

\begin{tikzpicture}
	\card
	\cardstrip
	\cardbanner{banner/white.png}
	\cardicon{banner/coin.png}
	\cardprice{6}
	\cardtitle{Altar}
	\cardcontent{Du musst eine Karte aus deiner Hand entsorgen, falls möglich. Dann nimmst du dir eine Karte, die bis zu \coin{5} kostet, aus dem Vorrat (auch wenn du keine Karte entsorgen konntest). Die neue Karte legst du auf deinen Ablagestapel.}
\end{tikzpicture}
\hspace{-1cm}
\begin{tikzpicture}
\card
\cardstrip
\cardbanner{banner/white.png}
\cardicon{banner/coin.png}
\cardprice{1}
\cardtitle{Armenhaus}
\cardcontent{Wenn du diese Karte ausspielst, erhältst du zunächst +\coin{4} für die folgende Kaufphase. Dann musst du alle deine Handkarten aufdecken. Für jede Geldkarte auf deiner Hand reduzieren sich die +\coin{4} um \coin{1} (niemals jedoch unter \coin{0}). Kombinierte Geldkarten sind auch Geldkarten. Hast du z. B. 2 Kupfer auf der Hand, so bleiben dir noch +\coin{2}.}
\end{tikzpicture}
\hspace{-1cm}
\begin{tikzpicture}
\card
\cardstrip
\cardbanner{banner/white.png}
\cardicon{banner/coin.png}
\cardprice{5}
\cardtitle{Banditenlager}
\cardcontent{Du ziehst zuerst eine Karte nach, dann nimmst du dir eine Beute-Karte vom Beute-Stapel neben dem Vorrat und legst die Karte auf deinen Ablagestapel. Ist der Beute-Stapel leer, nimmst du dir keine Beute-Karte. Danach darfst du bis zu 2 weitere Aktionen ausführen.}
\end{tikzpicture}
\hspace{-1cm}
\begin{tikzpicture}
\card
\cardstrip
\cardbanner{banner/white.png}
\cardicon{banner/coin.png}
\cardprice{4}
\cardtitle{Barde}
\cardcontent{Du ziehst zuerst eine Karte nach. Dann deckst du die obersten 3 Karten von deinem Nachziehstapel auf. Kannst du (auch nach dem Mischen des Ablagestapels) keine 3 Karten aufdecken, deckst du nur soviele Karten auf, wie möglich. Lege alle aufgedeckten Aktionskarten in beliebiger Reihenfolge zurück auf deinen Nachziehstapel und lege die übrigen aufgedeckten Karten ab. Kombinierte Aktionskarten sind auch Aktionskarten. Hast du keine Aktionskarten aufgedeckt, legst du auch keine Karten zurück auf deinen Nachziehstapel. Danach darfst du bis zu 2 weitere Aktionen ausführen.}
\end{tikzpicture}
\hspace{-1cm}
\begin{tikzpicture}
\card
\cardstrip
\cardbanner{banner/blue.png}
\cardicon{banner/coin.png}
\cardprice{2}
\cardtitle{Bettler}
\cardcontent{Wenn du den Bettler in deiner Aktionsphase ausspielst, nimmst du dir 3 Kupfer aus dem Vorrat direkt auf deine Hand. Sind nicht mehr genügend Kupfer im Vorrat, nimmst du dir so viele, wie möglich. Wenn ein Mitspieler eine Angriffskarte ausspielt, darfst du den Bettler aus deiner Hand ablegen. (\emph{Achtung:} Der Bettler wird nicht wie üblich nur vorgezeigt und kommt dann zurück auf die Hand, sondern er wird abgelegt.) Wenn du das machst, nimmst du dir 2 Silber aus dem Vorrat. Eines davon legst du auf deinen Nachziehstapel, das andere legst du auf deinen Ablagestapel. Ist nur noch ein Silber im Vorrat so nimmst du nur dieses und legst es auf deinen Nachziehstapel. Ist kein Silber mehr im Vorrat, nimmst du dir keines.}
\end{tikzpicture}
\hspace{-1cm}
\begin{tikzpicture}
\card
\cardstrip
\cardbanner{banner/gold.png}
\cardicon{banner/coin.png}
\cardprice{0*}
\cardtitle{Beute}
\cardcontent{Wird im Spiel mindestens eine der Königreichkarten \emph{Banditenlager}, \emph{Marodeur} oder \emph{Raubzug} verwendet, so wird auch der Beute-Stapel benötigt. Der Beute-Stapel ist niemals im Vorrat und wird \emph{neben} diesen bereit gelegt. Die Karten vom Beute-Stapel können nur durch die Anweisungen auf den 3 oben genannten Karten genommen werden. Auf andere Weise können keine Karten vom Beute-Stapel gekauft oder genommen werden. Der Botschafter (Dominion – Seaside) darf keine Karten auf den Beute-Stapel zurück legen. Der Beute-Stapel wird für die Spielende-Bedingung \emph{nicht} beachtet.\\ \smallskip
	Die Beute ist eine Geldkarte mit Wert \coin{3}, wie Gold. Wenn du die Beute ausspielst, musst du sie sofort zurück auf den Beute-Stapel legen. Du musst Geldkarten, die du auf der Hand hast, nicht ausspielen.}
\end{tikzpicture}
\hspace{-1cm}
\begin{tikzpicture}
\card
\cardstrip
\cardbanner{banner/white.png}
\cardicon{banner/coin.png}
\cardprice{4}
\cardtitle{Eisenhändler}
\cardcontent{\emph{Errata:} Die Reihenfolge der Anweisungen ist falsch, es sollte „[...] Decke die oberste Karte von deinem Nachziehstapel auf. Du darfst sie ablegen. Unabhängig davon: Ist es eine...“ heißen.\\ \smallskip
	Du ziehst zuerst eine Karte nach, dann deckst du die oberste Karte von deinem Nachziehstapel auf. Dann entscheidest du dich, ob du die aufgedeckte Karte zurück auf den Nachziehstapel oder auf deinen Ablagestapel legst. Du erhältst einen Bonus entsprechend dem Kartentyp der aufgedeckten Karte. Ist die aufgedeckte Karte eine Karte mit kombinierten Kartentypen, erhältst du für jeden Kartentyp den angegebenen Bonus. Deckst du z. B. den Harem (Dominion – Die Intrige) auf, erhältst du +\coin{1} und +1 Karte. Danach darfst du eine weitere Aktion ausführen.}
\end{tikzpicture}
\hspace{-1cm}
\begin{tikzpicture}
\card
\cardstrip
\cardbanner{banner/white.png}
\cardicon{banner/coin.png}
\cardprice{3}
\cardtitle{Eremit}
\cardcontent{Wenn du den Eremiten ausspielst, siehst du dir zunächst deinen Ablagestapel durch. Dann darfst du eine Karte, die keine Geldkarte ist, aus dem Ablagestapel oder aus deiner Hand entsorgen. Du musst keine Karte entsorgen und du darfst keine Geldkarte entsorgen. Kombinierte Geldkarten, wie z. B. der Harem (Dominion – Die Intrige) sind Geldkarten. Egal ob du eine Karte entsorgt hast oder nicht, musst du dir eine Karte nehmen, die bis zu \coin{3} kostet. Du nimmst diese Karte aus dem Vorrat und legst sie auf deinen Ablagestapel. Du musst eine Karte nehmen, wenn möglich. Du darfst nicht darauf verzichten.\\ \smallskip
	Wenn du den Eremiten aus dem Spiel ablegst (normalerweise in der Aufräumphase am Ende der Runde, in der du die Karte ausgespielt hast) und du in diesem Zug keine Karte gekauft hast, entsorge den Eremiten und nimm dir einen Verrückten. Du nimmst den Verrückten vom Stapel neben dem Vorrat und legst ihn auf deinen Ablagestapel. Karten, die du auf andere Weise genommen hast, als sie zu kaufen, haben keinen Einfluss darauf, ob du den Eremiten entsorgst. Ist der Verrückten-Stapel leer, so nimmst du dir keinen. Wird der Eremit in der Aufräumphase nicht regulär abgelegt, sondern z. B. durch das Komplott (Dominion – Hinterland) zurück auf den Nachziehstapel gelegt, entsorgst du den Eremiten nicht und nimmst dir auch keinen Verrückten.}
\end{tikzpicture}
\hspace{-1cm}
\begin{tikzpicture}
\card
\cardstrip
\cardbanner{banner/gold.png}
\cardicon{banner/coin.png}
\cardprice{5}
\cardtitle{Falschgeld}
\cardcontent{Diese Geldkarte hat den Wert \coin{1}. Du legst das Falschgeld wie üblich in der Kaufphase aus. Du erhältst +1 Kauf, dann darfst du eine Geldkarte aus deiner Hand wählen und diese zweimal ausspielen. Du erhältst also zweimal den Wert der gewählten Geldkarte (zusätzlich zu den \coin{1} Geld vom Falschgeld selbst) und führst auch zusätzliche Anweisungen auf der Karte zweimal aus. Danach musst du die gewählte Karte entsorgen. Wenn du das Falschgeld nutzt, um die Beute zweimal auszuspielen, erhältst du also \coin{6} (zusätzlich zu den \coin{1} vom Falschgeld selbst) und legst die Beute zurück auf den Beute-Stapel. Du kannst die Beute nicht entsorgen, da sie durch die Anweisung bereits anderweitig entfernt wurde. Kombinierte Geldkarten sind auch Geldkarten und können mit Falschgeld gewählt werden.}
\end{tikzpicture}
\hspace{-1cm}
\begin{tikzpicture}
\card
\cardstrip
\cardbanner{banner/white.png}
\cardicon{banner/coin.png}
\cardprice{4}
\cardtitle{Festung}
\cardcontent{Wenn du diese Karte ausspielst, ziehst du zunächst eine Karte nach. Dann darfst du bis zu 2 weitere Aktionen ausführen. \\ \smallskip
	Wenn du die Festung entsorgst, nimmst du die Karte zurück auf die Hand. Es spielt dabei keine Rolle, ob du die Festung in deinem eigenen Zug oder im Zug eines Mitspielers entsorgst. Die Festung gilt als entsorgt, obwohl du sie sofort zurück auf die Hand nimmst. Spielst du z. B. den Leichenkarren aus und entscheidest dich dafür, die Festung zu entsorgen, gilt die Bedingung als erfüllt und du muss den Leichenkarren nicht entsorgen.}
\end{tikzpicture}
\hspace{-1cm}
\begin{tikzpicture}
\card
\cardstrip
\cardbanner{banner/white.png}
\cardicon{banner/coin.png}
\cardprice{3}
\cardtitle{Gassenjunge}
\cardcontent{Wenn du den Gassenjungen ausspielst, ziehst du zunächst eine Karte nach, dann muss jeder deiner Mitspieler seine Kartenhand auf 4 reduzieren. Mitspieler, die bereits 4 oder weniger Karten auf der Hand haben, legen keine Karten ab.\\ \smallskip
	Wenn der Gassenjunge im Spiel ist und du eine weitere Angriffskarte ausspielst, darfst du (bevor du die neue Karte ausführst) den Gassenjungen entsorgen und dir einen Söldner nehmen. Du nimmst den Söldner vom Stapel neben dem Vorrat und legst ihn auf deinen Ablagestapel. Ist der Söldner-Stapel leer, so nimmst du dir keinen. Wenn du den selben Gassenjungen zweimal ausspielst, z. B. durch die Prozession, darfst du ihn nicht entsorgen um dir einen Söldner zu nehmen. Wenn du zwei einzelne Gassenjungen spielst, darfst du den ersten entsorgen und dir einen Söldner nehmen.}
\end{tikzpicture}
\hspace{-1cm}
\begin{tikzpicture}
\card
\cardstrip
\cardbanner{banner/white.png}
\cardicon{banner/coin.png}
\cardprice{5}
\cardtitle{Grabräuber}
\cardcontent{Du musst eine der beiden Optionen wählen und diese ausführen, soweit möglich. Du darfst auch eine Option wählen, die du nicht ausführen kannst. Du darfst den Müll-Stapel jederzeit durchsehen. Wenn du dir eine Karte vom Müll-Stapel nimmst, zeigst du diese deinen Mitspielern und legst sie dann sofort auf deinen Nachziehstapel. Wenn keine Karte im Müll-Stapel ist, die \coin{3} - \coin{6} kostet, nimmst du dir keine Karte. Du darfst keine Karten mit Trank in den Kosten (Dominion – Die Alchemisten) nehmen.\\ \smallskip
	Wenn du dich dafür entscheidest, eine Aktionskarte zu entsorgen, nimmst du dir eine Karte aus dem Vorrat und legst sie auf deinen Ablagestapel.}
\end{tikzpicture}
\hspace{-1cm}
\begin{tikzpicture}
\card
\cardstrip
\cardbanner{banner/white.png}
\cardicon{banner/coin.png}
\cardprice{5}
\cardtitle{Graf}
\cardcontent{Diese Karte gibt dir 2 voneinander unabhängige Wahlmöglichkeiten. Du entscheidest dich zuerst, ob du 2 Handkarten ablegst oder eine Handkarte auf deinen Nachziehstapel legst oder dir ein Kupfer vom Vorrat nimmst und auf deinen Ablagestapel legst. Danach entscheidest du dich, ob du +\coin{3} möchtest oder alle deine verbliebenen Handkarten entsorgst oder dir ein Herzogtum vom Vorrat nimmst und auf deinen Ablagestapel legst. Du musst beide Wahlmöglichkeiten in dieser Reihenfolge ausführen. Du kannst z. B. zuerst 2 Handkarten ablegen und dir danach ein Herzogtum nehmen. Du darfst auch eine Option wählen, die du nicht ausführen kannst. Wenn du mehrere Karten entsorgst, die eine „Wenn du diese Karte entsorgst ...“-Anweisung haben, ensorgst du zuerst die Karten und führst die „Wenn du diese Karte entsorgst ...“-Anweisungen dann in beliebiger Reihenfolge aus.}
\end{tikzpicture}
\hspace{-1cm}
\begin{tikzpicture}
\card
\cardstrip
\cardbanner{banner/white.png}
\cardicon{banner/coin.png}
\cardprice{6}
\cardtitle{Jagdgründe}
\cardcontent{Wenn du diese Karte ausspielst, ziehst du 4 Karten nach. Wenn du die Jagdgründe entsorgst, musst du wählen und entweder 1 Herzogtum oder 3 Anwesen vom Vorrat nehmen und auf deinen Ablagestapel legen. Du darfst auch eine Option wählen, die du nicht oder nur teilweise ausführen kannst, z. B. weil nur noch 2 Anwesen im Vorrat sind.}
\end{tikzpicture}
\hspace{-1cm}
\begin{tikzpicture}
\card
\cardstrip
\cardbanner{banner/white.png}
\cardicon{banner/coin.png}
\cardprice{5}
\cardtitle{Katakomben}
\cardcontent{Wenn du diese Karte ausspielst, siehst du dir zuerst die obersten 3 Karten von deinem Nachziehstapel an. Dann musst du dich entscheiden, ob du entweder alle 3 Karten auf die Hand nimmst oder alle 3 Karten auf deinen Ablagestapel legst. Wenn du die 3 Karten ablegst (und nur dann), ziehst du die nächsten 3 Karten von deinem Nachziehstapel. Wenn du die Katakomben entsorgst, nimmst du dir eine Karte, die weniger kostet als die Katakomben selbst, vom Vorrat und legst sie auf deinen Ablagestapel. Dabei ist egal, ob du die Katakomben in deinem eigenen Zug oder im Zug eines Mitspielers entsorgst.}
\end{tikzpicture}
\hspace{-1cm}
\begin{tikzpicture}
\card
\cardstrip
\cardbanner{banner/white.png}
\cardicon{banner/coin.png}
\cardprice{2}
\cardtitle{Knappe}
\cardcontent{Wenn du diese Karte ausspielst erhältst du zuerst +\coin{1}. Dann wählst du eine der 3 Optionen: +2 Aktionen, +2 Kauf oder du nimmst dir ein Silber vom Vorrat und legst es auf deinen Ablagestapel. Wenn du den Knappen entsorgst, nimmst du dir eine beliebige Angriffskarte aus dem Vorrat und legst sie auf deinen Ablagestapel. Du darfst nur sichtbare Karten aus dem Vorrat nehmen, also z. B. nur den offen liegenden Ritter.}
\end{tikzpicture}
\hspace{-1cm}
\begin{tikzpicture}
\card
\cardstrip
\cardbanner{banner/white.png}
\cardicon{banner/coin.png}
\cardprice{5}
\cardtitle{Kultist}
\cardcontent{Wenn du diese Karte ausspielst, ziehst du zuerst 2 Karten nach. Dann muss sich, beginnend mit dem Spieler zu deiner Linken, reihum jeder Mitspieler die oberste Ruinen-Karte vom Ruinen-Stapel nehmen. Ist oder wird der Ruinen-Stapel während des Nehmens leer, können keine weiteren Ruinen mehr genommen werden. Danach darfst du einen weiteren Kultisten aus deiner Hand ausspielen. Das Ausspielen des ursprünglichen Kultisten kostet wie üblich eine Aktion. Das Ausspielen weiterer Kultisten durch die Kartenanweisung benötigt und verbraucht keine Aktion(en). \\ \smallskip
	Wenn du einen Kultisten entsorgst, ziehst du 3 Karten nach. Dabei ist egal, ob du den Kultisten in deinem eigenen Zug oder im Zug eines Mitspielers entsorgst. Wenn du den Kultisten entsorgst, während du Karten aufdeckst, wie z. B. bei einem Ritter-Angriff, handelst du zuerst diese Aktion ab und legst die übrigen aufgedeckten Karten ab. Du ziehst also nicht die aufgedeckten Karten nach, sondern die nachfolgenden.}
\end{tikzpicture}
\hspace{-1cm}
\begin{tikzpicture}
\card
\cardstrip
\cardbanner{banner/white.png}
\cardicon{banner/coin.png}
\cardprice{3}
\cardtitle{Lagerraum}
\cardcontent{Wenn du diese Karte ausspielst, darfst du zuerst eine beliebige Anzahl Karten (auch 0) aus deiner Hand auf den Ablagestapel legen und ebensoviele Karten nachziehen. Unabhängig davon, ob und wie viele Karten du abgelegt hast, darfst du nochmals beliebig viele deiner verbleibenden Handkarten (auch 0) ablegen und erhältst für jede dabei abgelegten Karte +\coin{1}. Weiterhin erhältst du +1 Kauf für die folgende Kaufphase.}
\end{tikzpicture}
\hspace{-1cm}
\begin{tikzpicture}
\card
\cardstrip
\cardbanner{banner/white.png}
\cardicon{banner/coin.png}
\cardprice{2}
\cardtitle{Landstreicher}
\cardcontent{Wenn du diese Karte ausspielst, ziehst du zuerst eine Karte nach. Danach deckst du die oberste Karte von deinem Nachziehstapel auf. Ist es ein Fluch, eine Ruinen-Karte, eine Unterschlupf-Karte oder eine Punktekarte, nimm diese Karte auf die Hand. Dies gilt jeweils auch für kombinierte Kartentypen, die mindestens einen der genannten Typen enthalten. Ansonsten lege die Karte zurück auf den Nachziehstapel.}
\end{tikzpicture}
\hspace{-1cm}
\begin{tikzpicture}
\card
\cardstrip
\cardbanner{banner/green.png}
\cardicon{banner/coin.png}
\cardprice{4}
\cardtitle{Lehen}
\cardcontent{Diese Königreichkarte ist eine Punktekarte, keine Aktionskarte. Sie hat bis zum Ende des Spiels keine Funktion. Bei der Wertung zählt sie 1 Siegpunkt für je volle 3 Silber in deinem gesamten Kartensatz (Nachziehstapel, Ablagestapel und Handkarten). Bei Spielende suchst du aus deinem gesamten Kartensatz alle Silber heraus und zählst sie. Das Ergebnis durch 3 geteilt und abgerundet ergibt die Punkte. Im Spiel zu 3. und 4. werden 12 Karten Lehen verwendet, im Spiel zu 2. werden 8 Karten Lehen verwendet. Wenn du das Lehen entsorgst (egal ob in deinem eigenen Zug oder im Zug eines Mitspielers), nimmst du dir 3 Silber vom Vorrat und legst diese auf deinen Ablagestapel. Sind nicht mehr genug Silber im Vorrat, nimmst du dir nur so viele wie möglich.}
\end{tikzpicture}
\hspace{-1cm}
\begin{tikzpicture}
\card
\cardstrip
\cardbanner{banner/white.png}
\cardicon{banner/coin.png}
\cardprice{4}
\cardtitle{Leichenkarren}
\cardcontent{Wenn du den Leichenkarren ausspielst, musst du dich entscheiden, entweder genau eine Aktionskarte aus deiner Hand zu entsorgen oder diesen ausgespielten Leichenkarren zu entsorgen. Kombinierte Aktionskarten sind auch Aktionskarten. \\ \smallskip
	Wenn du den Leichenkarren nimmst (durch Kauf oder durch eine andere Aktion), nimmst du dir zusätzlich die beiden obersten Ruinen-Karten vom Ruinen-Stapel und legst diese auf deinen Ablagestapel. Sind weniger als 2 Ruinen-Karten übrig, nimmst du die restlichen. Der Spieler, der den Leichenkarren nimmt, nimmt auch die Ruinen-Karten. Spielt ein Spieler die Besessenheit (Dominion – Die Alchemisten) aus und lässt sein „Opfer“ einen Leichenkarren kaufen, so nimmt der Spieler, der die Besessenheit gespielt hat, am Ende den Leichenkarren und damit auch die beiden Ruinen. Verwendet ein Spieler den Fahrenden Händler (Dominion – Hinterland) und nimmt ein Silber anstatt dem Leichenkarren, so nimmt er sich keine Ruinen. Spielst du einen Botschafter (Dominion – Seaside) um Leichenkarren an deine Mitspieler zu verteilen, so nimmt sich jeder Spieler, der einen Leichenkarren nimmt, auch 2 Ruinen (wenn möglich). Wird ein Leichenkarren durch die Maskerade (Dominion – Die Intrige) weitergegeben, nimmt sich der Spieler keine Ruinen, da dies nicht als „nehmen“ gilt. Weiterhin erhältst du +\coin{5} für die folgende Kaufphase.}
\end{tikzpicture}
\hspace{-1cm}
\begin{tikzpicture}
\card
\cardstrip
\cardbanner{banner/white.png}
\cardicon{banner/coin.png}
\cardprice{4}
\cardtitle{\footnotesize{Lumpensammler}}
\cardcontent{\emph{Errata:} Der Kartentext ist falsch, es sollte „[...] Du darfst sofort deinen kompletten Nachziehstapel auf den Ablagestapel legen.“ statt „[...] Du darfst sofort deinen kompletten Nachziehstapel ablegen.“ heißen.\\ \smallskip
	Du darfst deinen kompletten Nachziehstapel auf den Ablagestapel ablegen, musst dies jedoch nicht. Du musst jedoch eine Karte aus deinem Ablagestapel wählen und auf deinen Nachziehstapel legen, ausser der Ablagestapel ist leer. Ist dein Nachziehstapel leer, z. B. weil du ihn gerade komplett abgelegt hast, legst du die Karte an die Stelle des Nachziehstapels. Weiterhin erhältst du +\coin{2} für die folgende Kaufphase. \\ \smallskip
	\emph{Anmerkung:} Durch das direkte Ablegen wird die Karte Tunnel (Dominion – Hinterland) nicht ausgelöst.}
\end{tikzpicture}
\hspace{-1cm}
\begin{tikzpicture}
\card
\cardstrip
\cardbanner{banner/blue.png}
\cardicon{banner/coin.png}
\cardprice{3}
\cardtitle{Marktplatz}
\cardcontent{Wenn du diese Karte ausspielst ziehst du eine Karte nach. Dann darfst du eine weitere Aktion ausführen und erhältst +1 Kauf für die folgende Kaufphase.\\ \smallskip
	Wenn du eine Karte entsorgst, darfst du den Marktplatz aus deiner Hand auf deinen Ablagestapel legen. Wenn du das machst, nimmst du 1 Gold vom Vorrat und legst es auf deinen Ablagestapel. Ist kein Gold mehr im Vorrat, nimmst du dir keines. Du darfst mehrere Karten Marktplatz aus deiner Hand ablegen, auch wenn du nur eine Karte entsorgst.}
\end{tikzpicture}
\hspace{-1cm}
\begin{tikzpicture}
\card
\cardstrip
\cardbanner{banner/white.png}
\cardicon{banner/coin.png}
\cardprice{4}
\cardtitle{Marodeur}
\cardcontent{Wenn du diese Karte spielst, nimmst du dir zuerst eine Karte vom Beute-Stapel und legst sie auf deinen Ablagestapel. Ist der Beute-Stapel aufgebraucht, nimmst du dir keine Beute. Dann muss sich, beginnend mit dem Spieler zu deiner Linken, reihum jeder Mitspieler die oberste Karte vom Ruinen-Stapel nehmen. Ist oder wird der Ruinen-Stapel während des Nehmens leer, können keine weiteren Ruinen mehr genommen werden.}
\end{tikzpicture}
\hspace{-1cm}
\begin{tikzpicture}
\card
\cardstrip
\cardbanner{banner/white.png}
\cardicon{banner/coin.png}
\cardprice{5}
\cardtitle{Medium}
\cardcontent{Wenn du diese Karte spielst, benennst du eine Karte (z. B. „Kupfer“, \emph{nicht} „Geld“) und deckst die oberste Karte von deinem Nachziehstapel auf. Wenn es sich um die benannte Karte handelt, nimmst du sie auf die Hand. Wenn nicht, legst du sie zurück auf den Nachziehstapel. Die benannte Karte muss eindeutig erkennbar sein, Sir Destry ist z. B. \emph{nicht} identisch mit Sir Martin. Du darfst auch eine Karte benennen, die nicht im Spiel ist.\\ \smallskip
	Danach darfst du eine weitere Aktion ausführen und erhältst +\coin{2} für die folgende Kaufphase.}
\end{tikzpicture}
\hspace{-1cm}
\begin{tikzpicture}
\card
\cardstrip
\cardbanner{banner/white.png}
\cardicon{banner/coin.png}
\cardprice{3}
\cardtitle{Mundraub}
\cardcontent{Wenn du diese Karte ausspielst erhältst du +1 Aktion und +1 Kauf für die folgende Kaufphase. Dann musst du eine beliebige Karte aus deiner Hand entsorgen. Wenn du keine Handkarten hast, musst du nichts entsorgen. Danach siehst du dir den Müll-Stapel durch und zählst wie viele Geldkarten mit unterschiedlichem Namen sich darin befinden. Für jede davon erhältst du +\coin{1} für die folgende Kaufphase. Befinden sich im Müll-Stapel z. B. 4 Kupfer, 1 Falschgeld und 6 Anwesen, so erhältst du +\coin{2}. Kombinierte Geldkarten (z. B. Harem, Dominion – Die Intrige) sind auch Geldkarten.}
\end{tikzpicture}
\hspace{-1cm}
\begin{tikzpicture}
\card
\cardstrip
\cardbanner{banner/white.png}
\cardicon{banner/coin.png}
\cardprice{5}
\cardtitle{Neubau}
\cardcontent{Du benennst zuerst eine beliebige Karte. Es muss keine Punktekarte sein. Du darfst sogar eine Karte benennen, die in diesem Spiel nicht verwendet wird. Dann deckst du solange Karten von deinem Nachziehstapel auf, bis eine Punktekarte offen liegt, die nicht die zuvor benannte Karte ist. (Das Benennen einer Karte dient dazu, diese Karte zu schützen.) Wenn du auch nach dem Mischen des Ablagestapels keine solche Karte aufdecken kannst, passiert nichts weiter. Wenn du eine Punktekarte aufdeckst, die nicht die zuvor benannte Karte ist, entsorgst du diese Karte und nimmst dir eine Punktekarte, die bis zu \coin{3} mehr kostet, als die entsorgte Karte aus dem Vorrat und legst diese auf deinen Ablagestapel.}
\end{tikzpicture}
\hspace{-1cm}
\begin{tikzpicture}
\card
\cardstrip
\cardbanner{banner/white.png}
\cardicon{banner/coin.png}
\cardprice{4}
\cardtitle{Prozession}
\cardcontent{
	\begin{singlespace}\tiny{
		\emph{Errata:} Auf der Karte steht, dass man sich eine Karte nehmen muss, die genau \coin{1} mehr kostet. Richtig wäre: ...nimm dir eine Aktionskarte, die genau \coin{1} mehr kostet. \\ \smallskip
		Du darfst eine Aktionskarte, die du noch auf der Hand hast, zweimal ausspielen. Du spielst die Karte zuerst aus und führst die Anweisungen komplett aus. Dann nimmst du die Karte zurück auf die Hand, spielst sie ein zweites Mal aus und führst die Anweisungen nochmal komplett aus. Das zweimalige Ausspielen der Aktionskarte kostet keine weiteren Aktionen. Das ursprüngliche Ausspielen der Prozession kostet wie üblich eine Aktion. Nachdem du die Karte zum zweiten Mal ausgeführt hast musst du sie entsorgen und dir dafür eine Aktionskarte, die genau \coin{1} mehr kostet als die entsorgte Karte, aus dem Vorrat nehmen (wenn möglich) und auf deinen Ablagestapel legen. Du musst dies auch machen, wenn du im Verlauf der Aktion den \emph{„Anschluss verloren“} hast (siehe: Neue Regeln). Du darfst andere Karten erst ausspielen, wenn die Prozession vollständig abgehandelt ist. Spielst du z. B. eine Festung auf eine Prozession, so führst du die zusätzlichen insgesamt 4 Aktionen erst aus, nachdem du die Festung entsorgt und dir eine neue Karte vom Vorrat genommen hast. Wenn du eine Prozession auf eine andere Prozession spielst, spielst du 2 weitere Aktionskarten aus deiner Hand je 2mal aus, entsorgst diese beiden Aktionskarten und nimmst dir 2 Karten vom Vorrat. Dann entsorgst du die zweite Prozession und nimmst dir auch dafür eine Karte vom Vorrat. Wenn du eine Dauer-Karte auf eine Prozession spielst, bleibt die Prozession bis zum nächsten Zug liegen, dann führst du den Effekt der Dauerkarte zweimal aus (obwohl die Dauerkarte bereits entsorgt ist).
		} \end{singlespace}
	}
\end{tikzpicture}
\hspace{-1cm}
\begin{tikzpicture}
\card
\cardstrip
\cardbanner{banner/white.png}
\cardicon{banner/coin.png}
\cardprice{4}
\cardtitle{Ratten}
\cardcontent{Wenn du diese Karte ausspielst, ziehst du zuerst eine Karte nach. Dann nimmst du dir eine weitere Karte Ratten aus dem Vorrat und legst diese auf deinen Ablagestapel. Sind keine Ratten mehr im Vorrat, nimmst du dir keine. Danach musst du eine Karte aus deiner Hand entsorgen, jedoch keine Ratten. Hast du nur Ratten oder keine Karten auf der Hand, entsorgst du keine Karte. Hast du nur Ratten auf der Hand, musst du deine Kartenhand den anderen Spielern vorzeigen.\\ \smallskip
	Wenn du die Ratten entsorgst, ziehst du sofort eine Karte nach. Dabei ist egal, ob du die Ratten in deinem eigenen Zug oder im Zug eines Mitspielers entsorgst.}
\end{tikzpicture}
\hspace{-1cm}
\begin{tikzpicture}
\card
\cardstrip
\cardbanner{banner/white.png}
\cardicon{banner/coin.png}
\cardprice{5}
\cardtitle{Raubzug}
\cardcontent{Du entsorgst diese Karte direkt nachdem du sie ausgespielt hast. Dann muss jeder Mitspieler, beginnend mit dem Spieler zu deiner Linken, seine komplette Kartenhand aufdecken. Du wählst bei jedem Mitspieler eine seiner Handkarten, die dieser ablegen muss. Danach nimmst du dir 2 Karten vom Beute-Stapel neben dem Vorrat und legst diese auf deinen Ablagestapel. Sind nicht mehr genügend Karten im Beute-Stapel, nimmst du nur so viele wie möglich.}
\end{tikzpicture}
\hspace{-1cm}
\begin{tikzpicture}
\card
\cardstrip
\cardbanner{banner/white.png}
\cardicon{banner/coin.png}
\cardprice{5}
\cardtitle{Ritter}
\cardcontent{\tiny{\begin{spacing}{0.1cm}Im Ritter-Stapel befinden sich bei Spielbeginn 10 Karten mit jeweils individuellem Namen. Alle Ritter haben die selbe Basisfähigkeit, zusätzlich jedoch noch eine individuelle Anweisung. Wenn die Ritter als Königreichkarte für ein Spiel gewählt werden, werden die Karten mit dem Typ „Ritter“ gemischt, als verdeckter Stapel in den Vorrat gelegt und die oberste Karte davon aufgedeckt. Siehe hierzu auch unter Spielvorbereitung „gemischte Stapel“. Auf den Ritter-Stapel wird kein Marker für die Handelsroute (Dominion – Blütezeit) gelegt, auch wenn die oberste Karte eine Punktekarte ist. Bis auf diese Ausnahmen wird der Ritter-Stapel behandelt, wie jeder andere der 10 Kartenstapel. Er ist Teil des Vorrats und wird bei der Spielendebedingung beachtet. Die Anweisungen auf den Ritterkarten werden wie üblich von oben nach unten ausgeführt. Sir Michael z. B. lässt die Mitspieler Karten ablegen, bevor seine Basisfähigkeit ausgeführt wird. Die Basisfähigkeit, die alle Ritter gemeinsam haben, wird wie folgt ausgeführt: Beginnend mit dem Spieler links von dir, muss jeder Mitspieler die obersten beiden Karten von seinem Nachziehstapel aufdecken. Hat der Mitspieler keine Karte aufgedeckt, die 3 - 6 Geld kostet, muss er keine Karte entsorgen. Er legt beide aufgedeckten Karten ab. Kostet genau eine der beiden aufgedeckten Karten 3 - 6 Geld, so muss er diese Karte entsorgen, die andere Karte legt er ab. Kosten beide aufgedeckten Karten 3 - 6 Geld, so darf der Mitspieler entscheiden, welche der beiden Karten er entsorgen muss. Die andere aufgedeckte Karte legt er ab. Entsorgt ein Mitspieler auf diese Weise eine andere Ritter-Karte, so musst auch du den Ritter entsorgen, der den Angriff ausgelöst hat. Die Mitspieler dürfen keine Karten mit Trank in den Kosten (Dominion – Die Alchemisten) entsorgen. Sollen die Ritter im Schwarzmarktstapel verwendet werden, so wird einer der 10 Ritter zufällig gezogen und in den Schwarzmarktstapel eingemischt.\\ \smallskip
		Die individuellen Fähigkeiten der einzelnen Ritter sind selbsterklärend. Darum hier nur kurze Anmerkungen zu einzel- nen Rittern. \emph{Sir Martin:} Es ist beabsichtigt, dass Sir Martin nur 4 kostet. Sir Martin ist der einzige Ritter, den ein Spieler z. B. durch die Waffenkammer nehmen kann (natürlich nur, wenn er offen liegt). \emph{Sir Vander:} Das Gold nimmst du aus dem Vorrat und legst es auf deinen Ablagestapel. Es ist egal, ob du Sir Vander in deinem eigenen Zug oder im Zug eines Mitspielers entsorgst. \emph{Dame Natalie:} Die Karte darfst (emph{Errata:} auf der Karte ist die Anweisung fälschlicherweise nicht optional) du aus dem Vorrat nehmen und sie dann auf deinen Ablagestapel legen. \emph{Dame Anna:} Du musst keine Karten entsorgen. \emph{Dame Josephine} ist eine Punktekarte. Der Ritter-Stapel ist jedoch dadurch kein Punktestapel. Wird die Handelsroute (Dominion – Blütezeit) verwendet, wird kein Marker auf den Ritter-Stapel gelegt.\end{spacing}}}
\end{tikzpicture}
\hspace{-1cm}
\begin{tikzpicture}
\card
\cardstrip
\cardbanner{banner/brown.png}
\cardicon{banner/coin.png}
\cardprice{0}
\cardtitle{Ruinen}
\cardcontent{Der Ruinen-Stapel besteht aus bis zu 5 unterschiedlichen Karten. Siehe hierzu auch unter Spielvorbereitung „gemischte Stapel“ und „Ruinen“. Der Ruinen-Stapel ist Teil des Vorrats. Die offenliegende Karte (und nur diese) kann bzw. muss durch eine entsprechende Kartenanweisung genommen oder auch gekauft werden. \\ \smallskip
	Wenn du die Überlebenden ausspielst, siehst du dir die obersten beiden Karten von deinem Nachziehstapel an und entscheidest dich dann, ob du beide Karten auf deinen Ablagestapel legst oder ob du beide Karten in beliebiger Reihenfolge zurück auf deinen Nachziehstapel legst. Du darfst nicht eine Karte ablegen und die andere zurück auf den Nachziehstapel legen. \\ \smallskip
	Die übrigen 4 Ruinen-Karten sind selbsterklärend.}
\end{tikzpicture}
\hspace{-1cm}
\begin{tikzpicture}
\card
\cardstrip
\cardbanner{banner/white.png}
\cardicon{banner/coin.png}
\cardprice{5}
\cardtitle{\footnotesize{Schrotthändler}}
\cardcontent{Wenn du diese Karte ausspielst, ziehst du zuerst eine Karte nach. Dann musst du eine Karte entsorgen, wenn du mindestens eine Karte auf der Hand hast. Danach darfst du eine weitere Aktion ausführen und erhältst +\coin{1} für die folgende Kaufphase.}
\end{tikzpicture}
\hspace{-1cm}
\begin{tikzpicture}
\card
\cardstrip
\cardbanner{banner/white.png}
\cardicon{banner/coin.png}
\cardprice{5}
\cardtitle{Schurke}
\cardcontent{Wenn im Müll-Stapel mindestens eine Karte ist, die \coin{3} - \coin{6} kostet, musst du eine davon nehmen und auf deinen Ablagestapel legen. Du darfst nicht darauf verzichten. Du darfst den Müllstapel durchsehen, bevor du den Schurken ausspielst. Du zeigst die Karte, die du nimmst deinen Mitspielern. Du darfst keine Karten mit Trank in den Kosten (Dominion – Die Alchemisten) nehmen.\\ \smallskip
	Wenn keine Karte im Müll-Stapel ist, die \coin{3} - \coin{6} kostet, nimmst du dir keine Karte. Stattdessen muss jeder Mitspieler, beginnend mit dem Spieler links von dir, die obersten beiden Karten von seinem Nachziehstapel aufdecken und eine Karte davon entsorgen, die \coin{3} - \coin{6} kostet. Hat der Mitspieler keine Karte aufgedeckt, die 3 - 6 Geld kostet, muss er keine Karte entsorgen. Er legt beide aufgedeckten Karten ab. Kostet genau eine der beiden aufgedeckten Karten \coin{3} - \coin{6}, so muss er diese Karte entsorgen, die andere Karte legt er ab. Kosten beide aufgedeckten Karten \coin{3} - \coin{6}, so darf der Mitspieler entscheiden, welche der beiden Karten er entsorgen muss. Die andere aufgedeckte Karte legt er ab.}
\end{tikzpicture}
\hspace{-1cm}
\begin{tikzpicture}
\card
\cardstrip
\cardbanner{banner/white.png}
\cardicon{banner/coin.png}
\cardprice{0*}
\cardtitle{Söldner}
\cardcontent{\begin{singlespace}\tiny{Wird der Gassenjunge im Spiel verwendet, so wird auch der Söldner-Stapel benötigt. Der Söldner-Stapel wird \emph{neben} dem Vorrat bereit gelegt. Die Karten vom Söldner-Stapel können nur durch die Anweisung auf dem Gassenjunge genommen werden. Auf andere Weise können keine Karten vom Söldner-Stapel gekauft oder genommen werden. Der Botschafter (Dominion – Seaside) darf keine Karten auf den Söldner-Stapel zurück legen. Der Söldner-Stapel wird für die Spielende-Bedingung nicht beachtet.\\ \smallskip
	Wenn du den Söldner ausspielst, darfst du 2 Karten aus deiner Hand entsorgen. Wenn du das machst, musst du zuerst 2 Karten nachziehen und erhältst +\coin{2} für die folgende Kaufphase. Dann müssen alle deine Mitspieler, beginnend mit dem Spieler links von dir, solange Handkarten ablegen, bis sie nur noch höchstens 3 Karten auf der Hand haben. Mitspieler, die bereits 3 oder weniger Karten auf der Hand haben, legen keine Karten mehr ab. Spieler, die mit dem Bettler auf diesen Angriff reagieren wollen, müssen diesen ablegen, bevor du dich entscheidest 2 Karten zu entsorgen oder nicht. Wenn du dich dafür entscheidest, 2 Karten zu entsorgen, jedoch nur noch 1 Karte auf der Hand hast. Entsorgst du nur diese 1 Karte. Da die nachfolgenden Anweisungen jedoch an das entsorgen von 2 Karten gebunden ist, ziehst du keine Karten nach und erhältst auch kein virtuelles Geld. Wenn die Karten, die du entsorgst, Entsorgungs-Anweisungen haben, entsorgst du zuerst beide und führst die Anweisungen dann in beliebiger Reihenfolge aus.}\end{singlespace}}
\end{tikzpicture}
\hspace{-1cm}
\begin{tikzpicture}
\card
\cardstrip
\cardbanner{banner/red.png}
\cardicon{banner/coin.png}
\cardprice{1}
\cardtitle{Unterschlupf}
\cardcontent{\begin{singlespace}\tiny{Unterschlupfkarten haben keinen Stapel (weder im Vorrat noch ausserhalb des Vorrats), sie können niemals gekauft werden.\\ \smallskip
		\emph{Hütte:} Diese Karte ist ein Unterschlupf, siehe „Besondere Karten“ und eine Reaktionskarte. Die Hütte ist niemals im Vorrat. Wenn du eine Punktekarte kaufst und die Hütte auf der Hand hast, darfst du die Hütte entsorgen. Dies gilt auch, wenn du eine kombinierte Punktekarte kaufst. Du erhältst keine weiteren Vorteile durch das Entsorgen der Hütte. Du wirst nur eine Karte los.\\ \smallskip
		\emph{Totenstadt:} Diese Karte ist ein Unterschlupf, siehe „Besondere Karten“ und eine Aktionskarte. Die Totenstadt ist niemals im Vorrat. Wenn du diese Karte ausspielst, erhältst du +2 Aktionen.\\ \smallskip
		\emph{Verfallenes Anwesen:} Diese Karte ist ein Unterschlupf, siehe „Besondere Karten“ und eine Punktekarte mit dem Wert 0 Siegpunkte. Das Verfallene Anwesen ist niemals im Vorrat. Wenn du diese Karte entsorgst, ziehst du sofort eine Karte nach, auch wenn du gerade eine andere Aktionskarte ausführst. Wenn du das Verfallene Anwesen z. B. durch den Altar entsorgst, ziehst du zuerst eine Karte nach und nimmst dir danach eine Karte vom Vorrat. Du darfst das Verfallene Anwesen nicht „freiwillig“ entsorgen, sondern benötigst wie üblich eine andere Karte mit der Anweisung, die dir erlaubt, eine Karte zu entsorgen.}\end{singlespace}}
\end{tikzpicture}
\hspace{-1cm}
\begin{tikzpicture}
\card
\cardstrip
\cardbanner{banner/white.png}
\cardicon{banner/coin.png}
\cardprice{0*}
\cardtitle{Verrückter}
\cardcontent{Wird der Eremit im Spiel verwendet, so wird auch der Verrückten-Stapel benötigt. Der Verrückten-Stapel wird \emph{neben} dem Vorrat bereit gelegt. Die Karten vom Verrückten-Stapel können nur durch die Anweisung auf dem Eremiten genommen werden. Auf andere Weise können keine Karten vom Verrückten-Stapel gekauft oder genommen werden. Der Botschafter (Dominion – Seaside) darf keine Karten auf den Verrückten-Stapel zurück legen. Der Verrückten-Stapel wird für die Spielende-Bedingung \emph{nicht} beachtet.\\ \smallskip
	Diese Karte ist nicht Teil des Vorrats. Sie kann nur durch die Anweisung auf dem Eremiten genommen werden. Auf eine andere Art kann der Verrückte nicht genommen oder gekauft werden. Wenn du den Verrückten ausspielst, erhältst du +2 Aktionen. Dann legst du den Verrückten normalerweise auf den Verrückten-Stapel zurück und ziehst pro Karte, die du noch auf der Hand hast, eine Karte nach (du verdoppelst also die Anzahl deiner Handkarten). Es kann jedoch vorkommen, dass du den Verrückten nicht zurück legen kannst, weil du den „Anschluss verloren“ hast (siehe Neue Regeln), z. B. weil du den Verrückten auf eine Prozession oder einen Thronsaal (Dominion – Basisspiel) ausgespielt hast.}
\end{tikzpicture}
\hspace{-1cm}
\begin{tikzpicture}
\card
\cardstrip
\cardbanner{banner/white.png}
\cardicon{banner/coin.png}
\cardprice{5}
\cardtitle{Vogelfreie}
\cardcontent{\begin{singlespace}\tiny{Wenn du diese Karte ausspielst, wählst du eine Karte aus dem Vorrat, die weniger kostet, als die Vogelfreien selbst (also normalerweise bis zu 4 Geld). Dann behandelst du die vor dir liegende Karte Vogelfreie genau so, als wäre es die gewählte Karte. Normalerweise führst du einfach nur die Anweisungen auf der gewählten Karte aus. Wählst du z. B. die Festung, so ziehst du eine Karte nach und führst danach bis zu 2 weitere Aktionen aus. Die Vogelfreien nehmen auch die Kosten, den Namen und den oder die Kartentypen der gewählten Karte an.\\ \smallskip
		Wenn du mit den Vogelfreien eine Karte wählst, die sich selbst entsorgt, so entsorgst du die Vogelfreien. Wenn die Karte nicht mehr vor dir liegt, wird sie wieder zu der Karte Vogelfreie selbst. Wenn du die Vogelfreien als Dauer-Karte (Dominion – Seaside) verwendest, bleibt die Karte bis zu deinem nächsten Zug im Spiel. Wenn du die Vogelfreien als Thronsaal (Dominion – Basisspiel), Königshof (Dominion – Blütezeit) oder Prozession verwendest um eine Dauer-Karte mehrmals auszuspielen, bleiben die Vogelfreien ebenso im Spiel. Wenn du die Vogelfreien mehrmals ausspielst (z. B. durch den Thronsaal), wählst du nur beim ersten Ausspielen eine Karte. Die Vogelfreien bleiben diese Karte auch beim nachfolgenden Ausspielen. Nutzt du z. B. die Prozession, um die Vogelfreien zweimal auszuspielen und kopierst die Festung, so ziehst du insgesamt 2 Karten nach und erhältst +4 Aktionen. Danach entsorgst du die ausliegende Karte, nimmst sie aber sofort wieder auf die Hand zurück, da es noch immer eine Festung ist. Zurück auf deiner Hand ist die Karte aus dem Spiel und wird wieder zu den Vogelfreien. Nun nimmst du dir noch eine Karte aus dem Vorrat, die genau 1 Geld mehr kostet, als die entsorgte Karte. Die Karte, die du entsorgt hast, ist inzwischen wieder die Karte Vogelfreie 5 Geld. Du nimmst dir also eine Karte die 6 Geld kostet. Wenn du die Vogelfreien als eine Karte nutzt, die etwas in der Aufräumphase macht (z. B. der Eremit), so führst du diese Anweisung aus. Wenn du das Füllhorn (Dominion – Reiche Ernte) ausspielst, werden die Vogelfreien als die Karte, die sie kopieren angesehen. Wenn du z. B. mit einer Karte Vogelfreie die Festung wählst und mit einer weiteren Karte Vogelfreie den Lumpensammler, nimmst du dir eine Karte, die bis zu 3 Geld kostet. Du darfst mit den Vogelfreien nur eine Karte wählen, die sichtbar im Vorrat liegt. Du darfst keine Karte wählen, deren Stapel leer ist. Du darfst keine Karte wählen, die nicht im Vorrat ist, wie z. B. den Söldner. In gemischten Stapeln, wie z. B. Ruinen oder Ritter, darfst du nur die oben liegende wählen.}\end{singlespace}}
\end{tikzpicture}
\hspace{-1cm}
\begin{tikzpicture}
\card
\cardstrip
\cardbanner{banner/white.png}
\cardicon{banner/coin.png}
\cardprice{4}
\cardtitle{Waffenkammer}
\cardcontent{Wenn du diese Karte ausspielst, nimmst du dir eine Karte, die bis zu \coin{4} kostet vom Vorrat. Du legst diese sofort verdeckt auf deinen Nachziehstapel, anstatt auf den Ablagestapel wie üblich.}
\end{tikzpicture}
\hspace{-1cm}
\begin{tikzpicture}
\card
\cardstrip
\cardbanner{banner/white.png}
\cardicon{banner/coin.png}
\cardprice{3}
\cardtitle{Weiser}
\cardcontent{Wenn du diese Karte ausspielst, erhältst du immer +1 Aktion, dann deckst du solange Karten von deinem Nachziehstapel auf, bis du eine Karte aufdeckst, die mindestens \coin{3} kostet und nimmst diese Karte auf die Hand. Lege die übrigen durch diese Aktion aufgedeckten Karten ab. Wenn du auch nach dem Mischen deines Ablagestapels keine Karte, die mindestens \coin{3} kostet, aufdecken kannst, legst du alle aufgedeckten Karten ab und nimmst keine Karte auf die Hand. Deckst du z. B. ein Kupfer, dann einen Fluch und dann eine Provinz auf, nimmst du die Provinz auf die Hand und legst Kupfer und Fluch ab.}
\end{tikzpicture}
\hspace{-1cm}
\begin{tikzpicture}
\card
\cardstrip
\cardbanner{banner/white.png}
\cardicon{}
\cardprice{}
\cardtitle{\scriptsize{Empfohlene 10er Sätze\qquad}}
\cardcontent{\emph{Dark Ages:}\\ \smallskip
	\emph{Leichenzug:} \\ 
	Festung, Jagdgründe, Katakomben, Kultist, Marktplatz, Mundraub, Prozession, Ritter, Vogelfreie, Waffenkammer \\
	\smallskip \emph{Spiel mit dem Teufel} \\ 
	Banditenlager, Grabräuber, Lagerraum, Landstreicher, Lumpensammler, Medium, Ratten, Raubzug, Schrotthändler, Weiser\\
	\smallskip \emph{Dark Ages und Basisspiel:} \\ 
	\smallskip \emph{Auf und Ab:} \\ 
	Armenhaus, Barde, Eremit, Jagdgründe, Medium / Geldverleiher, Hexe, Keller, Thronsaal, Werkstatt \\
	\smallskip \emph{Ritterspiele:} \\ 
	Altar, Knappe, Lumpensammler, Ratten, Ritter / Bibliothek, Gärten, Jahrmarkt, Laboratorium, Umbau \\
}
\end{tikzpicture}
\hspace{-1cm}
\begin{tikzpicture}
\card
\cardstrip
\cardbanner{banner/white.png}
\cardicon{}
\cardprice{}
\cardtitle{\scriptsize{Empfohlene 10er Sätze\qquad}}
\cardcontent{\emph{Dark Ages und Die Intrige:}\\ \smallskip
	\emph{Prophezeiung:} \\ 
	Eisenhändler, Landstreicher, Medium, Neubau, Waffenkammer / Adelige, Baron, Große Halle, Verschwörer, Wunschbrunnen \\
	\smallskip \emph{Invasion:} \\ 
	Bettler, Gassenjunge, Knappe, Marodeur, Schurke / Anbau, Bergwerksdorf, Harem, Kerkermeister, Trickser \\
	\smallskip \emph{Dark Ages und Seaside:} \\ 
	\smallskip \emph{Nasses Grab:} \\ 
	Eremit, Gassenjunge, Grabräuber, Graf, Lumpensammler / Eingeborenendorf, Müllverwerter, Piratenschiff, Schatzkarte, Schatzkammer \\
	\smallskip \emph{Einfaches Volk:} \\ 
	Armenhaus, Gassenjunge, Landstreicher, Lehen, Leichenkarren / Fischerdorf, Hafen, Insel, Ausguck, Lagerhaus \\
}
\end{tikzpicture}
\hspace{-1cm}
\begin{tikzpicture}
\card
\cardstrip
\cardbanner{banner/white.png}
\cardicon{}
\cardprice{}
\cardtitle{\scriptsize{Empfohlene 10er Sätze\qquad}}
\cardcontent{\emph{Dark Ages und die Alchemisten:}\\ \smallskip
	\emph{Seuchenherd:} \\ 
	Barde, Kultist, Lehen, Marktplatz, Ratten, Waffenkammer / Lehrling, Vision, Verwandlung, Weinberg\\
	\smallskip \emph{Klagelied:} \\ 
	Bettler, Eisenhändler, Falschgeld, Katakomben, Mundraub, Raubzug / Apotheker, Golem, Kräuterkundiger, Universität \\
	\smallskip \emph{Dark Ages und Blütezeit:} \\ 
	\smallskip \emph{Des einen Müll ...:} \\ 
	Falschgeld, Grabräuber, Marktplatz, Mundraub, Schurke / Abenteuer, Denkmal, Großer Markt, Stadt, Talisman \\
	\smallskip \emph{Ehre unter Dieben:} \\ 
	Banditenlager, Knappe, Neubau, Prozession, Schurke / Kunstschmiede, Hort, Hausierer, Steinbruch, Wachturm \\
}
\end{tikzpicture}
\hspace{-1cm}
\begin{tikzpicture}
\card
\cardstrip
\cardbanner{banner/white.png}
\cardicon{}
\cardprice{}
\cardtitle{\scriptsize{Empfohlene 10er Sätze\qquad}}
\cardcontent{\emph{Dark Ages und Reiche Ernte:}\\ \smallskip
	\emph{Dunkler Karneval:} \\ 
	Eremit, Festung, Kultist, Ritter, Schrotthändler, Vogelfreie / Festplatz, Füllhorn, Menagerie, Weiler \\
	\smallskip \emph{Auf den Sieger:} \\ 
	Banditenlager, Barde, Falschgeld, Leichenkarren, Marodeur, Raubzug / Ernte, Treibjagd Nachbau, Turnier \\
	\smallskip \emph{Dark Ages und Hinterland:} \\ 
	\smallskip \emph{Fern von Zuhause:} \\ 
	Barde, Bettler, Graf, Lehen, Marodeur / Botschaft, Feilscher, Aufbau, Kartograph, Katzengold \\
	\smallskip \emph{Expeditionen:} \\ 
	Altar, Armenhaus, Eisenhändler, Katakomben, Lagerraum / Fruchtbares Land, Fernstraße, Gewürzhändler, Tunnel, Wegkreuzung\\
}
\end{tikzpicture}
\hspace{-1cm}
\begin{tikzpicture}
\card
\cardstrip
\cardbanner{banner/white.png}
\cardtitle{\scriptsize{Neue Regeln (1/4)}\qquad}
\cardcontent{\begin{singlespace} \tiny{
		\emph{„ausspielen“:} Ein Spieler, der eine Karte ausspielt, legt diese offen vor sich aus und führt dann die Anweisungen von oben nach unten aus. Auch wenn die Karte nicht vor dem Spieler liegen bleibt (sondern z. B. wie die Beute oder der Verrückte sofort auf den jeweiligen Stapel zurück gelegt werden), werden die Anweisungen soweit möglich vollständig ausgeführt.\\ \smallskip
		\emph{„Im Spiel“ („ausgespielt“):} Aktions- und Geldkarten, die ein Spieler offen vor sich ausgelegt hat, sind „im Spiel“ bis sie aufgeräumt werden (üblicherweise in der Aufräumphase abgelegt werden). Beiseite gelegte und entsorgte Karten, sowie Karten, welche die Spieler auf der Hand haben und Karten im Vorrat und in den Nachzieh- und Ablagestapeln der Spieler sind nicht „im Spiel“. Das Aufdecken einer Reaktionskarte bringt diese nicht „ins Spiel“. \\ \smallskip
		\emph{Kaufphase:} In der Kaufphase muss der Spieler seine Geldkarten einzeln auslegen. Er darf dabei jedoch die Reihenfolge selbst bestimmen. Legt ein Spieler eine Geldkarte mit zusätzlichen Anweisungen aus, führt er diese Anweisungen nach Möglichkeit aus, bevor er eine weitere Geldkarte auslegt. Der Spieler muss alle Geldkarten, die er in dieser Runde auslegen möchte auslegen, bevor er eine Karte kauft (auch wenn er mehrere Käufe zur Verfügung hat). Der Spieler darf in dieser Kaufphase keine Geldkarten mehr auslegen, nachdem er eine Karte gekauft hat. \\ \smallskip
		\emph{Geldwert:} Geldkarten „produzieren“ ihren Wert im Moment, in dem sie ausgelegt werden. Werden Geldkarten noch während der Kaufphase verändert oder abgelegt (z. B. Beute), bleibt der „produzierte“ Geldwert für diese Kaufphase trotzdem erhalten.}
	\end{singlespace}}
\end{tikzpicture}
\hspace{-1cm}
\begin{tikzpicture}
\card
\cardstrip
\cardbanner{banner/white.png}
\cardtitle{\scriptsize{Neue Regeln (2/4)}\qquad}
\cardcontent{\begin{singlespace} \tiny{
	\emph{Geldkarten mit zusätzlichen Anweisungen:} Dark Ages beinhaltet 2 Geldkarten mit zusätzlichen Anweisungen, Falschgeld und Beute. Es handelt sich dabei nicht um Basiskarten, die in jedem Spiel verwendet werden (wie z. B. Gold). Diese Geldkarten können wie üblich in der Kaufphase ausgelegt werden. Sobald ein Spieler eine Geldkarte mit Anweisungen auslegt, führt er die Anweisung auf der Karte aus. Diese Geldkarten sind, wie jede andere Geldkarte, von Anweisungen betroffen, die sich auf Geldkarten beziehen, z. B. Dieb (Dominion-Basisspiel). \\ \smallskip
	\emph{Wenn mehrere Anweisungen gleichzeitig ausgeführt werden können} darfst du entscheiden, in welcher Reihenfolge du diese Anweisungen ausführst. Entsorgst du z. B. eine Karte Ratten und legst den Marktplatz aus deiner Hand ab, so entscheidest du, ob du zuerst die Anweisung der Ratten oder die Anweisung auf dem Marktplatz ausführst. Wenn mehrere Spieler von einer Anweisung betroffen sind (meistens durch Angriffskarten), so wird die Anweisung normalerweise in Spielerreihenfolge ausgeführt.\\ \smallskip
	\emph{Mehrfache Reaktionen:} Jeder Spieler darf auf ein einziges Ereignis (meist ein Angriff ) auch mit mehreren Reaktionskarten reagieren, wenn er diese auf der Hand hat. Wie üblich muss jedoch zuerst eine Karte vollständig ausgeführt werden, bevor der Spieler eine weitere Reaktionskarte aus seiner Hand aufdeckt bzw. ablegt. Ein Spieler darf z. B. als Reaktion auf einen Angriff die Geheimkammer (Dominion – Die Intrige) aufdecken. Nachdem er die Geheimkammer vollständig ausgeführt hat, darf er auch noch den Bettler aus seiner Hand ablegen (auch wenn er diesen gerade durch die Geheimkammer nachgezogen hat) und auch diesen ausführen.}		
	\end{singlespace}}
\end{tikzpicture}
\hspace{-1cm}
\begin{tikzpicture}
\card
\cardstrip
\cardbanner{banner/white.png}
\cardtitle{\scriptsize{Neue Regeln (3/4)}\qquad}
\cardcontent{\begin{singlespace} \tiny{
	\emph{Anschluss verloren:} In einigen Fällen ist es möglich, dass eine Kartenanweisung verlangt, eine Karte zu bewegen (z. B. abzulegen), die Karte jedoch bereits durch eine andere Anweisung anderweitig bewegt (z. B. entsorgt) wurde. Das hat zur Folge, dass du die Karte nicht bewegen kannst. Du führst die Anweisungen auf der Karte trotzdem aus (soweit möglich). Wenn du z. B. auf eine Prozession einen Verrückten spielst, legst du den Verrückten zurück auf den Verrückten-Stapel. Du erhältst +2 Aktionen und ziehst Karten nach. Nun musst du den Verrückten zum zweiten Mal ausspielen. Die Karte liegt jedoch bereits wieder auf dem Verrückten-Stapel. Du erhältst trotzdem nochmals die +2 Aktionen, kannst die Karte jedoch nicht mehr zurück auf den Stapel legen (sie liegt bereits dort). Du ziehst also keine Karten nach, weil diese Anweisung nur aktiv wird, wenn du die Karte auf den Verrückten-Stapel zurück legst. Danach müsstest du den Verrückten entsorgen, da dieser bereits auf seinen Stapel zurück gelegt ist, kannst du ihn nicht auf den Müllstapel legen. Du nimmst dir jedoch trotzdem eine Karte die genau 1 Geld kostet, wenn möglich.\\ \smallskip
	\emph{„Wenn du diese Karte entsorgst: ...“:} Immer, wenn du eine deiner Karten entsorgst, die eine solche Anweisung haben, führst du diese Anweisung aus. Es ist egal, ob du die Karte in deinem Zug oder im Zug eines Mitspielers (z. B. durch eine Angriffskarte) entsorgst. Jeder Spieler entsorgt immer nur seine eigenen Karten. Im Falle eines Angriffs entsorgt also nicht der Spieler, der die Angriffskarte ausgespielt hat, die Karten seiner Mitspieler, sondern immer nur der durch den Angriff betroffene Spieler. Der Spieler, der seine Karte entsorgt, führt dann auch die Anweisung aus, direkt nachdem er die Karte auf den Müll-Stapel gelegt hat. Dies kann dazu führen, dass der Spieler die Anweisung ausführt, bevor eine andere Aktionskarte vollständig ausgeführt wurde. Spielt ein Spieler z. B. einen Grabräuber und nutzt diesen, um einen Kultisten zu entsorgen, so zieht er zuerst 3 Karten für den Kultisten nach und nimmt danach eine Karte durch den Grabräuber. In einigen Fällen werden Karten aus dem Stapel des Spielers entfernt, ohne entsorgt zu werden, z. B. durch den Botschafter (Dominion – Seaside) oder die Maskerade (Dominion – Die Intrige). Wenn durch eine Anweisung mehrere Karten gleichzeitig entsorgt werden, so entsorgst du zunächst die Karten und führst dann die Entsorgungs-Anweisungen in beliebiger Reihenfolge aus.
	}
	\end{singlespace}}
\end{tikzpicture}
\hspace{-1cm}
\begin{tikzpicture}
\card
\cardstrip
\cardbanner{banner/white.png}
\cardtitle{\scriptsize{Neue Regeln (4/4)}\qquad}
\cardcontent{\begin{singlespace} \tiny{
	\emph{Nur zur Erinnerung:} Immer wenn ein Spieler Karten nachziehen müsste, sein Nachziehstapel je- doch leer ist, mischt der Spieler seinen Ablagestapel und legt ihn als neuen Nachziehstapel bereit.\\ \smallskip
	Müsste der Spieler mehr Karten nachziehen, als sein Nachziehstapel noch enthält, so zieht er zuerst die verbleibenden Karten seines Nachziehstapels, mischt dann seinen Ablagestapel und zieht die restlichen Karten nach. Kann er, auch nachdem er seinen Ablagestapel gemischt und als neuen Nachziehstapel bereit gelegt hat, nicht genügend Karten nachziehen, so zieht er nur so viele Karten, wie möglich. \\ \smallskip
	Gleiches gilt, wenn der Spieler Karten von seinem Nachziehstapel aufdecken oder ansehen muss. Muss ein Spieler Karten auf seinen Nachziehstapel legen und ist sein Nachziehstapel in diesem Moment leer, so legt er diese Karten an die Stelle seines Nachziehstapels. Er mischt seinen Ablagestapel nicht.
	}
	\end{singlespace}}
\end{tikzpicture}
\hspace{1cm}
	    % Basic settings for this card set
\renewcommand{\cardcolor}{guilds}
\renewcommand{\cardextension}{Erweiterung VII}
\renewcommand{\cardextensiontitle}{Die Gilden}

\clearpage
\newpage
\section{\cardextension \ - \cardextensiontitle}

\begin{tikzpicture}
	\card
	\cardstrip
	\cardbanner{banner/white.png}
	\cardicon{banner/coin.png}
	\cardprice{3+}
	\cardtitle{Arzt}
	\cardcontent{Wenn du diese Karte kaufst, darfst du mehr dafür zahlen. Pro \coin{1}, das du zusätzlich zahlst (d. h. \enquote{überzahlst}), darfst du dir die oberste Karte deines Nachziehstapels ansehen und sie entsorgen, ablegen oder auf den Nachziehstapel zurücklegen. Wenn dein Nachziehstapel aufgebraucht ist, mischst du deinen Ablagestapel und machst ihn zum neuen Nachziehstapel. Wenn dann noch immer nicht genügend Karten im Nachziehstapel sind, siehst du dir keine Karten an. Auch wenn du mehr als \coin{1} überzahlst, wird jede oberste Karte vom Nachziehstapel einzeln komplett abgehandelt, bevor du die nächste Karte ziehst.
	
	\smallskip
	
	Wenn du diese Karte in der Aktionsphase spielst, nennst du den Namen einer beliebigen Karte, deckst die obersten drei Karten deines Nachziehstapels auf und entsorgst jede Karte mit dem von dir genannten Namen. Die übrigen Karten legst du in beliebiger Reihenfolge wieder oben auf deinen Nachziehstapel. Du musst keine Karte nennen, die in diesem Spiel verwendet wird. Wenn dein Nachziehstapel keine drei Karten mehr umfasst, deckst du die darin noch enthaltenen Karten auf, mischst dann deinen Ablagestapel (der die bereits aufgedeckten Karten nicht umfasst) und machst ihn zum neuen Nachziehstapel, von dem du die noch fehlenden Karten aufdeckst. Sind dann noch immer nicht genügend Karten aufgedeckt, belässt du es bei den bereits aufgedeckten Karten.}
\end{tikzpicture}
\hspace{-0.6cm}
\begin{tikzpicture}
	\card
	\cardstrip
	\cardbanner{banner/white.png}
	\cardicon{banner/coin.png}
	\cardprice{5}
	\cardtitle{Bäcker}
	\cardcontent{Wenn du diese Karte spielst, ziehst du eine Karte, darfst eine weitere Aktionskarte ausspielen und nimmst dir eine Münze.
	
	\smallskip

	Wird ein Spiel mit dieser Karte gespielt, erhält jeder Spieler zu Beginn des Spiels eine Münze. Das gilt auch für Partien mit dem Schwarzmarkt , bei denen der Bäcker sich im Schwarzmarkt-Stapel befindet.}
\end{tikzpicture}
\hspace{-0.6cm}
\begin{tikzpicture}
	\card
	\cardstrip
	\cardbanner{banner/white.png}
	\cardicon{banner/coin.png}
	\cardprice{4}
	\cardtitle{Berater}
	\cardcontent{Wenn dein Nachziehstapel keine drei Karten mehr umfasst, deckst du die darin noch enthaltenen Karten auf, mischst dann deinen Ablagestapel und machst ihn zum neuen Nachziehstapel, von dem du die noch fehlenden Karten aufdeckst. Sind dann noch immer nicht genügend Karten aufgedeckt, belässt du es bei den bereits aufgedeckten Karten. Unabhängig davon, wie viele Karten du aufgedeckt hast, wählt dein linker Nachbar eine davon aus, die du ablegst. Die verbleibenden Karten nimmst du auf die Hand.}
\end{tikzpicture}
\hspace{-0.6cm}
\begin{tikzpicture}
	\card
	\cardstrip
	\cardbanner{banner/white.png}
	\cardicon{banner/coin.png}
	\cardprice{4+}
	\cardtitle{Herold}
	\cardcontent{\tiny{\begin{Spacing}{1}
	\vspace{1em}
	Wenn du diese Karte kaufst, darfst du mehr dafür zahlen. Wenn du das tust legst du pro 1 Geld, das du überzahlst, eine beliebige Karte deines Ablagestapels oben auf deinen Nachziehstapel. Du darfst dir dafür die Karten deines Ablagestapels ansehen, was normalerweise nicht möglich ist. Allerdings darfst du nicht zuerst deinen Ablagestapel durchsehen, um zu entscheiden, wie viel du überzahlen willst. Sobald du überzahlt hast, musst du die entsprechende Anzahl an Karten in beliebiger Reihenfolge oben auf deinen Nachziehstapel legen, sofern möglich. Wenn du so viel überzahlst, dass du mehr Karten auf deinen Nachziehstapel legen müsstest, als sich in deinem Ablagestapel befinden, legst du einfach alle Karten deines Ablagestapels in beliebiger Reihenfolge auf deinen Nachziehstapel. Falls du den Herold kaufst, ohne zu überzahlen, darfst du deinen Ablagestapel nicht durchsehen.
	
	\smallskip

	Wenn du diese Karte in der Aktionsphase spielst, ziehst du zuerst eine Karte. Du erhältst dann eine weitere Aktion, die du spielen darfst, nachdem du alle Anweisungen dieser Karte (soweit möglich) erfüllt hast. Dann deckst du die oberste Karte deines Nachziehstapels auf. Wenn es sich um eine Aktionskarte handelt, musst du sie spielen. Die Karte zu spielen, verbraucht keine Aktion. Karten, die mehreren Kartentypen angehören, von denen einer Aktion ist (z. B. Große Halle aus Dominion – Die Intrige), sind Aktionskarten. Alle anderen Kartentypen werden auf den Nachziehstapel zurückgelegt, ohne sie zu spielen.
	
	\smallskip

	\emph{Hinweis:} Wenn durch den Herold eine Dauer-Karte (aus Dominion – Seaside) ausgespielt wird, wird der Herold dennoch am Ende der Runde wie gewohnt abgelegt, da er nicht gebraucht wird, um an etwas zu erinnern.
	\end{Spacing}}}
\end{tikzpicture}
\hspace{-0.6cm}
\begin{tikzpicture}
	\card
	\cardstrip
	\cardbanner{banner/white.png}
	\cardicon{banner/coin.png}
	\cardprice{5}
	\cardtitle{\scriptsize{Kaufmannsgilde}}
	\cardcontent{Wenn diese Karte im Spiel ist (d. h. du hast sie in der Aktionsphase gespielt), darfst du eine weitere Karte kaufen und hast \coin{1} zusätzlich zur Verfügung. Jedes Mal, wenn du eine Karte kaufst, nimmst du dir eine Münze (1 Münze, wenn du eine Karte kaufst; 2 Münzen, wenn du zwei Karten kaufst etc.). Denke daran, dass du Münzen nur einsetzen kannst, bevor du Karten kaufst: Du darfst also diese Münze nicht sofort aufwenden. Diese Anweisung ist kumulativ: Wenn du zwei Kaufmannsgilden ausgespielt hast, bringt dir jede Karte, die du kaufst, zwei Münzen ein. Wenn du jedoch eine Kaufmannsgilde mehrfach spielst, aber nur eine im Spiel hast – wie beim Thronsaal (Dominion – Basisspiel) oder Königshof (Dominion – Blütezeit) –, bekommst du beim Kauf einer Karte nur eine Münze.}
\end{tikzpicture}
\hspace{-0.6cm}
\begin{tikzpicture}
	\card
	\cardstrip
	\cardbanner{banner/white.png}
	\cardicon{banner/coin.png}
	\cardprice{2}
	\cardtitle{\scriptsize{Leuchtenmacher}}
	\cardcontent{Du darfst in der Aktionsphase eine weitere Aktionskarte ausspielen. Du darfst in der Kaufphase einen zusätzlichen Kauf tätigen. Nimm dir eine Münze.}
\end{tikzpicture}
\hspace{-0.6cm}
\begin{tikzpicture}
	\card
	\cardstrip
	\cardbanner{banner/gold.png}
	\cardicon{banner/coin.png}
	\cardprice{3+}
	\cardtitle{\footnotesize{Meisterstück}}
	\cardcontent{Dies ist eine Geldkarte mit dem Wert \coin{1}, wie Kupfer. Wenn du sie kaufst, nimmst du dir ein Silber pro \coin{1}, das du überzahlst. Falls du zum Beispiel \coin{6} für das Meisterstück zahlst, erhältst du drei Silber. Das Meisterstück ist eine Geldkarte und wird regeltechnisch auch so behandelt.}
\end{tikzpicture}
\hspace{-0.6cm}
\begin{tikzpicture}
	\card
	\cardstrip
	\cardbanner{banner/white.png}
	\cardicon{banner/coin.png}
	\cardprice{5}
	\cardtitle{Metzger}
	\cardcontent{Zuerst nimmst du dir 2 Münzen. Dann darfst du eine Karte aus deiner Hand entsorgen und eine beliebige Anzahl an Münzen einsetzen (auch 0 Münzen). Da du den Metzger nicht mehr auf der Hand hast, kannst du diese Karte nicht entsorgen. Allerdings kannst du eine andere Metzger-Karte entsorgen. Wenn du eine Karte entsorgt hast, nimmst du dir eine Karte, deren Kosten höchstens der Summe aus den Kosten der entsorgten Karte und der Anzahl der eingesetzten Münzen entsprechen darf. So könntest du beispielsweise ein Anwesen entsorgen und sechs Münzen zahlen, um dir eine Provinz zu nehmen; oder du könntest einen weiteren Metzger entsorgen und null Münzen zahlen, um dir ein Herzogtum zu nehmen. Die Münzen, die du einsetzt, werden in den Vorrat zurückgelegt und dürfen in der Kaufphase nicht mehr benutzt werden, um andere Karten zu kaufen.}
\end{tikzpicture}
\hspace{-0.6cm}
\begin{tikzpicture}
	\card
	\cardstrip
	\cardbanner{banner/white.png}
	\cardicon{banner/coin.png}
	\cardprice{4}
	\cardtitle{Platz}
	\cardcontent{Zuerst ziehst du eine Karte. Du erhältst zwei weitere Aktionen, die du spielen darfst, nachdem du alle Anweisungen dieser Karte (soweit möglich) erfüllt hast. Dann darfst du eine Geldkarte (auch ein Meisterstück) ablegen. Du kannst die Karte ablegen, die du eben gezogen hast, falls es sich um eine Geldkarte handelt. Wenn du eine Geldkarte abgelegt hast, nimmst du dir eine Münze. Karten, die mehreren Kartentypen angehören, von denen einer Geld ist (wie der Harem aus Dominion – Die Intrige), sind Geldkarten.}
\end{tikzpicture}
\hspace{-0.6cm}
\begin{tikzpicture}
	\card
	\cardstrip
	\cardbanner{banner/white.png}
	\cardicon{banner/coin.png}
	\cardprice{2+}
	\cardtitle{Steinmetz}
	\cardcontent{\tiny{\begin{Spacing}{1}
	\vspace{1em}
	Wenn du die Karte Steinmetz kaufst, darfst du mehr dafür bezahlen. Wenn du das tust, nimmst du dir zwei Aktionskarten, die beide jeweils genau so viel kosten, wie du überzahlt hast. Du kannst dir zweimal die gleiche oder zwei verschiedene Karten nehmen. Wenn du zum Beispiel den Steinmetz für \coin{6} kaufst, könntest du dir zwei Herolde nehmen. Die Aktionskarten stammen aus dem Vorrat und werden auf deinen Ablagestapel gelegt. Sollten keine Karten mit den entsprechenden Kosten im Vorrat sein, bekommst du keine. Wenn du mit einem Trank (Dominion – Die Alchemisten) überzahlst, bekommst du Karten im Wert des Tranks . Karten, die mehreren Kartentypen angehören, von denen einer Aktion ist (wie die Große Halle aus Dominion – Die Intrige), sind Aktionskarten. Wenn du beschließt, nicht mehr als die normalen Kosten zu zahlen, bekommst du keine Karten. Es ist nicht möglich, dann Aktionskarten zu nehmen, die 0 Geld kosten. 
	
	\smallskip

	Wenn du diese Karte spielst, entsorgst du eine Karte aus deiner Hand und nimmst dir zwei Karten, die jeweils weniger kosten als die Karte, die du entsorgt hast. Wenn du keine Karte zum Entsorgen mehr auf der Hand hast, bekommst du keine Karten. Du kannst dir zweimal die gleiche oder zwei verschiedene Karten aus dem Vorrat nehmen. Diese werden auf deinen Ablagestapel gelegt. Falls es keine preiswerteren Karten im Vorrat gibt (weil du beispielsweise ein Kupfer entsorgst), bekommst du keine Karten. Wenn es im Vorrat nur noch eine einzige Karte gibt, die preiswerter ist als die entsorgte Karte, nimmst du dir diese. Du nimmst dir die Karten einzeln nacheinander; das kann bei Karten eine Rolle spielen, die beim Nehmen eine Auswirkung haben (wie das Gasthaus aus Dominion – Hinterland).
	\end{Spacing}}}
\end{tikzpicture}
\hspace{-0.6cm}
\begin{tikzpicture}
	\card
	\cardstrip
	\cardbanner{banner/white.png}
	\cardicon{banner/coin.png}
	\cardprice{4}
	\cardtitle{\scriptsize{Steuereintreiber}}
	\cardcontent{Du darfst eine Geldkarte aus deiner Hand entsorgen. Karten, die mehreren Kartentypen angehören, von denen einer Geld ist (wie der Harem aus Dominion – Die Intrige), sind Geldkarten. Wenn du eine Geldkarte entsorgst, muss jeder Mitspieler, der mindestens fünf Karten in der Hand hat, die gleiche Karte aus seiner Hand auf seinen Ablagestapel legen oder seine Handkarten vorzeigen, damit die anderen Spieler sehen, dass er diese Karte nicht auf der Hand hat. Nimm dir für die entsorgte Karte eine Geldkarte, die bis zu \coin{3} mehr kostet als die entsorgte Geldkarte, und lege sie oben auf deinen Nachziehstapel. Wenn dein Nachziehstapel aufgebraucht ist, wird dies die einzige Karte deines Nachziehstapels. Du musst dir keine teurere Geldkarte nehmen, sondern darfst dir auch eine gleich teure oder preiswertere Geldkarte nehmen.}
\end{tikzpicture}
\hspace{-0.6cm}
\begin{tikzpicture}
	\card
	\cardstrip
	\cardbanner{banner/white.png}
	\cardicon{banner/coin.png}
	\cardprice{5}
	\cardtitle{Wahrsager}
	\cardcontent{Das Gold und die Fluchkarten stammen aus dem Vorrat und werden auf den Ablagestapel gelegt. Wenn kein Gold mehr vorhanden ist, erhältst du keins. Sind nicht mehr genügend Fluchkarten da, teilst du sie reihum aus, beginnend mit deinem linken Nachbarn. Jeder Spieler, der eine Fluchkarte erhalten hat, muss eine Karte vom Nachziehstapel ziehen. Falls ein Spieler keine Fluchkarte bekommen hat – weil nicht genügend Fluchkarten da waren oder aus einem anderen Grund –, zieht er keine Karte. Nutzt ein Spieler den Wachturm (Dominion – Blütezeit), um den Fluch zu entsorgen, hat er dennoch eine Fluchkarte erhalten und zieht deshalb eine Karte. Nutzt ein Spieler den Fahrenden Händler (Dominion – Hinterland), um stattdessen ein Silber zu nehmen, hat er keine Fluchkarte erhalten und zieht infolgedessen auch keine Karte.}
\end{tikzpicture}
\hspace{-0.6cm}
\begin{tikzpicture}
	\card
	\cardstrip
	\cardbanner{banner/white.png}
	\cardicon{banner/coin.png}
	\cardprice{5}
	\cardtitle{\footnotesize{Wandergeselle}}
	\cardcontent{Zunächst nennst du eine Karte. Dabei muss es sich nicht um den Namen einer Karte handeln, die in diesem Spiel verwendet wird. Dann deckst du solange Karten vom Nachziehstapel auf, bis du drei Karten aufgedeckt hast, die nicht dem von dir genannten Namen entsprechen. Nimm diese Karten auf die Hand und lege die anderen ab. Sollte dabei der Nachziehstapel aufgebraucht werden, bevor du auf drei entsprechende Karten gestoßen bist, mischst du deinen Ablagestapel und nutzt ihn als neuen Nachziehstapel. Wenn du beim weiteren Aufdecken noch immer nicht auf drei Karten mit einem anderen als dem genannten Namen kommst, beendest du das Aufdecken. Nimm die gefundenen Karten anderen Namens auf die Hand und lege den Rest auf deinen Ablagestapel.}
\end{tikzpicture}
\hspace{-0.6cm}
\begin{tikzpicture}
	\card
	\cardstrip
	\cardbanner{banner/white.png}
	\cardtitle{\scriptsize{Empfohlene 10er Sätze\qquad}}
	\cardcontent{\emph{Die Gilden und Basisspiel:}
	
	\smallskip

	\emph{Kunsthandwerk:} \\ 
	Steinmetz, Berater, Bäcker, Wandergeselle, Kaufmannsgilde / Laboratorium, Keller, Werkstatt, Jahrmarkt, Geldverleiher 

	\smallskip 
	
	\emph{Rechtschaffen und anständig:} \\ 
	Metzger, Bäcker, Leuchtenmacher, Arzt, Wahrsager / Miliz, Dieb, Geldverleiher, Gärten, Dorf 
	
	\smallskip 
	
	\emph{Des Guten zuviel:} \\ 
	Platz, Meisterstück, Leuchtenmacher, Steuereintreiber, Herold / Bibliothek, Umbau, Abenteurer, Markt, Kanzler}
\end{tikzpicture}
\hspace{-0.6cm}
\begin{tikzpicture}
	\card
	\cardstrip
	\cardbanner{banner/white.png}
	\cardtitle{\scriptsize{Empfohlene 10er Sätze\qquad}}
	\cardcontent{\emph{Die Gilden und Die Intrigen:}
	
	\smallskip

	\emph{Nenne diese Karte:} \\ 
	Bäcker, Arzt, Platz, Berater, Meisterstück / Burghof, Wunschbrunnen, Harem, Tribut, Adelige 

	\smallskip 
	
	\emph{Geschäftstricks:} \\ 
	Steinmetz, Herold, Wahrsager, Wandergeselle, Metzger / Große Halle, Adelige, Verschwörer, Maskerade, Kupferschmied 

	\smallskip 
	
	\emph{Entscheidungen, Entscheidungen:} \\ 
	Kaufmannsgilde, Leuchtenmacher, Meisterstück, Steuereintreiber, Metzger / Brücke, Handlanger, Bergwerk, Anbau, Herzog}
\end{tikzpicture}
\hspace{-0.6cm}
\begin{tikzpicture}
	\card
	\cardstrip
	\cardbanner{banner/white.png}
	\cardtitle{\footnotesize{Neue Regeln (1/2)}\qquad}
	\cardcontent{\tiny{\begin{Spacing}{1}
	\vspace{1em}
	\emph{Die Münzen:}\\
	Einige Karten in Die Gilden erlauben den Spielern, sich Münzen zu nehmen. Münzen werden immer vom Vorrat genommen, nicht von anderen Spielern. Diese Münzen können – anders als die Geldkarten – über den Zug, in dem ein Spieler sie erhält, hinaus aufbewahrt werden. In der Kaufphase eines Spielers kann dieser, bevor er Karten kauft, eine beliebige Anzahl an Münzen einsetzen; jede eingesetzte Münze bringt dem Spieler +1 Geld. Eingesetzte Münzen werden in den Vorrat zurückgelegt. Die Anzahl der Münzen im Vorrat ist nicht begrenzt: Sollten einmal nicht ausreichend Münzen vorhanden sein, verwenden die Spieler einen beliebigen Ersatz. Zwar handelt es sich um die gleichen Marker / Münzen wie in Seaside und Blütezeit, allerdings können die durch Anweisungen auf den Karten aus Die Gilden erworbenen Münzen ausschließlich in der Kaufphase ausgegeben werden. Sie dürfen weder zum Kauf einer Karte mithilfe der Karte Schwarzmarkt benutzt werden, noch dürfen sie auf ein Piratenschiff-Tableau (Dominion – Seaside) oder die Handelsroute (Dominion – Blütezeit) gelegt werden.
	\end{Spacing}}}
\end{tikzpicture}
\hspace{-0.6cm}
\begin{tikzpicture}
	\card
	\cardstrip
	\cardbanner{banner/white.png}
	\cardtitle{\footnotesize{Neue Regeln (2/2)}\qquad}
	\cardcontent{\tiny{\begin{Spacing}{1}
	\vspace{1em}
	\emph{Überzahlen:}\\
	Für einige Karten in Die Gilden darf man mehr zahlen, als sie eigentlich kosten. Bei diesen Karten steht hinter den Kosten ein \enquote{+}, z. B. 2+ Geld. Wenn ein Spieler vor dem Kauf einen beliebigen zusätzlichen Betrag für eine solche Karte zahlt, tritt ein auf der jeweiligen Karte genannter Effekt ein, der von der Höhe der Überzahlung abhängt. Überzahlt werden kann mit allen Geldwerten, mit denen Karten gekauft werden: Geldkarten, Münzen und Geldwerte auf Aktionskarten. Die Karten Trank (Dominion – Die Alchemisten) können ebenfalls zum Überzahlen verwendet werden (allerdings ist das nicht immer sinnvoll). Das zum Überzahlen verwendete Geld ist ausgegeben und kann dann nicht mehr eingesetzt werden. Geldkarten werden abgelegt und Münzen kommen zurück in den Vorrat. Spieler können eine Karte ausschließlich beim Kauf überzahlen, nicht wenn sie diese auf eine andere Weise erhalten. Das \enquote{+} ist lediglich eine Erinnerung: Für eine Karte, die hinter den Kosten das Symbol \enquote{+} aufweist, gelten nach wie vor für alle Zwecke die normalen Kosten. Grundsätzlich gibt es keine Wechselwirkungen zwischen Karten, welche die Kartenkosten reduzieren – wie die Brücke (Dominion – Die Intrige) oder die Fernstraße (Dominion – Hinterland) – und dem \enquote{Überzahlen}. 
	
	\smallskip

	\emph{Zwei Dinge passieren zur gleichen Zeit:} \\
	Wenn einem Spieler zwei Dinge gleichzeitig passieren, entscheidet er selbst, in welcher Reihenfolge sie eintreten. Wenn unterschiedlichen Spielern zwei Dinge gleichzeitig passieren, werden diese reihum in Spielreihenfolge abgehandelt, beginnend mit dem Spieler, der gerade an der Reihe ist.
	\end{Spacing}}}
\end{tikzpicture}
\hspace{0.6cm}

	    % Basic settings for this card set
\renewcommand{\cardcolor}{adventures}
\renewcommand{\cardextension}{Erweiterung VIII}
\renewcommand{\cardextensiontitle}{Abenteuer}

\clearpage
\newpage
\section{\cardextension \ - \cardextensiontitle}

\begin{tikzpicture}
	\card
	\cardstrip
	\cardbanner{banner/white.png}
	\cardicon{banner/coin.png}
	\cardprice{}
	\cardtitle{}
	\cardcontent{}
\end{tikzpicture}

\hspace{-1cm}
\begin{tikzpicture}
\card
\cardstrip
\cardbanner{banner/white.png}
\cardicon{}
\cardprice{}
\cardtitle{\scriptsize{Empfohlene 10er Sätze\qquad}}
\cardcontent{\emph{Name:}
	\\
	Karten ...
	\\
	\smallskip
	\\
	\emph{Name:}
	\\
	Karten ...
	\\
	\smallskip
	\\
	\emph{Name:}
	\\
	Karten ...
	\\
	\smallskip
	\\
	\emph{Name:}
	\\
	Karten ...
	\\
	\smallskip
	\\
	\emph{Name:}
	\\
	Karten ...
	\\
	\smallskip
	\\
	\emph{Name:}
	\\
	Karten ...
	\\
}
\end{tikzpicture}
\hspace{1cm}
	    % Basic settings for this card set
\renewcommand{\cardcolor}{empires}
\renewcommand{\cardextension}{Erweiterung IX}
\renewcommand{\cardextensiontitle}{Empires}

\clearpage
\newpage
\section{\cardextension \ - \cardextensiontitle}

\begin{tikzpicture}
	\card
	\cardstrip
	\cardbanner{banner/white.png}
	\cardicon{banner/coin.png}
	\cardprice{2}
	\cardtitle{Feldlager}
	\cardcontent{}
\end{tikzpicture}
\hspace{-1cm}
\begin{tikzpicture}
	\card
	\cardstrip
	\cardbanner{banner/white.png}
	\cardicon{banner/coin.png}
	\cardprice{2}
	\cardtitle{Patrizier}
	\cardcontent{}
\end{tikzpicture}
\hspace{-1cm}
\begin{tikzpicture}
	\card
	\cardstrip
	\cardbanner{banner/white.png}
	\cardicon{banner/coin.png}
	\cardprice{2}
	\cardtitle{Siedler}
	\cardcontent{}
\end{tikzpicture}
\hspace{-1cm}
\begin{tikzpicture}
	\card
	\cardstrip
	\cardbanner{banner/white.png}
	\cardicon{banner/coin.png}
	\cardprice{3}
	\cardtitle{Bauernmarkt}
	\cardcontent{}
\end{tikzpicture}
\hspace{-1cm}
\begin{tikzpicture}
	\card
	\cardstrip
	\cardbanner{banner/white.png}
	\cardicon{banner/coin.png}
	\cardprice{3}
	\cardtitle{Gladiator}
	\cardcontent{}
\end{tikzpicture}
\hspace{-1cm}
\begin{tikzpicture}
	\card
	\cardstrip
	\cardbanner{banner/white.png}
	\cardicon{banner/coin.png}
	\cardprice{3}
	\cardtitle{Katapult}
	\cardcontent{}
\end{tikzpicture}
\hspace{-1cm}
\begin{tikzpicture}
	\card
	\cardstrip
	\cardbanner{banner/white.png}
	\cardicon{banner/coin.png}
	\cardprice{3}
	\cardtitle{Wagenrennen}
	\cardcontent{}
\end{tikzpicture}
\hspace{-1cm}
\begin{tikzpicture}
	\card
	\cardstrip
	\cardbanner{banner/orange.png}
	\cardicon{banner/coin.png}
	\cardprice{3}
	\cardtitle{Zauberin}
	\cardcontent{}
\end{tikzpicture}
\hspace{-1cm}
\begin{tikzpicture}
	\card
	\cardstrip
	\cardbanner{banner/gold.png}
	\cardicon{banner/coin.png}
	\cardprice{4}
	\cardtitle{Felsen}
	\cardcontent{}
\end{tikzpicture}
\hspace{-1cm}
\begin{tikzpicture}
	\card
	\cardstrip
	\cardbanner{banner/white.png}
	\cardicon{banner/coin.png}
	\cardprice{4}
	\cardtitle{Opfer}
	\cardcontent{}
\end{tikzpicture}
\hspace{-1cm}
\begin{tikzpicture}
	\card
	\cardstrip
	\cardbanner{banner/white.png}
	\cardicon{banner/coin.png}
	\cardprice{4}
	\cardtitle{Tempel}
	\cardcontent{}
\end{tikzpicture}
\hspace{-1cm}
\begin{tikzpicture}
	\card
	\cardstrip
	\cardbanner{banner/white.png}
	\cardicon{banner/coin.png}
	\cardprice{4}
	\cardtitle{Villa}
	\cardcontent{}
\end{tikzpicture}
\hspace{-1cm}
\begin{tikzpicture}
	\card
	\cardstrip
	\cardbanner{banner/orange.png}
	\cardicon{banner/coin.png}
	\cardprice{5}
	\cardtitle{Archiv}
	\cardcontent{}
\end{tikzpicture}
\hspace{-1cm}
\begin{tikzpicture}
	\card
	\cardstrip
	\cardbanner{banner/gold.png}
	\cardicon{banner/coin.png}
	\cardprice{5}
	\cardtitle{Diebesgut}
	\cardcontent{}
\end{tikzpicture}
\hspace{-1cm}
\begin{tikzpicture}
	\card
	\cardstrip
	\cardbanner{banner/white.png}
	\cardicon{banner/coin.png}
	\cardprice{5}
	\cardtitle{Emsiges Dorf}
	\cardcontent{}
\end{tikzpicture}
\hspace{-1cm}
\begin{tikzpicture}
	\card
	\cardstrip
	\cardbanner{banner/white.png}
	\cardicon{banner/coin.png}
	\cardprice{5}
	\cardtitle{Forum}
	\cardcontent{}
\end{tikzpicture}
\hspace{-1cm}
\begin{tikzpicture}
	\card
	\cardstrip
	\cardbanner{banner/white.png}
	\cardicon{banner/coin.png}
	\cardprice{5}
	\cardtitle{Gärtnerin}
	\cardcontent{}
\end{tikzpicture}
\hspace{-1cm}
\begin{tikzpicture}
	\card
	\cardstrip
	\cardbanner{banner/white.png}
	\cardicon{banner/coin.png}
	\cardprice{5}
	\cardtitle{Handelsplatz}
	\cardcontent{}
\end{tikzpicture}
\hspace{-1cm}
\begin{tikzpicture}
	\card
	\cardstrip
	\cardbanner{banner/whitegold.png}
	\cardicon{banner/coin.png}
	\cardprice{5}
	\cardtitle{Krone}
	\cardcontent{}
\end{tikzpicture}
\hspace{-1cm}
\begin{tikzpicture}
	\card
	\cardstrip
	\cardbanner{banner/white.png}
	\cardicon{banner/coin.png}
	\cardprice{5}
	\cardtitle{Legionär}
	\cardcontent{}
\end{tikzpicture}
\hspace{-1cm}
\begin{tikzpicture}
	\card
	\cardstrip
	\cardbanner{banner/gold.png}
	\cardicon{banner/coin.png}
	\cardprice{5}
	\cardtitle{Vermögen}
	\cardcontent{}
\end{tikzpicture}
\hspace{-1cm}
\begin{tikzpicture}
	\card
	\cardstrip
	\cardbanner{banner/white.png}
	\cardicon{banner/coin.png}
	\cardprice{5}
	\cardtitle{Wilde Jagd}
	\cardcontent{}
\end{tikzpicture}
\hspace{-1cm}
\begin{tikzpicture}
	\card
	\cardstrip
	\cardbanner{banner/gold.png}
	\cardicon{banner/coin.png}
	\cardprice{5}
	\cardtitle{Zauber}
	\cardcontent{}
\end{tikzpicture}
\hspace{-1cm}
\begin{tikzpicture}
	\card
	\cardstrip
	\cardbanner{banner/white.png}
	\cardicon{banner/hex.png}
	\cardprice{\textcolor{white}{4}}
	\cardtitle{Ingenieurin}
	\cardcontent{}
\end{tikzpicture}
\hspace{-1cm}
\begin{tikzpicture}
	\card
	\cardstrip
	\cardbanner{banner/white.png}
	\cardicon{banner/hex.png}
	\cardprice{\textcolor{white}{8}}
	\cardtitle{\scriptsize{Königlicher\,Schmied}}
	\cardcontent{}
\end{tikzpicture}
\hspace{-1cm}
\begin{tikzpicture}
	\card
	\cardstrip
	\cardbanner{banner/white.png}
	\cardicon{banner/hex.png}
	\cardprice{\textcolor{white}{8}}
	\cardtitle{Lehnsherr}
	\cardcontent{}
\end{tikzpicture}
\hspace{-1cm}
\begin{tikzpicture}
	\card
	\cardstrip
	\cardbanner{banner/white.png}
	\cardicon{banner/hex.png}
	\cardprice{\textcolor{white}{8}}
	\cardtitle{Stadtviertel}
	\cardcontent{}
\end{tikzpicture}
\hspace{-1cm}
\begin{tikzpicture}
	\card
	\cardstrip
	\cardbanner{banner/white.png}
	\cardicon{banner/coin.png}
	\cardprice{8}
	\cardiconaddition{banner/hex.png}
	\cardpriceaddition{\textcolor{white}{8}}
	\cardtitle{\quad Reichtum}
	\cardcontent{}
\end{tikzpicture}
\hspace{-1cm}
\begin{tikzpicture}
	\card
	\cardstrip
	\cardbanner{banner/green.png}
	\cardtitle{Schlösser (1/2)\quad}
	\cardcontent{}
\end{tikzpicture}
\hspace{-1cm}
\begin{tikzpicture}
	\card
	\cardstrip
	\cardbanner{banner/green.png}
	\cardtitle{Schlösser (2/2)\quad}
	\cardcontent{}
\end{tikzpicture}
\hspace{-1cm}
\begin{tikzpicture}
	\card
	\cardstrip
	\cardbanner{banner/white.png}
	\cardtitle{Katapult/Felsen\qquad}
	\cardcontent{}
\end{tikzpicture}
\hspace{-1cm}
\begin{tikzpicture}
	\card
	\cardstrip
	\cardbanner{banner/white.png}
	\cardtitle{\scriptsize{Gladiator/Reichtum}\qquad}
	\cardcontent{}
\end{tikzpicture}
\hspace{-1cm}
\begin{tikzpicture}
	\card
	\cardstrip
	\cardbanner{banner/white.png}
	\cardtitle{\scriptsize{Siedler/Emsiges Dorf}\qquad}
	\cardcontent{}
\end{tikzpicture}
\hspace{-1cm}
\begin{tikzpicture}
	\card
	\cardstrip
	\cardbanner{banner/white.png}
	\cardtitle{\scriptsize{Patrizier/Handelsplatz}\qquad}
	\cardcontent{}
\end{tikzpicture}
\hspace{-1cm}
\begin{tikzpicture}
	\card
	\cardstrip
	\cardbanner{banner/white.png}
	\cardtitle{\scriptsize{Feldlager/Diebesgut}\qquad}
	\cardcontent{}
\end{tikzpicture}
\hspace{-1cm}
\begin{tikzpicture}
	\card
	\cardstrip
	\cardbanner{banner/white.png}
	\cardtitle{Ereignisse (1/4)\quad}
	\cardcontent{}
\end{tikzpicture}
\hspace{-1cm}
\begin{tikzpicture}
	\card
	\cardstrip
	\cardbanner{banner/white.png}
	\cardtitle{Ereignisse (2/4)\quad}
	\cardcontent{}
\end{tikzpicture}
\hspace{-1cm}
\begin{tikzpicture}
	\card
	\cardstrip
	\cardbanner{banner/white.png}
	\cardtitle{Ereignisse (3/4)\quad}
	\cardcontent{}
\end{tikzpicture}
\hspace{-1cm}
\begin{tikzpicture}
	\card
	\cardstrip
	\cardbanner{banner/white.png}
	\cardtitle{Ereignisse (4/4)\quad}
	\cardcontent{}
\end{tikzpicture}
\hspace{-1cm}
\begin{tikzpicture}
	\card
	\cardstrip
	\cardbanner{banner/green.png}
	\cardtitle{Landmarken (1/7)\quad}
	\cardcontent{}
\end{tikzpicture}
\hspace{-1cm}
\begin{tikzpicture}
	\card
	\cardstrip
	\cardbanner{banner/green.png}
	\cardtitle{Landmarken (2/7)\quad}
	\cardcontent{}
\end{tikzpicture}
\hspace{-1cm}
\begin{tikzpicture}
	\card
	\cardstrip
	\cardbanner{banner/green.png}
	\cardtitle{Landmarken (3/7)\quad}
	\cardcontent{}
\end{tikzpicture}
\hspace{-1cm}
\begin{tikzpicture}
	\card
	\cardstrip
	\cardbanner{banner/green.png}
	\cardtitle{Landmarken (4/7)\quad}
	\cardcontent{}
\end{tikzpicture}
\hspace{-1cm}
\begin{tikzpicture}
	\card
	\cardstrip
	\cardbanner{banner/green.png}
	\cardtitle{Landmarken (5/7)\quad}
	\cardcontent{}
\end{tikzpicture}
\hspace{-1cm}
\begin{tikzpicture}
	\card
	\cardstrip
	\cardbanner{banner/green.png}
	\cardtitle{Landmarken (6/7)\quad}
	\cardcontent{}
\end{tikzpicture}
\hspace{-1cm}
\begin{tikzpicture}
	\card
	\cardstrip
	\cardbanner{banner/green.png}
	\cardtitle{Landmarken (7/7)\quad}
	\cardcontent{}
\end{tikzpicture}
\hspace{-1cm}
\begin{tikzpicture}
	\card
	\cardstrip
	\cardbanner{banner/white.png}
	\cardtitle{\scriptsize{Empfohlene 10er Sätze\qquad}}
	\cardcontent{\emph{}
	\\
	
	\\
	\smallskip
	\\
	\emph{}
	\\
	
	\\
	\smallskip
	\\
	\emph{}
	\\
	
	\\
	\smallskip
	\\
	\emph{}
	\\
	
	\\
	\smallskip
	\\
	\emph{}
	\\
	
	\\
	\smallskip
	\\
	\emph{}
	\\
	
	\\}
\end{tikzpicture}
\hspace{-1cm}
\begin{tikzpicture}
	\card
	\cardstrip
	\cardbanner{banner/white.png}
	\cardtitle{\scriptsize{Empfohlene 10er Sätze\qquad}}
	\cardcontent{\emph{}
	\\
	
	\\
	\smallskip
	\\
	\emph{}
	\\
	
	\\
	\smallskip
	\\
	\emph{}
	\\
	
	\\
	\smallskip
	\\
	\emph{}
	\\
	
	\\
	\smallskip
	\\
	\emph{}
	\\
	
	\\
	\smallskip
	\\
	\emph{}
	\\
	
	\\}
\end{tikzpicture}
\hspace{-1cm}
\begin{tikzpicture}
	\card
	\cardstrip
	\cardbanner{banner/white.png}
	\cardtitle{\scriptsize{Empfohlene 10er Sätze\qquad}}
	\cardcontent{\emph{}
	\\
	
	\\
	\smallskip
	\\
	\emph{}
	\\
	
	\\
	\smallskip
	\\
	\emph{}
	\\
	
	\\
	\smallskip
	\\
	\emph{}
	\\
	
	\\
	\smallskip
	\\
	\emph{}
	\\
	
	\\
	\smallskip
	\\
	\emph{}
	\\
	
	\\}
\end{tikzpicture}
\hspace{-1cm}
\begin{tikzpicture}
	\card
	\cardstrip
	\cardbanner{banner/white.png}
	\cardtitle{\scriptsize{Empfohlene 10er Sätze\qquad}}
	\cardcontent{\emph{}
	\\
	
	\\
	\smallskip
	\\
	\emph{}
	\\
	
	\\
	\smallskip
	\\
	\emph{}
	\\
	
	\\
	\smallskip
	\\
	\emph{}
	\\
	
	\\
	\smallskip
	\\
	\emph{}
	\\
	
	\\
	\smallskip
	\\
	\emph{}
	\\
	
	\\}
\end{tikzpicture}
\hspace{1cm}
	    % Basic settings for this card set
\renewcommand{\cardcolor}{nocturne}
\renewcommand{\cardextension}{Erweiterung X}
\renewcommand{\cardextensiontitle}{Nocturne}

\clearpage
\newpage
\section{\cardextension \ - \cardextensiontitle}

\begin{tikzpicture}
	\card
	\cardstrip
	\cardbanner{banner/white.png}
	\cardicon{banner/coin.png}
	\cardprice{}
	\cardtitle{}
	\cardcontent{}
\end{tikzpicture}

\hspace{-1cm}
\begin{tikzpicture}
\card
\cardstrip
\cardbanner{banner/white.png}
\cardicon{}
\cardprice{}
\cardtitle{\scriptsize{Empfohlene 10er Sätze\qquad}}
\cardcontent{\emph{Name:}
	\\
	Karten ...
	\\
	\smallskip
	\\
	\emph{Name:}
	\\
	Karten ...
	\\
	\smallskip
	\\
	\emph{Name:}
	\\
	Karten ...
	\\
	\smallskip
	\\
	\emph{Name:}
	\\
	Karten ...
	\\
	\smallskip
	\\
	\emph{Name:}
	\\
	Karten ...
	\\
	\smallskip
	\\
	\emph{Name:}
	\\
	Karten ...
	\\
}
\end{tikzpicture}
\hspace{1cm}
	    % Basic settings for this card set
\renewcommand{\cardcolor}{promo}
\renewcommand{\cardextension}{Promokarte}
\renewcommand{\cardextensiontitle}{}

\clearpage
\newpage
\section{\cardextension}

\begin{tikzpicture}
	\card
	\cardstrip
	\cardbanner{banner/white.png}
	\cardicon{banner/coin.png}
	\cardprice{4}
	\cardtitle{Gesandter}
	\cardcontent{Die aufgedeckten Karten legst du zunächst offen vor dir aus. Kannst du auch nach dem Mischen deines Ablagestapels nur weniger als 5 Karten aufdecken, deckst du nur so viele auf, wie möglich. Dann wählt der Spieler links von dir eine dieser offenen Karten. Die gewählte Karte legst du auf deinen Ablagestapel, die restlichen Karten nimmst du auf die Hand.}
\end{tikzpicture}
\hspace{-0.6cm}
\begin{tikzpicture}
	\card
	\cardstrip
	\cardbanner{banner/white.png}
	\cardicon{banner/coin.png}
	\cardprice{3}
	\cardtitle{\footnotesize{Schwarzmarkt}}
	\cardcontent{\miniscule{\begin{Spacing}{1}
	\vspace{1em}
	Vor dem Spiel mit der Aktionskarte Schwarzmarkt muss der dazugehörige Schwarzmarkt-Stapel zusammengestellt werden. Die Spieler einigen sich welche Karten verwendet werden sollen.

	\medskip

	Für die Zusammenstellung ist Folgendes zu beachten:
	\begin{itemize}
	\item Der Schwarzmarkt-Stapel muss aus mindestens 15 Karten bestehen.
	\item Es dürfen nur Königreichkarten verwendet werden, die nicht im Vorrat sind.
	\item Jede Karte darf nur einmal verwendet werden.
	\end{itemize}
	Die Spieler dürfen sich die Karten vor dem Spiel ansehen. Dann werden die verwendeten Karten gemischt und verdeckt neben dem Vorrat bereit gelegt. Der Schwarzmarkt-Stapel ist nicht Teil des Vorrats. Er wird weder für die Spielende-Bedingung beachtet, noch für andere Zwecke, die auf den Vorrat Bezug nehmen, z.B. Aktionskarten, die erlauben eine Karte zu nehmen (Werkstatt).

	\medskip

	Spielst du den Schwarzmarkt aus, musst du zunächst die obersten 3 Karten vom Schwarzmarkt-Stapel aufdecken. Nun darfst du eine dieser Karten kaufen. Es handelt sich hierbei um einen Kauf in der Aktionsphase, d. h. du darfst sowohl Geldkarten als auch virtuelles Geld verwenden. Der Kauf läuft also in gleicher Weise ab, wie in der Kaufphase. (Der einzige Unterschied ist, dass du nicht wie üblich eine Karte aus dem Vorrat kaufst, sondern eine der 3 aufgedeckten Karten vom Schwarzmarktstapel.) Die nicht gekauften Karten legst du in beliebiger Reihenfolge verdeckt unter den Schwarzmarkt-Stapel zurück. Du musst deinen Mitspielern nicht zeigen in welcher Reihenfolge du die Karten zurücklegst. Die ausgelspielten Geldkarten lässt du bis zur Aufräumphase vor dir liegen. Dieser Kauf verbraucht nicht deinen freien Kauf, du darfst also in der Kaufphase (mindestens) eine weitere Karte kaufen. Nicht verwendetes virtuelles Geld und auch überzählige Münzen ausgespielter Geldkarten stehen dir in der Kaufphase zur Verfügung.

	\medskip

	Wenn du den Schwarzmarkt ausspielst, aber keine Karte kaufen willst oder kannst (z.B. weil der Schwarzmarktstapel leer ist oder du nicht genügend Geld hast), erhältst du trotzdem +\coin{2} für die Kaufphase.
	\end{Spacing}}}
\end{tikzpicture}
\hspace{-0.6cm}
\begin{tikzpicture}
	\card
	\cardstrip
	\cardbanner{banner/gold.png}
	\cardicon{banner/coin.png}
	\cardprice{5}
	\cardtitle{Geldversteck}
	\cardcontent{Das Geldversteck ist eine Geldkarte mit dem Wert \coin{2}, wie ein Silber. In der Aufräumphase wird das Geldversteck wie üblich abgelegt. Immer wenn ein Spieler seinen Ablagestapel mischt, führt er folgende Schritte aus:
	\begin{noindlist}
	\item Er sucht zunächst alle Geldversteck-Karten aus seinem Ablagestapel und legt sie beiseite.
	\item Dann mischt er die verbliebenen Karten des Ablagestapels.
	\item Nun sortiert er die Karten Geldversteck wieder an beliebigen Stellen seiner Wahl (auch ganz oben oder unten) in den Stapel ein. Dabei darf er Karten des Stapels abzählen, aber nicht ansehen.
	\item Zuletzt legt er den Stapel als neuen Nachziehstapel bereit.
	\end{noindlist}}
\end{tikzpicture}
\hspace{-0.6cm}
\begin{tikzpicture}
	\card
	\cardstrip
	\cardbanner{banner/white.png}
	\cardicon{banner/coin.png}
	\cardprice{4}
	\cardtitle{Carcasonne}
	\cardcontent{\emph{Errata:} Der Kartentext ist falsch, es sollte \enquote{Wenn du zu Beginn deiner Aufräumphase nicht mehr als eine weitere Aktionskarte im Spiel hast\dots} statt \enquote{Wenn du zu Beginn deiner Aufräumphase nur noch eine weitere Aktionskarte im Spiel hast\dots} heißen.

	\medskip

	Zuerst ziehst du immer eine Karte nach und erhältst +2 Aktionen. Wenn du zu Beginn deiner Aufräumphase die Karte Carcassonne und nicht mehr als eine weitere Aktionskarte im Spiel hast, darfst du dich entscheiden, Carcassonne oben auf deinen Nachziehstapel zu legen oder wie üblich auf den Ablagestapel zu legen. Hast du die Karte Carcassonne genau zweimal im Spiel und ansonsten keine weitere Aktionskarte im Spiel, darfst du eine oder beide Karten Carcassonne auf deinen Nachziehstapel legen.}
\end{tikzpicture}
\hspace{-0.6cm}
\begin{tikzpicture}
	\card
	\cardstrip
	\cardbanner{banner/white.png}
	\cardicon{banner/coin.png}
	\cardprice{5}
	\cardtitle{Gouverneur}
	\cardcontent{Zuerst erhältst du +1 Aktion. Dann wählst du eine der folgenden Optionen:
	\begin{noindlist}
	\item Du ziehst 3 Karten und jeder andere Spieler zieht 1 Karte.
	\item Du nimmst dir 1 Gold und jeder andere Spieler nimmt sich 1 Silber.
	\item Du darfst 1 Karte aus deiner Hand entsorgen und nimmst dir 1 Karte, die genau
	\coin{2} mehr kostet als die entsorgte Karte und jeder andere Spieler darf 1 Karte entsorgen, die genau \coin{1} mehr kostet als die entsorgte Karte.
	\end{noindlist}
	Geh nach der Reihenfolge, beginnend mit dir selbst. Die Karten werden vom Vorrat genommen und auf den Ablagestapel gelegt. Sind im Vorrat keine Karten mehr übrig, können keine entsprechenden Karten genommen werden. Wählst du z.B. die zweite Option und es ist nur noch 1 Silber im Vorrat, bekommt es der Spieler links von dir und die anderen Spieler erhalten nichts. Bei der dritten Option nimmst du dir nur dann 1 Karte, wenn du vorher 1 Karte entsorgt hast und wenn 1 Karte mit den genauen geforderten Kosten im Vorrat verfügbar ist. Wenn du 1 Karte entsorgst, musst du dir 1 Karte nehmen, sofern möglich. Du kannst den ausgespielten Gouverneur nicht entsorgen, da er sich nicht mehr auf deiner Hand befindet. Du kannst aber einen anderen Gouverneur aus deiner Hand entsorgen.}
\end{tikzpicture}
\hspace{-0.6cm}
\begin{tikzpicture}
	\card
	\cardstrip
	\cardbanner{banner/white.png}
	\cardicon{banner/coin.png}
	\cardprice{8}
	\cardtitle{Prinz}
	\cardcontent{\tiny{\begin{Spacing}{1}
	\vspace{1em}
	Wenn du dich entscheidest, einen Prinzen zu spielen, legst du ihn sofort zur Seite. Er befindet sich damit allerdings nicht im Spiel. Anschließend wählst du eine Aktionskarte von deiner Hand, die zu diesem Zeitpunkt maximal \coin{4} kostet und legst diese ebenfalls zur Seite.

	\medskip

	Auch Karten mit Kosten von \coin{0}, wie die Preiskarten aus Reiche Ernte sowie Karten mit Kosten von \coin{X} + aus Die Gilden dürfen mit Hilfe des Prinzen zur Seite gelegt werden. Karten, deren Kosten einen Trank enthalten, dürfen dagegen nicht zur Seite gelegt werden.

	\medskip

	Zu Beginn eines Zuges musst du die zur Seite gelegte Aktionskarte spielen. Sie verbraucht dabei nicht deine freie Aktion. Sobald du die Aktionskarte ablegen musst, legst du sie stattdessen wieder zur Seite. Wenn du die Aktionskarte während deines Spielzugs aus dem Spielbereich entfernst (z.B. sie entsorgen musst) und sie dementsprechend in der Aufräumphase nicht mehr im Spiel ist, darfst du die Karte nicht wieder zur Seite legen. Die Wirkung des Prinzen wird sofort aufgehoben.

	\medskip

	Wenn du zu Beginn deines Zuges mehrere Karten spielen musst (z.B. Aktionskarten durch mehrere Prinzen oder Dauerkarten aus Seaside), darfst du selbst entscheiden, in welcher Reihenfolge du sie ausspielst. Die Karte \emph{PRINZ} muss zur Seite gelegt werden, damit sie einen Effekt hat. Den Prinzen zum Beispiel auf einen Thronsaal zu spielen erlaubt dir nicht, zwei Karten zur Seite zu legen, da du den Prinzen nur einmal zur Seite legen kannst. Alle zur Seite gelegten Prinzen und Aktionskarten gehören zum Kartensatz eines Spielers.
	\end{Spacing}}}
\end{tikzpicture}
\hspace{-0.6cm}
\begin{tikzpicture}
	\card
	\cardstrip
	\cardbanner{banner/white.png}
	\cardtitle{Sauna/Eisloch\qquad}
	\cardcontent{Spielvorbereitung: Legt auf diese Karte 5 Eisloch und oben darauf 5 Sauna.

	\medskip

	Es darf immer nur die oberste Karte des Stapels genommen oder gekauft werden.}
\end{tikzpicture}
\hspace{-0.6cm}
\begin{tikzpicture}
	\card
	\cardstrip
	\cardbanner{banner/white.png}
	\cardicon{banner/coin.png}
	\cardprice{4}
	\cardtitle{Sauna}
	\cardcontent{\tiny{\begin{Spacing}{1}
	\vspace{1em}
	\emph{Siehe auch die Hinweise zur Karte Eisloch!}

	\medskip

	Wenn du die Sauna ausspielst, ziehst du zuerst eine Karte und bekommst +1 Aktion. Du kannst dann sofort ein Eisloch aus deiner Hand ausspielen. Das verbraucht keine deiner Aktionen, einschließlich der Aktion, die die Sauna gewährt. Du darfst ein Eisloch auf diese Weise nur direkt nach dem Ausspielen der Sauna spielen, nicht, wenn du zwischendurch eine andere Aktionskarte ausgespielt hast, selbst wenn du eine Sauna im Spiel hast.

	\medskip

	Solange die Sauna im Spiel ist, darfst du jedes Mal, wenn du ein Silber ausspielst, eine Karte aus deiner Hand entsorgen. Wenn du das gleiche Silber mehrmals spielst, wie z.B. mit dem Falschgeld (Dominion - Dark Ages) oder der Krone (Dominion - Empires), darfst du jedes Mal eine Karte entsorgen, wenn du das Silber spielst.

	\medskip

	Wenn du ein Silber ausspielst, kannst du dir jedes Mal überlegen, ob du eine Karte entsorgen möchtest, du musst diese Entscheidung nicht einmal für den gesamten Zug treffen. Wenn du mehrere Saunen im Spiel hast und ein Silber ausspielst, kannst du für jede Sauna, die du im Spiel hast, eine Karte aus deiner Hand entsorgen. Du kannst dich jedes Mal immer noch dazu entschließen, keine Karte zu entsorgen.

	\medskip

	Wenn die Sauna das Spiel verlässt, weil sie zum Beispiel mit der Prozession (Dominion - Dark Ages) entsorgt wurde, kann ihr Effekt nicht mehr genutzt werden.
	\end{Spacing}}}
\end{tikzpicture}
\hspace{-0.6cm}
\begin{tikzpicture}
\card
	\cardstrip
	\cardbanner{banner/white.png}
	\cardicon{banner/coin.png}
	\cardprice{5}
	\cardtitle{Eisloch}
	\cardcontent{Wenn du das Eisloch ausspielst, ziehst du zuerst 3 Karten. Du kannst dann sofort eine Sauna aus deiner Hand ausspielen. Das verbraucht keine deiner Aktionen, und du erhältst trotzdem die +1 Aktion der Sauna, wenn du sie auf diese Weise spielst.

	\medskip

	Du darfst eine Sauna auf diese Weise nur direkt nach dem Ausspielen des Eislochs spielen, nicht, wenn du zwischendurch eine andere Aktionskarte ausgespielt hast, selbst wenn du ein Eisloch im Spiel hast.

	\medskip

	\emph{Folgendes gilt sowohl für das Eisloch als auch für die Sauna:}

	\medskip

	Du kannst die Sauna und das Eisloch durch die Effekte der jeweils anderen Karte abwechselnd spielen, wobei du nur die Aktion für die erste ausgespielte Karte verbrauchst. Du kannst damit fortfahren, bis du nach dem Ausspielen der einen Karte die entsprechende andere Karte nicht mehr auf der Hand hast.

	\medskip

	Wenn du eine Sauna ausspielst, kannst du nicht sofort eine weitere Sauna aus deiner Hand ausspielen, ohne eine Aktion zu verbrauchen. Das gleiche gilt für das Ausspielen eines Eislochs nach einem anderen Eisloch.}
\end{tikzpicture}
\hspace{-0.6cm}
\begin{tikzpicture}
	\card
	\cardstrip
	\cardbanner{banner/white.png}
	\cardtitle{Ereignisse\qquad}
	\cardcontent{\tiny{\begin{Spacing}{1}
	\vspace{1em}
	\emph{Einladung:} Wenn du das Ereignis kaufst, nimmst du dir vom Vorrat eine Aktionskarte, die bis zu \coin{4} kostet und legst sie offen zur Seite. Wenn du sie beiseite gelegt hast, dann spielst du die Aktionskarte zu Beginn des nächsten Zuges aus. Das Ausspielen verbraucht nicht deine Standardaktion für den Zug. Um dich daran zu erinnern, dass du die Karte in deinem nächsten Zug ausspielst, kannst du sie seitwärts oder diagonal drehen, und sie dann richtig herum drehen, sobald du sie ausspielst.

	\medskip

	Wenn du die Aktionskarte bewegst, nachdem du sie genommen, aber bevor du sie zur Seite gelegt hast (z.B. indem du sie mit dem Wachturm (Dominion – Blütezeit) auf den Nachziehstapel legst), dann wird die Einladung zu der Aktionskarte den \enquote{Anschluss verlieren} und nicht in der Lage sein, sie zur Seite zu legen; in diesem Fall wirst du sie zu Beginn deines nächsten Zuges nicht ausspielen.

	\medskip

	Wenn du die Einladung nutzt, um ein Nomadencamp (Dominion – Hinterland) zu nehmen, wird die Einladung wissen, dass das Nomadencamp auf deinem Nachziehstapel zu finden ist, so dass du es in diesem Fall zur Seite legst (sofern du es nicht über einen anderen Effekt an eine andere Stelle verschoben hast).

	\medskip

	\emph{Errata:} Der letzte Satz auf der Karte müsste heißen: \enquote{Wenn du das tust, spiele sie zu Beginn deines nächsten Zuges.}
	\end{Spacing}}}
\end{tikzpicture}
\hspace{-0.6cm}
\begin{tikzpicture}
	\card
	\cardstrip
	\cardbanner{banner/white.png}
	\cardicon{banner/coin.png}
	\cardprice{5}
	\cardtitle{Höflinge}
	\cardcontent{}
\end{tikzpicture}
\hspace{-0.6cm}
\begin{tikzpicture}
	\card
	\cardstrip
	\cardbanner{banner/white.png}
	\cardicon{banner/coin.png}
	\cardprice{4}
	\cardtitle{Abbruch}
	\cardcontent{}
\end{tikzpicture}
\hspace{0.6cm}

	    % Basic settings for this card set
\renewcommand{\cardcolor}{basicgame}
\renewcommand{\cardextension}{Template}
\renewcommand{\cardextensiontitle}{Template}
\renewcommand{\seticon}{basic.png}

\clearpage
\newpage
\section{\cardextension \ - \cardextensiontitle}

\begin{tikzpicture}
	\card
	\cardstrip
	\cardbanner{banner/white.png}
	\cardicon{icons/coin.png}
	\cardprice{}
	\cardtitle{}
	\cardcontent{}
\end{tikzpicture}
\hspace{-0.6cm}
\begin{tikzpicture}
	\card
	\cardstrip
	\cardbanner{banner/gold.png}
	\cardicon{icons/coin.png}
	\cardprice{}
	\cardtitle{}
	\cardcontent{}
\end{tikzpicture}
\hspace{-0.6cm}
\begin{tikzpicture}
	\card
	\cardstrip
	\cardbanner{banner/purple.png}
	\cardicon{icons/coin.png}
	\cardprice{}
	\cardtitle{}
	\cardcontent{}
\end{tikzpicture}
\hspace{-0.6cm}
\begin{tikzpicture}
	\card
	\cardstrip
	\cardbanner{banner/blue.png}
	\cardicon{icons/coin.png}
	\cardprice{}
	\cardiconaddition{icons/hex.png}
	\cardpriceaddition{}
	\cardtitle{}
	\cardcontent{}
\end{tikzpicture}
\hspace{-0.6cm}
\begin{tikzpicture}
	\card
	\cardstrip
	\cardbanner{banner/green.png}
	\cardicon{icons/coin.png}
	\cardprice{}
	\cardiconaddition{icons/potion.png}
	\cardpriceaddition{}
	\cardtitle{}
	\cardcontent{}
\end{tikzpicture}
\hspace{-0.6cm}
\begin{tikzpicture}
	\card
	\cardstrip
	\cardbanner{banner/white.png}
	\cardtitle{\scriptsize{Empfohlene 10er Sätze\qquad}}
	\cardcontent{\emph{Name:}\\
	Karten ...

	\smallskip

	\emph{Name:}\\
	Karten ...

	\smallskip

	\emph{Name:}\\
	Karten ...

	\smallskip

	\emph{Name:}\\
	Karten ...

	\smallskip

	\emph{Name:}\\
	Karten ...

	\smallskip

	\emph{Name:}\\
	Karten ...}
\end{tikzpicture}
\hspace{0.6cm}
	\end{center}
\end{document}
